% !TeX program = pdflatex
\documentclass[11pt,a4paper]{article}
\usepackage[margin=25mm,headheight=15pt,footskip=27pt]{geometry}
\usepackage[utf8]{inputenc}
\usepackage{microtype,amsmath,amssymb,amsthm,mathtools,mathrsfs,bm}
\usepackage{booktabs,tabularx,array,enumitem,xcolor,graphicx,tikz,caption}
\usepackage{titlesec,fancyhdr,needspace,xurl,pgfplots,float,placeins}
\pgfplotsset{compat=1.18}
\usepgfplotslibrary{groupplots,fillbetween}
\usetikzlibrary{arrows.meta,positioning,calc}

\definecolor{ContentsGreen}{HTML}{176B4A}
\usepackage[colorlinks=true,linkcolor=ContentsGreen,citecolor=ContentsGreen,urlcolor=black,
 bookmarksnumbered=true]{hyperref}
\urlstyle{same}
\titleformat{\section}{\sffamily\large\bfseries}{\thesection}{.65em}{}
\titleformat{\subsection}{\sffamily\normalsize\bfseries}{\thesubsection}{.65em}{}
\titleformat{\subsubsection}{\sffamily\normalsize\bfseries}{\thesubsubsection}{.65em}{}
\titlespacing*{\section}{0pt}{2.5ex plus .8ex}{1ex}
\titlespacing*{\subsection}{0pt}{1.8ex plus .5ex}{.7ex}
\setlength{\parindent}{1.1em}\setlength{\parskip}{2pt plus .4pt}
\setlength{\emergencystretch}{2em}
\setlength{\abovedisplayskip}{7pt plus 2pt minus 2pt}
\setlength{\belowdisplayskip}{7pt plus 2pt minus 2pt}
\allowdisplaybreaks[2]\raggedbottom
\clubpenalty=10000\widowpenalty=10000
\captionsetup{font=small,labelfont=bf,skip=6pt}
\setlist{topsep=4pt,itemsep=2pt,parsep=0pt,leftmargin=1.6em}
\newtheorem{theorem}{Theorem}[section]
\newtheorem{lemma}[theorem]{Lemma}
\newtheorem{proposition}[theorem]{Proposition}
\newtheorem{corollary}[theorem]{Corollary}
\theoremstyle{definition}\newtheorem{definition}[theorem]{Definition}
\theoremstyle{remark}\newtheorem{remark}[theorem]{Remark}
\numberwithin{equation}{section}
\newcommand{\R}{\mathbb R}\newcommand{\E}{\mathbb E}\newcommand{\Pp}{\mathbb P}
\newcommand{\Id}{I}\newcommand{\F}{\mathrm F}
\newcommand{\Ksg}{{K_{\mathrm{sg}}}}\newcommand{\Cfour}{{C_4}}
\newcommand{\gmom}{\mathfrak g}
\newcommand{\Rmom}{\mathfrak r}
\newcommand{\norm}[1]{\lVert#1\rVert}\newcommand{\ip}[2]{\langle#1,#2\rangle}
\DeclareMathOperator{\tr}{tr}\DeclareMathOperator{\rank}{rank}
\DeclareMathOperator{\range}{range}\DeclareMathOperator{\diag}{diag}
\DeclareMathOperator{\sym}{sym}\DeclareMathOperator{\Var}{Var}\DeclareMathOperator{\op}{op}
\newcolumntype{Y}{>{\raggedright\arraybackslash}X}
\pagestyle{fancy}\fancyhf{}
\fancyhead[L]{\sffamily\footnotesize The logarithmic cost of Khatri--Rao embeddings}
\fancyfoot[C]{\small\thepage}\renewcommand{\headrulewidth}{.3pt}
\title{The Sharp Logarithmic Cost of Khatri--Rao Subspace Embeddings}
\author{Diar Heidary}
\date{14 October 2026}
\hypersetup{pdftitle={The Sharp Logarithmic Cost of Khatri-Rao Subspace Embeddings},
 pdfauthor={Diar Heidary}}
\begin{document}
\thispagestyle{plain}
{\centering\LARGE The Sharp Logarithmic Cost of\\[3pt]
Khatri--Rao Subspace Embeddings\par}
\vspace{12pt}
\begingroup
\renewcommand{\thefootnote}{\ensuremath{\ast}}
{\centering\large Diar Heidary\footnotemark[1]\par}
\footnotetext[1]{Unaffiliated work. This manuscript is the result of independent research
and was not conducted as part of a research program at Sharif University
of Technology, a master's thesis, or a doctoral dissertation.}
\endgroup
\vspace{3pt}
{\centering\href{mailto:diyar.heydari83@sharif.edu}{\texttt{diyar.heydari83@sharif.edu}}\par}

\vspace{8pt}
{\centering 14 October 2026\par}
\vspace{17pt}
\begingroup\setlength{\fboxsep}{9pt}\setlength{\fboxrule}{.7pt}
\noindent\fcolorbox{red!75!black}{red!2}{%
\begin{minipage}{\dimexpr\linewidth-2\fboxsep-2\fboxrule\relax}
\small
\textcolor{red!75!black}{\textbf{AI usage disclosure.}}
The mathematical arguments were developed using the large language models
Astra, 6-Sol, Fable 5.1, and Claude Opus 5.5 under the author's direction. This manuscript
was drafted in collaboration with Astra. The author takes full
responsibility for the content of the paper, including its mathematical
claims and proofs, all citations and attributions, and any remaining errors.
\end{minipage}}
\endgroup
\vspace{13pt}
\begin{abstract}
Khatri--Rao sketches use tensor products of small random vectors.
For fixed tensor order $d\ge2$, accuracy, confidence, and directional
sub-Gaussian parameters, we establish the worst-case subspace embedding
dimension $\Theta(r\log^{d-1}(er))$. Matching Gaussian, Rademacher,
and spherical lower bounds establish sharpness. In contrast, bounded
directional fourth moments suffice for $(1+\varepsilon)$-relative-error
low-rank approximation and regression with $O(r)$ samples at fixed
parameters. This distinction holds on the same input: approximation
can succeed with high probability while the embedding condition fails.
On specified Gaussian support classes, range-based outputs have exactly
the dense-Gaussian distribution. Our analysis also yields explicit
finite-sample moment certificates, polylogarithmic confidence bounds
for raw regression under sub-Gaussian factors, and angular-moment
bounds for expected full-range error.
\end{abstract}

\noindent\textbf{Keywords:} Khatri--Rao sketches; subspace embeddings;
relative-error approximation; matrix moments; tensor geometry.
\par\smallskip
\noindent\textbf{MSC 2020:} 65F55, 68W20, 60B20, 15A69.
\clearpage
\begingroup
\hypersetup{linkcolor=ContentsGreen}
\small\setlength{\parskip}{0pt}\setcounter{tocdepth}{2}
\tableofcontents
\endgroup
\clearpage


\section{Introduction}\label{sec:introduction}

Khatri--Rao sketches represent each randomized measurement probe as a
rank-one tensor $z=\bigotimes_{j=1}^d g^{(j)}$. Storing the small factor
vectors instead of an ambient random matrix reduces storage, while
Kronecker, canonical polyadic (CP), or tensor-train structure can make
the sketch efficient to apply. These advantages motivate applications
to tensor compression, polynomial kernels, regression, and eigensolvers
\cite{Ahle2020,RakhshanRabusseau2020,ChenJin2021,BGKL2026,SVB2026}.

The computational advantages raise two sample-complexity questions:
how many probes are needed to preserve every vector in a fixed
$r$-dimensional subspace, and must downstream approximation tasks use
the same number? Beretta--Musco~\cite[Theorem~1.3]{BM2026} establish
the nearly linear embedding bound $O(r\log^{d+1}(er))$ for
sub-Gaussian factors at fixed parameters. We determine the sharp
worst-case order $\Theta(r\log^{d-1}(er))$, saving two rank logarithms,
and prove that the remaining logarithmic factor is necessary for
Gaussian, Rademacher, and spherical factors.

Relative-error approximation has a different sample requirement.
Embedding demands two-sided spectral control, including an upper bound
that can fail because of a single large projected probe. The
approximation guarantees instead follow from lower injection and
control of the average interaction with the residual. Bounded
directional fourth moments suffice for these two ingredients, yielding
$(1+\varepsilon)$-relative-error low-rank approximation and regression
with $O(r)$ samples at fixed parameters. We prove that approximation
can succeed while embedding fails on the same input, and identify
Gaussian supports on which range-based outputs have exactly the
dense-Gaussian law.

\subsection{Setting and guarantees}\label{sec:setting}
Let $N=\prod_{j=1}^d n_j$ and let the tensor order satisfy $d\ge2$.
We draw iid probes with independent isotropic factors:
\begin{equation}\label{eq:intro:model}
 z_i=\bigotimes_{j=1}^d g_i^{(j)}\in\R^N,
 \qquad\E g_i^{(j)}g_i^{(j)\mathsf T}=I_{n_j},\qquad
 \Omega=m^{-1/2}[z_1\ \cdots\ z_m],\quad\Pi=\Omega^\top.
\end{equation}
The sketch size $m$ counts probe vectors, equivalently columns of
$\Omega$ or rows of $\Pi$. Coordinates within a factor may be dependent
unless specified otherwise. Every input and subspace is deterministic
before sampling. For two-sided embedding, we assume the directional
sub-Gaussian condition
\begin{equation}\label{core:eq:model}
 \E\exp\!\left(\frac{\ip{u}{g^{(j)}}^2}{\Ksg^2}\right)\le2
 \qquad(\norm u=1),
\end{equation}
in addition to the isotropy and independence in \eqref{eq:intro:model}.
This fixes our directional sub-Gaussian convention.
The approximation theorem needs only the weaker ambient condition
\begin{equation}\label{app:eq:fourth}
 \E zz^\top=I_N,\qquad
 \E|\ip a z|^4\le\Cfour\norm a^4\quad(a\in\R^N).
\end{equation}

\begin{definition}[Embedding, injection, and relative approximation]
\label{def:guarantees}
A sketch distribution is an \emph{oblivious subspace embedding} (OSE)
with squared-norm distortion $\varepsilon$ and failure probability
$\delta$ if, for every fixed isometry $U\in\R^{N\times r}$,
\begin{equation}\label{eq:intro:ose}
 \Pp\{\norm{U^\top\Omega\Omega^\top U-I_r}\le\varepsilon\}
 \ge1-\delta.
\end{equation}
A \emph{lower injection} with parameter $\alpha>0$ requires only
$U^\top\Omega\Omega^\top U\succeq\alpha I_r$ with that probability.
For a fixed matrix $A$, a \emph{relative approximation} satisfies
$\norm{A-\widehat A}_{\F}^2\le(1+\varepsilon)\norm{A-A_r}_{\F}^2$,
where $A_r$ is a best rank-at-most-$r$ approximation.
For regression, relative error refers to the original least-squares
residual, as stated in Theorem~\ref{app:thm:relative}.
\end{definition}

Table~\ref{tab:literature} separates the embedding and approximation
guarantees. All rank-dependence comparisons fix tensor order, factor-law
parameters, accuracy, and confidence.

\begin{table}[tbp]
\centering\small
\renewcommand{\arraystretch}{1.22}
\begin{tabularx}{\linewidth}{@{}>{\raggedright\arraybackslash}p{.31\linewidth}>{\raggedright\arraybackslash}p{.23\linewidth}Y@{}}
\toprule
Result & Factor law & Sketch size / guarantee\\
\midrule
\multicolumn{3}{@{}l}{\textbf{Two-sided embedding}}\\\addlinespace[2pt]
Bujanovi\'c et al.~\cite{BGKL2026} & Gaussian, $d=2$ & $O(r^{3/2})$\\\addlinespace[3pt]
Beretta--Musco~\cite{BM2026} & Sub-Gaussian & $O(r\log^{d+1}(er))$\\\addlinespace[3pt]
This work: matching bounds & Sub-Gaussian upper; Gaussian, Rademacher, spherical lower & $\Theta(r\log^{d-1}(er))$\\\addlinespace[3pt]
\midrule
\multicolumn{3}{@{}l}{\textbf{Approximation}}\\\addlinespace[2pt]
Cama\~no et al.~\cite{CEMT2025} & Fourth-moment product laws & $O(r)$, constant factor\\\addlinespace[3pt]
Saibaba et al.~\cite{SVB2026} & Sub-Gaussian entries & $O(r^d)$, constant factor\\\addlinespace[3pt]
Theorem~\ref{app:thm:relative} & Bounded fourth moments & $O(r)$, $1+\varepsilon$ error\\\addlinespace[3pt]
Proposition~\ref{exact:prop:separation} & Gaussian, specified supports & Exact Gaussian law despite failure on the leading subspace\\\addlinespace[3pt]
\bottomrule
\end{tabularx}
\caption{Rank dependence at fixed parameters.
The $\Theta$ statement is worst-case over fixed subspaces and admissible
mode dimensions; it does not assert a lower bound for every factor law.
For $r\ge2$, $\log(er)$ and $\log r$ give the same asymptotic order.
The last row is a distributional comparison on the specified support class.}
\label{tab:literature}
\end{table}

\subsection{Why the logarithm appears}\label{sec:heuristics}
Write $y_i=U^\top z_i$ and
$M=U^\top\Omega\Omega^\top U=m^{-1}\sum_i y_i y_i^\top$ for the
restricted Gram matrix. The lower-bound mechanism has three steps.
\begin{enumerate}
\item \textbf{One large row breaks embedding.} Positivity gives
\[
 \lambda_{\max}(M)\ge\frac{\max_i\norm{y_i}^2}{m}.
\]
Thus $\norm{y_i}^2>(1+\varepsilon)m$ for even one probe violates
the embedding condition.

\item \textbf{A shared scalar amplifies all coordinates.}
On the Gaussian common-factor witness
$Ue_a=e_a\otimes e_1\otimes\cdots\otimes e_1$,
\[
 y_i=\left(\prod_{j=2}^d\eta_{i,j}\right)\xi_i,
 \qquad \xi_i\sim N(0,I_r),\quad \eta_{i,j}\sim N(0,1),
\]
with all displayed Gaussian variables independent. The weight
$w_i=\prod_{j=2}^d\eta_{i,j}^2$ multiplies the entire squared norm:
$\norm{y_i}^2=w_i\norm{\xi_i}^2$.

\item \textbf{The largest weight supplies the logarithm.}
At fixed $d$, the maximum weight has order $(\log m)^{d-1}$,
while $\norm{\xi_i}^2$ concentrates around $r$. Together these give
the necessary term $r\log^{d-1}r$ in the witness constructions.
Rademacher factors reproduce the mechanism through rare alignment
in common modes of dimension proportional to $\log r$
(Proposition~\ref{lower:prop:rademacher}).
\end{enumerate}

The upper bound must control these extremes uniformly over the target
subspace, as well as the collective covariance fluctuations. Under
the present moment assumptions, the covariance route gives, at its
stated confidence,
\[
 \norm{M-I_r}\lesssim_{d,\Ksg}
 \sqrt{r/m}+\frac{\max_i\norm{y_i}^2}{m}.
\]
Gaussian averaging bounds the largest squared projected row norm by
$O_{d,\Ksg}(rs^{d-1}+s^d)$ with $s=\log(m/\beta)$ and failure
probability $\beta$. The resulting term
$(rt^{d-1}+t^d)/\varepsilon$ in \eqref{eq:intro:sharp} reflects the
large-row scale. Its term $rt/\varepsilon^2$ includes a logarithm
from the trace-moment proof, rather than an unavoidable
dense-Gaussian fluctuation cost: the covariance route reduces it to
$r/\varepsilon^2$ when $\delta\ge2/r$. These estimates are bounds,
not an additive formula for the top eigenvalue; matching extremes
are asserted only for the constructed witnesses.

Subspace geometry can reduce the shared fluctuation. In the order-two
Gaussian common-Schmidt model, the weight becomes
$w=\sum_b\lambda_bh_b^2$, and $L=\max_b\lambda_b$ gives the scale
$1+L\log m$ for its maximum, up to constants at fixed confidence.
Spreading the weights therefore explains the transition
$r[1+L\log(er)]$.

\paragraph{Why approximation needs less.}
Relative-error approximation is controlled by two quantities: the
inverse $\norm{M^{-1}}$ on the leading subspace and a mean-square
interaction between that subspace and the residual. Lower injection
bounds the inverse, while fourth moments bound the interaction in
Frobenius norm without estimating a sample maximum. Neither step
requires an upper bound on $\lambda_{\max}(M)$, which allows
$O(r)$ samples at fixed accuracy and confidence. This is a statement
about the requirements of the error bound: a large leading eigenvalue
can still accompany a large residual interaction, and the realized
approximation error need not be monotone in the sample weights.

\subsection{Main results}
\begin{enumerate}
\item \textbf{Sharp embedding dimension.}
With $t=\log(4r/\delta)$, Theorem~\ref{core:thm:ose} gives
\begin{equation}\label{eq:intro:sharp}
 m\ge C_{d,\Ksg}\left[\frac{rt}{\varepsilon^2}
             +\frac{rt^{d-1}+t^d}{\varepsilon}\right].
\end{equation}
The covariance proof replaces $rt/\varepsilon^2$ by
$r/\varepsilon^2$ when $\delta\ge2/r$.
Gaussian, Rademacher, and spherical constructions match the fixed-parameter
rank order; the latter two use common-mode dimensions growing as $\log r$.
Compared with Beretta--Musco~\cite[Theorem~1.3]{BM2026}, the bound saves
two rank logarithms and gives the large-row term accuracy cost
$\varepsilon^{-1}$.

\item \textbf{Relative approximation with linear sketch size.}
Theorem~\ref{app:thm:relative} gives rangefinding, fixed-rank output,
and regression guarantees under bounded directional fourth moments, with
\begin{equation}\label{eq:intro:relative}
 m=O\!\left(\Cfour r+\Cfour\log(2/\delta)
                    +\frac{\Cfour r}{\varepsilon\delta}\right).
\end{equation}
For range-based objectives, concatenating independent sketches and
using range containment improves the confidence dependence to
$O(\Cfour r\varepsilon^{-1}[1+\log(1/\delta)])$.
For raw regression, the inverse-confidence cost is necessary over
general fourth-moment laws, while sub-Gaussian product laws admit
polylogarithmic confidence cost
(Theorem~\ref{app:thm:subgauss-confidence}).

\item \textbf{Linear-size approximation despite embedding failure.}
At $m=O(r)$, Gaussian and Rademacher sketches achieve relative error
on fixed inputs whose leading-space Gram matrices have a bounded
lower edge and an excessive upper edge
(Corollary~\ref{separation:cor:same-input}). On Gaussian common-factor
supports, Proposition~\ref{exact:prop:separation} strengthens this to an
exact dense-Gaussian law for range-based outputs, including the classical
expected-error bound.

\item \textbf{Subspace geometry.}
Theorem~\ref{geom:thm:schmidt} identifies the order-two common-Schmidt
transition $r[1+L\log(er)]$, where $L$ is the largest Schmidt weight.
Input-dependent approximation bounds retain the leading-space fourth
moment and its interaction with the residual. Spherical normalization
improves moment constants while retaining the worst-case logarithm.
\end{enumerate}

Section~\ref{core:sec:embedding} proves the sharp embedding order,
giving the covariance proof and stating a separate explicit
moment certificate; the detailed rank-one argument is in
Appendix~\ref{appendix:covariance}. Section~\ref{geom:subsec:geometry} explains the effect of
subspace geometry and spherical normalization. Section~\ref{app:sec:relative}
develops approximation, the same-input separation, and exact range laws.
Numerics and discussion follow in Sections~\ref{sec:numerics}--\ref{sec:discussion}.

\subsection{Related work}
\paragraph{Tensor sketches.}
Independent sign-product sketches are analyzed by
Meister--Sarl\'os--Woodruff~\cite{MSW2019}; direct product and recursive
polynomial-kernel sketches are analyzed by Ahle et al.~\cite{Ahle2020}; Rakhshan--Rabusseau~\cite{RakhshanRabusseau2020}
study tensorized random projections. Fast Kronecker transforms and
Kronecker sign constructions appear in
\cite{JinKoldaWard2021,BambergerKrahmerWard2022}, with sampling laws
different from the independent product probes here.
The Gaussian OSE analysis of Bujanovi\'c et al.~\cite{BGKL2026}, the
rangefinding and tensor algorithms of Saibaba et al.~\cite{SVB2026},
and Beretta--Musco's general OSE theorem~\cite{BM2026} give the direct
comparisons in Table~\ref{tab:literature}. The present upper and lower
bounds determine the fixed-parameter worst-case rank dependence.
TTStack interpolates between Khatri--Rao and Gaussian tensor-train
sketches and gives OSE and injection budgets linear in tensor order
through a richer sketch family~\cite{TTStack2026}. These guarantees
do not remove the obstruction for the independent rank-one product
probes studied here.

\paragraph{Covariance and moments.}
Rank-one covariance approximation is developed in
\cite{Rudelson1999,Oliveira2010}; scalar moment inequalities for sums are
studied in~\cite{Latala1997}. Gaussian pairings and trace inequalities
are classical~\cite{Isserlis,Janson,TroppElementary}.
Our matrix argument retains the rank-one direction, preserves centering
in the first two insertion orders, and uses only a logarithmic terminal
moment order. The covariance proof applies Tikhomirov's covariance
theorem~\cite{Tikhomirov2018}; Bousquet's empirical-process inequality
\cite{Bousquet2002} supplies arbitrary confidence. Gaussian smoothing
and Catoni's influence function give the optimized lower-covariance
profile~\cite{Oliveira2016,Catoni2012}; its rate and method are classical.
Matrix-chaos inequalities also include Khatri--Rao applications
\cite{MatrixChaos2025}. Sparse-chaos estimates of
Dai--He--Wang--Zhu~\cite{DHWZ2026} account for sparsity and input
geometry in general; their scope differs from the dense worst-case
rank dependence considered here.

\paragraph{Approximation.}
Lower covariance and small-ball arguments appear in
\cite{Oliveira2016,KM2015}; Cama\~no et al.~\cite[Theorem~5.1]{CEMT2025}
give fourth-moment Khatri--Rao injection.
The range, inverse-moment, and regression reductions come from
\cite{Sarlos2006,CW2009,HMT2011,DMMW2011,CNW2016,TYUCPSD2017}; see also tensor regression
and rank-one estimation in~\cite{ChenJin2021,BK2021}.
Range invariance for Gaussian and spherical factors is noted in
\cite[Section~5.4]{CEMT2025}. We use these ingredients to obtain
explicit fourth-moment relative guarantees and a same-input embedding
separation, and identify support classes on which the full range output
has the dense-Gaussian law.

The August 2026 update to the Simons workshop problems
\cite[Problems~5.1--5.2]{Simons2026} records the counterexamples of
Townsend--Wang~\cite{TW2026} and asks for convenient additional assumptions.
Our fourth-moment theorem supplies one sufficient condition for product
probes; it does not identify the weakest possible condition.

\subsection{Notation and scope}
Vectors use the Euclidean norm and matrices the spectral norm;
$\norm{\cdot}_{\F}$ is the Frobenius norm. Traces are unnormalized,
$\preceq$ is the Loewner order, and the dagger denotes the Moore--Penrose
inverse. $P_{\mathcal L}$ is the orthogonal projector onto $\mathcal L$.
\begin{center}\begin{minipage}{\linewidth}\small
\begin{tabularx}{\linewidth}{@{}p{.28\linewidth}Y@{}}
\toprule Symbol & Meaning\\\midrule
$d,n_j,N$ & Tensor order, mode dimensions, and ambient dimension.\\
$r,m$ & Target rank and sketch size.\\
$\Omega,\Pi$ & Normalized probe matrix and sketching map $\Pi=\Omega^\top$.\\
$\Ksg,\Cfour$ & Directional sub-Gaussian and fourth-moment parameters.\\
$q,\gmom_q$ & Half of the terminal trace-moment order and the scalar Gaussian moment $(2q-1)!!$.\\
$Z,\Delta_\Omega$ & Restricted covariance error and ambient covariance error.\\
$H_i,X_i$ & Projected covariance sample and its centered version.\\
$C_Q,\kappa(Q,W)$ & Leading-space fourth moment and residual interaction.\\
$\mathfrak m_{\rm inj}$ & Sample budget for lower injection.\\
$\tau_r^{\rm F}(A),\tau_r^{\rm tr}(K)$ & Squared Frobenius tail $\|A-A_r\|_\F^2$ and PSD trace tail $\sum_{j>r}\lambda_j(K)$.\\
$\nu,\Rmom$ & Variance bound and higher-moment radius in the appendix comparison.\\
$\lambda_b,L,\rho$ & Schmidt weights, their maximum, and their purity $\sum_b\lambda_b^2$.\\
\bottomrule
\end{tabularx}
\end{minipage}\end{center}
Constants depend only on their displayed subscripts. Inputs are fixed
before sampling, and the guarantees concern exact arithmetic.

\section{Sharp embedding dimension}
\label{core:sec:embedding}

Two effects govern the embedding dimension: accumulated covariance
fluctuations and unusually large projected rows. The covariance route
controls both and proves the sharp fixed-parameter order. A separate
rank-one moment certificate supplies explicit constants and a complementary
high-confidence bound; its proof is deferred to
Appendix~\ref{appendix:covariance}. Gaussian and Rademacher witnesses
show that a single large row can force the logarithmic overhead;
spherical factors are treated in Section~\ref{geom:subsec:geometry}.

Work in the tensor model of Section~\ref{sec:setting}, under
the sub-Gaussian condition \eqref{core:eq:model}. Fix an isometry
$U:\R^r\to\mathcal H=\bigotimes_{j=1}^d\R^{n_j}$, and set
\begin{equation}\label{core:eq:rows}
 y_i=U^{\mathsf T}\bigotimes_{j=1}^d g_i^{(j)},\qquad
 H_i=y_iy_i^{\mathsf T},\qquad X_i=H_i-I_r,\qquad
 Z=\frac1m\sum_{i=1}^mX_i.
\end{equation}
Independence across the factors gives $\E H_i=I_r$. The sketch with rows
$m^{-1/2}(\bigotimes_jg_i^{(j)})^{\mathsf T}$ embeds $\range(U)$ with
squared-norm distortion at most $\varepsilon$ precisely when
$\norm Z\le\varepsilon$.

\begin{theorem}[Sharp dimension order]\label{core:thm:ose}
For every fixed $d\ge2$ and $\Ksg$, there is a constant $C_{d,\Ksg}$ such that,
for $0<\varepsilon,\delta<1$ and $t=\log(4r/\delta)$,
\begin{equation}\label{core:eq:ose}
 m\ge C_{d,\Ksg}\left[
 \frac{rt}{\varepsilon^2}+
 \frac{rt^{d-1}+t^d}{\varepsilon}\right]
 \quad\Longrightarrow\quad
 \Pp\{\norm Z>\varepsilon\}\le\delta.
\end{equation}
For each of the standard Gaussian and independent-coordinate Rademacher
factor laws, there are deterministic
$r$-dimensional subspaces and $c_d>0$ for which
\begin{equation}\label{core:eq:lower-main}
 \inf_{1\le m\le c_dr(\log r)^{d-1}}
 \Pp\{\norm Z>1/2\}\longrightarrow1
 \qquad(r\longrightarrow\infty).
\end{equation}
For either law, and for any sub-Gaussian class containing one of them,
the worst-case embedding dimension over admissible mode dimensions
and deterministic $r$-dimensional subspaces is
$\Theta_{d,\Ksg,\varepsilon,\delta}(r\log^{d-1}(er))$ as $r\to\infty$,
with $d,\Ksg,\varepsilon,\delta$ fixed.
\end{theorem}

\begin{proof}[Upper bound and conclusion of Theorem~\ref{core:thm:ose}]
For the upper bound, take $q=\lceil t\rceil$ in
Proposition~\ref{core:prop:certificate}. Its parameters satisfy
$h_q\le\sqrt{2q}$, $\eta_q\le\sqrt e$, and $c_q\le\sqrt q$
(the last inequality is AM--GM on $1,3,\ldots,2q-1$), with $q\le2t$.
Thus, outside an event of probability $\delta$,
\[
 \norm Z\le\sqrt e\sqrt{vq/m}+\sqrt{2e}\,q\Rmom_q/m.
\]
Here $v=O_{d,\Ksg}(r)$ and
$q\Rmom_q=O_{d,\Ksg}(rq^{d-1}+q^d)$.
The square-root term controls covariance fluctuations, while the term
proportional to $m^{-1}$ accounts for large projected rows.
Choosing the constant in \eqref{core:eq:ose} so that each term is at
most $\varepsilon/2$ proves the claimed upper bound. For $d=2$, the
term $rt/\varepsilon$ is absorbed by $rt/\varepsilon^2$.
At fixed accuracy and confidence, this gives
$O_{d,\Ksg,\varepsilon,\delta}(r\log^{d-1}(er))$.

The Gaussian construction in Section~\ref{core:sec:lower} and the
Rademacher construction in Proposition~\ref{lower:prop:rademacher}
prove the lower bound. In fact, their covariance error is at least one
with probability tending to one uniformly over the indicated widths.
They therefore rule out every fixed $\varepsilon<1$ and $\delta<1$,
which proves the matching worst-case order.
\end{proof}

Section~\ref{core:sec:local} gives an independent covariance proof of
the fixed-parameter upper order, followed by the matching lower bounds
in Section~\ref{core:sec:necessity}. Section~\ref{core:sec:certificate}
closes the section with the explicit moment certificate used above;
its rank-one comparison argument is developed in
Appendix~\ref{appendix:covariance}.
The moment route improves the $\varepsilon^{-1}$ term of the
general-confidence covariance bound when $\delta<2/r$.

\subsection{Directional moments and covariance bounds}
\label{core:sec:local}

Isotropy fixes the quadratic moment of every unit projection.
Subtracting this known contribution, together with the constant term,
from the exponential moment leaves a smaller bound for the higher
moments. Set
\begin{equation}\label{core:eq:notation}
 \theta=1-\Ksg^{-2},\qquad
 \beta_p=\sup_{x>0}\frac{x^p}{e^x-1-x}\quad(p\ge2).
\end{equation}
Here $\beta_2=2$ and $\beta_p\le p!$. Retaining the full exponential
remainder, instead of a single Taylor coefficient, improves every higher
scalar budget. Jensen's inequality gives
$\Ksg^2\ge1/\log2$, so $\theta$ is positive. The proof then applies the
same scalar estimate successively to the independent tensor factors.

\begin{lemma}[Isotropy-anchored moments]\label{core:lem:factor}
For an isotropic factor $g$ satisfying \eqref{core:eq:model}, every $H\succeq0$ and
integer $p\ge2$ satisfy
\begin{equation}\label{core:eq:factor-moment}
 \E(g^{\mathsf T}Hg)^p\le
 \theta\beta_p\Ksg^{2p}(\tr H)^p.
\end{equation}
Consequently, for every deterministic tensor $a\in\mathcal H$,
\begin{equation}\label{core:eq:directional}
 \E\left|\ip a{\bigotimes_{j=1}^dg^{(j)}}\right|^{2p}
 \le[\theta\beta_p\Ksg^{2p}]^d\norm a^{2p}.
\end{equation}
\end{lemma}
\begin{proof}
If $H\ne0$, its normalized eigenvalues are convex weights. Convexity of
the exponential and \eqref{core:eq:model} give $\E e^Q\le2$, where
$Q=g^{\mathsf T}Hg/(\Ksg^2\tr H)$. Isotropy gives $\E Q=\Ksg^{-2}$.
Thus $\E(e^Q-1-Q)\le\theta$, and
$Q^p\le\beta_p(e^Q-1-Q)$ proves \eqref{core:eq:factor-moment}.
The case $H=0$ is immediate.

More generally, for a linear map $T$ from $\mathcal H$ to a Euclidean
space, induction on the number of factors proves
\[
 \E\norm{T(\bigotimes_jg^{(j)})}^{2p}
 \le[\theta\beta_p\Ksg^{2p}]^d\norm T_{\F}^{2p}.
\]
Indeed, after conditioning on the last factor, apply the induction
hypothesis to the partially contracted map $T_g$. Its squared Frobenius
norm is $g^{\mathsf T}Hg$ for a positive semidefinite $H$ with
$\tr H=\norm T_{\F}^2$. Apply \eqref{core:eq:factor-moment} to finish the
induction. A scalar target gives \eqref{core:eq:directional}.
\end{proof}

For Gaussian averaging, write the scalar and radial moments as follows.
\begin{equation}\label{core:eq:gaussian-moments}
 \begin{aligned}
 \gmom_p&=\E G^{2p}=(2p-1)!!,
       &&G\sim N(0,1),\quad\gmom_0=1,\\
 D_{r,p}&=\E\norm{G_r}^{2p}=\prod_{h=0}^{p-1}(r+2h),
       &&G_r\sim N(0,I_r),\quad D_{r,0}=1.
 \end{aligned}
\end{equation}
\paragraph{The covariance argument.}
A classical covariance estimate proves the sharp fixed-parameter
upper order.
\begin{theorem}[Sample covariance bounds]\label{variance:thm:refinement}
Under the assumptions of Theorem~\ref{core:thm:ose}, put
$t=\log(4r/\delta)$ and $u=\log(8/\delta)$.
There is a constant $C_{d,\Ksg}$ such that, if $\delta\ge2/r$, then
\begin{equation}\label{variance:eq:moderate}
 m\ge C_{d,\Ksg}\left[\frac r{\varepsilon^2}
       +\frac{rt^{d-1}+t^d}{\varepsilon}\right]
 \quad\Longrightarrow\quad\Pp\{\norm Z>\varepsilon\}\le\delta.
\end{equation}
For every $0<\varepsilon,\delta<1$, one also has
\begin{equation}\label{variance:eq:general}
 m\ge C_{d,\Ksg}\left[\frac{r+u}{\varepsilon^2}
       +\frac{u(rt^{d-1}+t^d)}{\varepsilon}\right]
 \quad\Longrightarrow\quad\Pp\{\norm Z>\varepsilon\}\le\delta.
\end{equation}
\end{theorem}
The covariance input is Tikhomirov's theorem
\cite[Corollary~2]{Tikhomirov2018}; Bousquet's inequality
\cite[Theorem~2.3]{Bousquet2002} supplies the second confidence bound.
The proof of both bounds is in Appendix~\ref{appendix:variance-refinement}.
At fixed confidence, \eqref{variance:eq:moderate} removes the logarithm
from the variance term. At arbitrary confidence,
\eqref{variance:eq:general} pays an additional confidence logarithm
in the term proportional to $\varepsilon^{-1}$.
The finite-moment certificate \eqref{core:eq:explicit-m} provides explicit constants.

\begin{proof}[Covariance proof of the fixed-parameter upper bound]
We prove the upper order in Theorem~\ref{core:thm:ose} directly.
Let $G\sim N(0,I_r)$ be independent of $y$. Integrate $|\ip Gy|^{2q}$
in the two orders using Lemma~\ref{core:lem:factor}. Since
$\gmom_q\ge q!$ and $D_{r,q}\le(r+2q)^q$,
\[
 (\E\norm y^{2q})^{1/q}\le\Ksg^{2d}q^{d-1}(r+2q),\qquad q\ge2.
\]
Markov and a union bound with $q\asymp\log(m/\delta)$ control the
maximum of the $m$ squared row norms. Independent random signs center
the rows without changing their Gram matrices. Their directional sixth
moments are bounded, so Tikhomirov's Corollary~2, for $m\ge2r$, gives
\[
 \norm Z\le C_{d,\Ksg}\left[\sqrt{r/m}
                  +m^{-1}\max_i\norm{y_i}^2\right]
\]
outside an event of probability $1/r$. At fixed $\varepsilon,\delta$,
the choice $m\asymp r\log^{d-1}(er)$ makes this at most
$\varepsilon$ with failure at most $\delta$ for all sufficiently
large $r$. A larger fixed constant handles the finitely many smaller
ranks. Appendix~\ref{appendix:variance-refinement} closes the
sample-size logarithm uniformly in accuracy and gives both displayed
confidence bounds, using Bousquet for the second one.
\end{proof}

\paragraph{Where the two logarithms are saved.}
Beretta--Musco's proof first bounds squared row norms by
$O(r\log^d(m/\delta))$, then applies matrix Chernoff, paying another
$\log(r/\delta)$~\cite[Section~2]{BM2026}. Gaussian averaging above
divides the directional moment by $\gmom_q$, reducing the leading
row scale to $r q^{d-1}+q^d$. The covariance estimate then uses the
maximum row divided by $m$, without a further rank logarithm.
Proposition~\ref{core:prop:certificate} gives the complementary
explicit moment certificate.

\subsection{Why the logarithms are necessary}
\label{core:sec:necessity}
Both constructions separate two regimes. When $m<r$, rank deficiency
forces covariance error at least one. When $m\ge r$, a single large
sample suffices: positivity gives
\[
 \lambda_{\max}\!\left(m^{-1}\sum_i y_iy_i^\top\right)
 \ge m^{-1}\max_i\norm{y_i}^2.
\]
On a common-factor subspace, $y_i$ is an isotropic $r$-vector multiplied
by a product of scalar projections. Gaussian extremes amplify this
common scalar weight. For signs, prescribing all $h$ coordinates in one
factor to be $+1$ has probability $2^{-h}$ and produces a squared
projection of size $h$. Choosing $h$ proportional to $\log r$ makes
simultaneous alignment across the common factors likely among $m\ge r$
rows, at the same logarithmic scale as the Gaussian extreme.
\subsubsection{Gaussian factors}\label{core:sec:lower}
\begin{proof}[Gaussian lower bound in Theorem~\ref{core:thm:ose}]
Choose
$Ue_a=e_a\otimes e_1\otimes\cdots\otimes e_1$. For standard Gaussian
factors the projected row has the form
\[
 y_i=\left(\prod_{\ell=1}^{d-1}\eta_{i\ell}\right)\xi_i,
 \qquad\xi_i\sim N(0,I_r),\quad\eta_{i\ell}\sim N(0,1),
\]
with all variables independent. Write
$w_i=\prod_{\ell=1}^{d-1}\eta_{i\ell}^2$.
If $m<r$, the sample covariance has a zero eigenvalue, so
$\norm Z\ge1$.

For $m\ge r$ sufficiently large, put
$a=\sqrt{\log m/(d-1)}$. Integration of the Gaussian density on
$[a,a+1/a]$ gives
\[
 \Pp\{\eta^2\ge a^2\}\ge
 c_d m^{-1/[2(d-1)]}(\log m)^{-1/2}.
\]
Indeed, the interval has length $1/a$, and its minimum density is
$(2\pi)^{-1/2}\exp[-(a+1/a)^2/2]$; the remaining factor is bounded
below for $a\ge1$. Independence implies
\[
 \Pp\left\{\max_iw_i<
       \left(\frac{\log m}{d-1}\right)^{d-1}\right\}
 \le\exp\left[-c'_d\frac{\sqrt m}{(\log m)^{(d-1)/2}}\right].
\]
Choose a maximizing index $J$ using only the weights. Conditional on
the weights, $\xi_J$ is still a standard Gaussian vector. Moreover,
\[
 \Pp\{\norm{\xi_J}^2<r/2\}
 \le e^{r/4}2^{-r/2}\le e^{-r/16},
\]
by Markov applied to $e^{-\norm{\xi_J}^2/2}$.
Positivity of every summand consequently gives, outside these two
failure events,
\[
 \lambda_{\max}\left(\frac1m\sum_iH_i\right)
 \ge\frac{r}{2m}\left(\frac{\log m}{d-1}\right)^{d-1}.
\]
For $r\le m\le[4(d-1)^{d-1}]^{-1}r(\log r)^{d-1}$, this is at least
$2$, so $\norm Z\ge1$. Both exceptional probabilities tend to zero
uniformly in this range. For example, throughout the range,
$\sqrt m\ge\sqrt r$ and $\log m\le2\log r$ for large $r$.
This proves \eqref{core:eq:lower-main}.
Finally, $\E e^{G^2/\Ksg^2}=(1-2/\Ksg^2)^{-1/2}\le2$ precisely when
$\Ksg^2\ge8/3$, so this witness belongs to every factor class with such
a Gaussian-admitting parameter.
\end{proof}


\subsubsection{Rademacher factors}
\label{core:sec:rademacher}

\begin{proposition}[A bounded-entry obstruction]\label{lower:prop:rademacher}
Fix $d\ge2$. For every sufficiently large integer $r$, there is a
deterministic isometry $U_r:\R^r\to\bigotimes_{j=1}^d\R^{n_j}$,
with $n_1=r$ and $n_2=\cdots=n_d=O_d(\log r)$, such that independent
factors with independent Rademacher coordinates satisfy
\begin{equation}\label{lower:eq:rademacher-uniform}
 \inf_{1\le m\le c_dr(\log r)^{d-1}}
 \Pp\left\{\left\|\frac1m\sum_{i=1}^m y_iy_i^{\mathsf T}-I_r\right\|
                  >\frac12\right\}
 \ge 1-\exp(-\sqrt r),
 \qquad y_i=U_r^{\mathsf T}\bigotimes_{j=1}^dg_i^{(j)},
\end{equation}
where the infimum is over integer $m$ and one may take
\[
 c_d=\frac{1}{2[4(d-1)\log2]^{d-1}}.
\]
These factors obey \eqref{core:eq:model} for every $\Ksg^2\ge8/3$.
Consequently, the worst-case dimension order in
Theorem~\ref{core:thm:ose} is also necessary when the factor law is
restricted to independent Rademacher coordinates.
\end{proposition}
\begin{proof}
For $m<r$, rank deficiency makes the covariance deviation at least $1$
deterministically. It remains to consider $m\ge r$.
Take $r\ge2^{4(d-1)}$, and put
\[
 h=\left\lfloor\frac{\log_2r}{2(d-1)}\right\rfloor,
 \qquad b=h^{-1/2}(1,\ldots,1)^{\mathsf T}\in\R^h,
 \qquad U_re_a=e_a\otimes b^{\otimes(d-1)}\quad(1\le a\le r).
\]
The images of the coordinate vectors are orthonormal. Set $n_1=r$
and $n_j=h$ for $j\ge2$. With $\xi_i=g_i^{(1)}$ and
$s_{ij}=\langle b,g_i^{(j)}\rangle$, the projected samples are
\[
 y_i=\left(\prod_{j=2}^d s_{ij}\right)\xi_i,
 \qquad \|\xi_i\|^2=r.
\]
Let $E_i$ be the event that all $h$ coordinates of $g_i^{(j)}$ are
$+1$ for every $2\le j\le d$. The events $E_i$ are independent,
and
\[
 \Pp(E_i)=2^{-h(d-1)}\ge r^{-1/2}.
\]
On $E_i$ one has $s_{ij}=\sqrt h$, hence
$\|y_i\|^2=rh^{d-1}$. Every sample covariance summand is positive
semidefinite, so the occurrence of any $E_i$ implies
\begin{equation}\label{lower:eq:rademacher-spike}
 \lambda_{\max}\left(\frac1m\sum_{i=1}^m y_iy_i^{\mathsf T}\right)
 \ge\frac{rh^{d-1}}m.
\end{equation}
The assumption on $r$ gives
$h\ge\log r/[4(d-1)\log2]$. Thus the right side of
\eqref{lower:eq:rademacher-spike} is at least $2$ throughout
$m\le c_dr(\log r)^{d-1}$. The probability that
no $E_i$ occurs is at most
\[
 (1-r^{-1/2})^m\le\exp(-m/\sqrt r)\le\exp(-\sqrt r).
\]
Thus the covariance deviation is at least $1$ except on an event of
probability at most $\exp(-\sqrt r)$. Together with the rank-deficient
regime, this proves the uniform bound.

Finally, Rademacher vectors are isotropic and satisfy
$\E e^{t\langle g,u\rangle}\le e^{t^2/2}$ for unit $u$
by Hoeffding's bound~\cite{Hoeffding1963}; Gaussian integration gives
$\E e^{\langle g,u\rangle^2/\Ksg^2}\le(1-2/\Ksg^2)^{-1/2}\le2$
when $\Ksg^2\ge8/3$, verifying \eqref{core:eq:model}.
\end{proof}

\begin{remark}\label{lower:rem:bounded-entries}
The witness uses only entries of absolute value one. Its logarithmic
overhead comes from directional alignment across tensor factors:
$\langle b,g\rangle$ can attain $\sqrt h$ when its $h$ coordinates align.
The prescribed simultaneous alignment is rare for an individual row,
but occurs with probability tending to one among $m\ge r$ independent
rows for the chosen $h$. The dimensions of the common factors grow
with $r$; the proposition does not address a restriction that keeps
those mode dimensions fixed. Coordinatewise boundedness therefore
does not supply a bound $\|y_i\|^2\le C_dr$ uniform in $r$.
\end{remark}


\paragraph{An accuracy-sensitive lower obstruction.}
The heavy-row argument can retain the unit background instead of discarding
it. Proposition~\ref{lower:prop:accuracy} proves that, for Gaussian factors,
\[
 m=\Omega_d\!\left(\frac{r(\log r)^{d-1}}{\varepsilon_r}\right)
\]
is necessary at fixed confidence whenever
$\varepsilon_r r(\log r)^{d-1}\to\infty$. Its finite failure bound and proof
are in Appendix~\ref{appendix:spectral-extensions}. This refinement does
not determine the full joint dependence on confidence and accuracy.

\subsection{An explicit moment certificate}
\label{core:sec:certificate}
We finish with numerical constants for the moment bound used in the
upper-bound argument for Theorem~\ref{core:thm:ose}. The covariance
bound in Theorem~\ref{variance:thm:refinement} has a smaller variance
term when $\delta\ge2/r$; the moment argument avoids the additional
confidence logarithm in the large-row term of \eqref{variance:eq:general}.

Set
\begin{equation}\label{core:eq:constants}
 s_r=\frac{r+2+\sqrt{3r(r+2)}}6,\qquad
 v=[2\theta\Ksg^4]^d s_r.
\end{equation}
A quadratic Gaussian averaging argument, proved in
Appendix~\ref{appendix:factorial}, gives
\begin{equation}\label{core:eq:variance}
 \E H^2\preceq vI_r,\qquad
 V:=\E(H-I_r)^2\preceq(v-1)I_r.
\end{equation}
Here $s_1=1$ and $(1+\sqrt3)r/6\le s_r\le r$.
The sharper fourth-moment constant in Lemma~\ref{core:lem:sharp-sg-fourth}
is used in the individual moment profiles and approximation bounds;
the all-order factorial certificate retains the normalization $v$.

\begin{proposition}[Explicit finite-moment certificate]\label{core:prop:certificate}
Under the factor assumptions of Theorem~\ref{core:thm:ose}, use $v$ from
\eqref{core:eq:constants}. For an integer $q\ge1$ and $0<\delta<1$, put
\begin{align}\label{core:eq:scales}
 \mathcal R_q&=\tfrac12\Ksg^{2d}(2q)^{d-2}(r+2q),&
 \Rmom_q&=\mathcal R_q+1+\frac4{3v},\\
 h_1&=0,&h_q&=\frac{2(q-1)}{\sqrt{2q-3}}\quad(q\ge2),\notag\\
 c_q&=\gmom_q^{1/(2q)},&
 \eta_q&=(r/\delta)^{1/(2q)}.\label{core:eq:probability-factors}
\end{align}
Then
\begin{equation}\label{core:eq:explicit-closed}
 \Pp\left\{\norm Z>
 \eta_qc_q\left[\sqrt{v/m}+h_q\Rmom_q/m\right]\right\}\le\delta.
\end{equation}
In particular, writing $\mathsf a_q=\eta_qc_q\sqrt v$ and
$\mathsf b_q=\eta_qc_qh_q\Rmom_q$, squared-norm distortion
$\varepsilon$ is ensured with failure at most $\delta$ by
\begin{equation}\label{core:eq:explicit-m}
 m\ge\left\lceil
 \left(\frac{\mathsf a_q+
 \sqrt{\mathsf a_q^2+4\varepsilon\mathsf b_q}}{2\varepsilon}\right)^2
 \right\rceil.
\end{equation}
The integer $q$ may be selected to minimize this deterministic certificate.
\end{proposition}

The \hyperref[proof:certificate]{proof} is in
Appendix~\ref{appendix:shift}, after the moment estimates in
Appendices~\ref{appendix:factorial}--\ref{appendix:closed}.
It retains the matrix direction when converting scalar to rank-one moments.
Centering removes single insertions; two insertions retain the variance,
and higher insertions are bounded by a positive moment envelope.
Gaussian pairing coefficients solve the resulting scalar recurrence.
Only directional moments through order $4q$ are needed.
The upper-bound argument following Theorem~\ref{core:thm:ose} makes the choice
$q=\lceil\log(4r/\delta)\rceil$ to obtain its all-confidence bound.
The constants above belong to this moment certificate, rather than
to the covariance theorem.

For Gaussian, independent-sign, or spherical factors,
Proposition~\ref{core:prop:gaussian-factorial} gives the alternative
$v_{\rm G}=3^dr$ and the smaller law-specific radius
\eqref{core:eq:gaussian-radius}. Appendix~\ref{sec:sharper} also retains
individual moment budgets, zero initial moments, and centered terminal
insertions. Its discrete and integral certificates retain these finer inputs
without new probabilistic assumptions.

\section{Geometry beyond the worst case}
\label{geom:subsec:geometry}
The following supports isolate the scalar amplification behind the
worst-case logarithm. On a common-factor support, alignment in the
shared modes amplifies every projected coordinate at once. The
common-Schmidt construction replaces this product weight by a weighted
sum of squared contractions. For order-two Gaussian probes, the largest
Schmidt weight determines how much of the logarithmic penalty remains.

\subsection{Common-factor and common-Schmidt supports}
\begin{definition}[Common-factor and common-Schmidt supports]
\label{exact:prop:supports}
Let $z=g\otimes h$, where $g\sim N(0,\Id_{n_1})$ is independent of
$h=h_2\otimes\cdots\otimes h_d$ with independent standard Gaussian factors.
For the following supports the projected probe, conditional on $h$, is
$N(0,w(h)I_R)$ with $w(h)>0$. Its line is therefore uniform.
\begin{enumerate}
\item For orthonormal $a_1,\ldots,a_R$ and fixed unit vectors
$b_2,\ldots,b_d$, the common-factor basis
\[
 v_a=a_a\otimes b_2\otimes\cdots\otimes b_d
\]
has $w(h)=\prod_{j=2}^d\ip{b_j}{h_j}^2>0$ almost surely.
\item More generally, let $f_1,\ldots,f_\ell$ be orthonormal vectors in
the remaining tensor modes, let $\lambda_b>0$ with $\sum_b \lambda_b=1$, and let the
$R\ell$ vectors $e_{a,b}$ be jointly orthonormal in the first mode. Then
\[
 v_a=\sum_{b=1}^{\ell}\sqrt{\lambda_b}\,e_{a,b}\otimes f_b,
 \qquad a=1,\ldots,R,
\]
are orthonormal and have
$w(h)=\sum_b \lambda_b\ip{f_b}{h}^2>0$ almost surely.
\end{enumerate}
\end{definition}
\paragraph{Projected row law.}
In the first case $\ip{v_a}{z}=\ip{a_a}{g}\prod_{j=2}^d\ip{b_j}{h_j}$.
In the second, conditional covariance between coordinates $a,a'$ is
\[
 \sum_{b,c}\sqrt{\lambda_b\lambda_c}\ip{f_b}{h}\ip{f_c}{h}
                  \ip{e_{a,b}}{e_{a',c}}
 =\mathbf1_{\{a=a'\}}\sum_b \lambda_b\ip{f_b}{h}^2.
\]
The coordinates are jointly Gaussian conditional on $h$. Positivity holds
because the contraction with any nonzero $f_b$ is a nonzero polynomial in
Gaussian factors: product tensors span the ambient tensor space, and a
nonzero polynomial has a Gaussian zero set of measure zero.


\subsection{The common-Schmidt transition}
The worst-case logarithm can be located more precisely on the order-two
common-Schmidt supports of Definition~\ref{exact:prop:supports}.
Write their Schmidt weights as $\lambda_b\ge0$, $\sum_b\lambda_b=1$, and
put $L=\max_b\lambda_b$ and $\rho=\sum_b\lambda_b^2$. Their projected
Gaussian row has the distribution
\[
 y=\sqrt w\,g,\qquad
 w=\sum_b\lambda_bh_b^2,\qquad
 g\sim N(0,\Id_r),\quad h_b\sim N(0,1),
\]
with $g$ independent of the weight. Direct Gaussian moments give
\begin{equation}\label{geom:eq:variance}
 \E(yy^\top-\Id_r)^2
 =\bigl[r+1+2(r+2)\rho\bigr]\Id_r.
\end{equation}
The variance depends on the quadratic weight sum $\rho$, but the
embedding width also depends on extreme rows. The largest Schmidt
weight $L$ controls the size of a rare row that can violate the upper
spectral bound.
\begin{theorem}[Common-Schmidt transition]\label{geom:thm:schmidt}
For independent rows of the order-two common-Schmidt law, a universal
constant $C$ has the following property. For $0<\varepsilon,\delta<1$,
write $\tau_\delta=r+\log(8/\delta)$. Then
\begin{equation}\label{geom:eq:all-confidence}
 m\ge C\left[\frac{\tau_\delta}{\varepsilon^2}
       +\frac{L\tau_\delta}{\varepsilon}
          \log\frac{e\tau_\delta}{\varepsilon\delta}\right]
 \quad\Longrightarrow\quad
 \Pp\left\{\left\|m^{-1}\sum_i y_iy_i^\top-I_r\right\|
          >\varepsilon\right\}\le\delta.
\end{equation}
In particular, for fixed accuracy and confidence, the sufficient width
is $C_{\varepsilon,\delta}r[1+L\log(er)]$.
Conversely, for a universal $c>0$, the probability of exceeding the same
fixed distortion $\varepsilon$ tends to one
uniformly over the Schmidt weights and
$1\le m\le cr[1+L\log(er)]$ as $r\to\infty$.
\end{theorem}
The \hyperref[proof:schmidt]{proof in Appendix~\ref*{geom:appendix}}
first derives a conditional weighted-Gaussian upper bound and then
uses one extreme row for the matching lower bound.
At $L=1$, only one Schmidt weight is nonzero, so the support reduces to
a common-factor support and the theorem recovers the sharp
$\Theta_{\varepsilon,\delta}(r\log(er))$ width. For flat weights
$\lambda_b=1/s$, the scalar weight $w$ is an average of $s$ independent
squared Gaussians, and the sharp width is
$\Theta_{\varepsilon,\delta}(r[1+\log(er)/s])$. Thus spreading the
Schmidt spectrum removes the rank logarithm once $s\gtrsim\log r$,
at fixed accuracy and confidence.

\subsection{Spherical normalization}
Spherical factors remove Gaussian radial randomness. Averaging it out
can only reduce convex trace moments. Alignment of directions remains,
which is why the worst-case logarithm survives.

For independent uniform
unit-sphere vectors $u^{(j)}\in\R^{n_j}$, use isotropic factors
$x^{(j)}=\sqrt{n_j}\,u^{(j)}$. Gaussian polar decomposition and conditional
Jensen transfer every centered even trace moment of their projected
sample covariance from the corresponding Gaussian tensor model.
Consequently the upper bound in Theorem~\ref{core:thm:ose} holds for these
spherical factors, with the Gaussian constants. The directional fourth
constant of one factor improves to $3n_j/(n_j+2)$.

This finite-dimensional gain does not change the worst-case rank order:
for fixed $d\ge2$ there are common-factor supports with
$n_1=r$ and $n_2=\cdots=n_d=O_d(\log r)$ for which spherical factors
require $\Omega_d(r\log^{d-1}r)$ rows. Their first-mode row length is
constant, but rare common coordinates still produce a large common
weight. Appendix~\ref{geom:appendix} gives the
\hyperref[proof:spherical-transfer]{moment transfer} and the
\hyperref[proof:spherical-lower]{matching obstruction}.
Uniformity of the individual factors does
not imply uniformity of the direction after projection onto an arbitrary
tensor subspace.


\section{Approximation without embedding}\label{app:sec:relative}
For each fixed input, bounded directional fourth moments suffice for
relative-error rangefinding, low-rank approximation, and regression with
sketch size linear in the target rank at fixed accuracy and confidence.
Use iid isotropic probes with the normalization \eqref{eq:intro:model}.
The first results assume \eqref{app:eq:fourth};
Corollary~\ref{app:cor:input} later retains only the leading-space
moment and its residual interaction. Product laws satisfy the global
condition by Lemma~\ref{app:lem:product} in Appendix~\ref{appendix:injection}.

\subsection{The relative-error guarantee}
For $s>0$ and $0<\alpha,\eta<1$, define
\begin{align}
 A_s^{\rm inj}(\beta)&=1+\frac{2(s+2)}\beta
                  +\frac{4s(s+2)}{3\beta^2},\qquad
 B_\eta^{\rm inj}(\beta)=\frac\beta2+2\log\frac2\eta,
                  \label{app:eq:injection-profile}\\
 \mathfrak m_{\rm inj}(C,s;\alpha,\eta)
  &=\frac{2C}{(1-\alpha)^2}
       \min_{\beta>0} A_s^{\rm inj}(\beta)B_\eta^{\rm inj}(\beta).
                  \label{app:eq:injection-certificate}
\end{align}
Set $\mathfrak m_{\rm inj}(C,0;\alpha,\eta)=0$.
The minimum exists for $s>0$.

For the approximation theorems, let $B_r$ bound the aggregate
interaction with a unit residual direction:
\begin{equation}\label{app:eq:aggregate-budget}
 \E\norm{Q^\top z}^2\ip vz^2\le B_r
 \quad(Q^\top Q=I_r,\ \norm v=1,\ Q^\top v=0).
\end{equation}
One may always take $B_r=r\gamma(\Cfour)$, where
$\gamma(c)=\min\{c/2,c-1\}$, by Lemma~\ref{core:lem:orthogonal-fourth}.
Appendix~\ref{appendix:wick} gives smaller law-specific choices.

\begin{theorem}[Relative error for rangefinding and regression]
\label{app:thm:relative}
Assume \eqref{app:eq:fourth}, fix $r\ge1$ and $\varepsilon,\delta\in(0,1)$,
and suppose
\begin{equation}\label{app:eq:relative-budget}
 m\ge\max\left\{\mathfrak m_{\rm inj}(\Cfour,r;1/2,\delta/2),
                       \frac{8B_r}{\varepsilon\delta}\right\}.
\end{equation}
For each of the following fixed inputs, the corresponding conclusion holds
with probability at least $1-\delta$.
\begin{enumerate}
\item For $A\in\R^{n\times N}$, put $Y=A\Omega$ and let $P_Y$ be its
orthogonal range projector. If $A_r$ is a best rank-at-most-$r$
approximation and $\tau=\norm{A-A_r}_\F^2$, then
\begin{equation}\label{app:eq:rsvd-relative}
 \norm{A-P_YA}_\F^2\le(1+\varepsilon)\tau,
 \qquad
 \norm{A-\widehat A_r}_\F^2\le(1+\varepsilon)\tau,
\end{equation}
where $\widehat A_r$ minimizes the Frobenius error over matrices of rank at
most $r$ whose columns lie in $\range Y$. If $Q_Y$ is an orthonormal basis
of that range, one may compute $\widehat A_r=Q_Y(Q_Y^\top A)_r$.
\item For a full-column-rank design $A\in\R^{N\times r}$ and any fixed
multiple-response matrix $B\in\R^{N\times s}$, let
\[
 \widetilde X\in\arg\min_X\norm{\Pi(AX-B)}_\F^2.
\]
Then the residual in the original space obeys
\begin{equation}\label{app:eq:ls-relative}
 \norm{A\widetilde X-B}_\F^2
 \le(1+\varepsilon)\min_X\norm{AX-B}_\F^2.
\end{equation}
The failure bound is independent of the number of response columns.
\end{enumerate}
\end{theorem}

\subsection{Lower injection and residual interaction}
The mechanism is visible in SVD coordinates. If the leading block $\Omega_1$ has full row
rank, the rank-$r$ competitor $B_*=(A\Omega)\Omega_1^\dagger V_1^\top$
satisfies
\[
 \norm{A-B_*}_{\F}^2
 =\norm{A-A_r}_{\F}^2+
   \norm{\Sigma_2\Omega_2\Omega_1^\dagger}_{\F}^2.
\]
Lower injection controls the inverse factor, while orthogonal fourth
moments control the interaction between the leading and trailing singular
directions. These estimates play separate roles: only $\norm{M^{-1}}$
is used, with $M$ the leading Gram matrix, so no upper spectral bound on
$M$ is needed. The residual interaction is controlled in mean square
and may be dependent on $M$.
The deterministic reduction and approximate-matrix-product step are
classical~\cite{Sarlos2006,CW2009,HMT2011,CNW2016}. The contribution
here is a single product fourth-moment hypothesis that supplies both
lower injection and orthogonal residual control, with explicit constants.

\begin{lemma}[Explicit lower injection]\label{app:lem:injection}
For a fixed isometry $Q\in\R^{N\times r}$ and $0<\alpha,\eta<1$,
$m\ge\mathfrak m_{\rm inj}(\Cfour,r;\alpha,\eta)$ implies
\[
 \Pp\{Q^\top\Omega\Omega^\top Q\succeq\alpha\Id_r\}\ge1-\eta.
\]
In particular, half injection with failure at most $\delta/2$ holds if
\begin{equation}\label{app:eq:injection-budget}
 m\ge\mathfrak m_{\rm inj}(\Cfour,r;1/2,\delta/2).
\end{equation}
A sufficient count requiring no minimization is
$\frac{56}{3}\Cfour[r+2+2\log(4/\delta)]$.
\end{lemma}
Its proof is in Appendix~\ref{appendix:injection}; that appendix also
gives the numerical profile and optimizes the injection level used in
Theorem~\ref{app:thm:relative}.

\begin{lemma}[Orthogonal cross moment]\label{app:lem:cross}
Let $Q\in\R^{N\times r}$ be an isometry and let $Q^\top W=0$.
For $W\ne0$ define
\begin{equation}\label{app:eq:mixed-parameter}
 \kappa(Q,W)=
 \frac{\E[\norm{Q^\top z}^2\norm{W^\top z}^2]}{\norm W_\F^2},
 \qquad \Omega=m^{-1/2}[z_1\ \cdots\ z_m].
\end{equation}
For iid isotropic $z_i$, whenever the displayed expectation is finite,
\begin{equation}\label{app:eq:cross-exact}
 \E\norm{Q^\top(\Omega\Omega^\top-\Id_N)W}_\F^2
 =\frac{\kappa(Q,W)}m\norm W_\F^2.
\end{equation}
Under \eqref{app:eq:aggregate-budget}, $\kappa(Q,W)\le B_r$, so
\begin{equation}\label{app:eq:cross}
 \E\norm{Q^\top(\Omega\Omega^\top-\Id_N)W}_\F^2
 \le\frac{B_r}{m}\norm W_\F^2.
\end{equation}
\end{lemma}
\begin{proof}
The matrix $(Q^\top z)(z^\top W)$ has mean zero and squared Frobenius
norm $\norm{Q^\top z}^2\norm{W^\top z}^2$. Independence across samples
gives \eqref{app:eq:cross-exact}. Apply \eqref{app:eq:aggregate-budget} to each nonzero residual
column $w_b/\norm{w_b}$ and sum to obtain the upper bound.
The case $W=0$ is immediate.
\end{proof}

\begin{proof}[Proof of Theorem~\ref{app:thm:relative}]
For the first assertion, initially suppose $\rank A>r$ and write a reduced
SVD $A=U_1\Sigma_1V_1^\top+U_2\Sigma_2V_2^\top$, with the first block of
size $r$. Define $\Omega_j=V_j^\top\Omega$, $\Delta_\Omega =\Omega\Omega^\top-\Id_N$,
$M=\Omega_1\Omega_1^\top$, and
\[
 \mathcal E=\norm{\Sigma_2V_2^\top \Delta_\Omega V_1}_\F^2.
\]
Lemma~\ref{app:lem:cross}, with $W=V_2\Sigma_2$, gives
$\E \mathcal E\le B_r\tau/m$. On $M\succeq\Id_r/2$, the matrix $\Omega_1$ has full
row rank and
\[
 \Omega_2\Omega_1^\dagger=V_2^\top \Delta_\Omega V_1M^{-1}.
\]
The matrix $B_*=(A\Omega)\Omega_1^\dagger V_1^\top$ has rank at most $r$
and columns in $\range Y$. Its two error blocks have orthogonal right
singular spaces, so
\[
 \norm{A-B_*}_\F^2
 =\tau+\norm{\Sigma_2\Omega_2\Omega_1^\dagger}_\F^2
 \le\tau+4\mathcal E.
\]
Both errors in \eqref{app:eq:rsvd-relative} are at most this competitor's
error. Lemma~\ref{app:lem:injection} bounds failure of the lower event by
$\delta/2$, while Markov's inequality gives
\[
 \Pp\{4\mathcal E>\varepsilon\tau\}\le\frac{4B_r}{m\varepsilon}\le\delta/2.
\]
A union bound proves the claim, without independence of these two events.
If $\rank A\le r$, apply the injection lemma to its actual row space:
the range is captured exactly on that event and both errors vanish.

For regression, write $A=QR$ with $Q$ orthonormal and $R$ invertible, and
put $E=(\Id_N-QQ^\top)B$. The optimum original residual is $\norm E_\F^2$.
On $M=Q^\top\Omega\Omega^\top Q\succeq\Id_r/2$, the sketched problem has a
unique solution and its normal equations give
\[
 R\widetilde X=Q^\top B+M^{-1}G,
 \qquad G=Q^\top \Delta_\Omega  E.
\]
Orthogonality yields
\[
 \norm{A\widetilde X-B}_\F^2
 =\norm E_\F^2+\norm{M^{-1}G}_\F^2
 \le\norm E_\F^2+4\norm G_\F^2.
\]
Lemma~\ref{app:lem:cross} and Markov give failure at most $\delta/2$ for
the second term exceeding $\varepsilon\norm E_\F^2$; the lower event
contributes another $\delta/2$. When $E=0$, the normal equations already
give exact original residual zero and no division by $\norm E_\F^2$ is
needed. 
\end{proof}

\begin{corollary}[Input-dependent relative error]\label{app:cor:input}
In the rangefinding assertion of Theorem~\ref{app:thm:relative}, suppose
$\rank A>r$ and take $Q=V_1$, $W=V_2\Sigma_2$ from its SVD. In the
regression assertion take $Q$ to span $\range A$ and
$W=(I-QQ^\top)B$. For the relevant fixed pair define
\[
 C_Q=\sup_{\norm u=1}\E\ip{Qu}{z}^4.
\]
For iid isotropic probes, the same relative-error conclusion holds if
$C_Q<\infty$, $\kappa(Q,W)<\infty$, and
\begin{equation}\label{app:eq:input-budget}
 m\ge\max\left\{\mathfrak m_{\rm inj}(C_Q,r;1/2,\delta/2),
                  \frac{8\kappa(Q,W)}{\varepsilon\delta}\right\}.
\end{equation}
When $W=0$, only the injection condition on the relevant row or design
space is needed, and the residual is exactly zero.
\end{corollary}
\begin{proof}
Apply the injection lemma to $Q^\top z$ and replace the cross estimate
in the preceding proof by \eqref{app:eq:cross-exact}. No fourth-moment
bound outside the leading space and its residual interaction is used.
\end{proof}

If the common-factor or common-Schmidt support contains both spaces,
the projected Gaussian probe is $\sqrt w\,g$, with $g$ independent of $w$.
Orthogonal Gaussian coordinates give
\[
 C_Q=3\E w^2,\qquad \kappa(Q,W)=r\E w^2.
\]
For the order-two common-Schmidt law, the interaction is therefore
$r(1+2\rho)$. The universal fourth-moment bound would give the larger
allowance $3r(1+2\rho)/2$.

\subsection{PSD approximation and range confidence}
\Needspace{14\baselineskip}
For a positive semidefinite matrix $K$, write
\[
 \mathcal N_K(\Omega)
 =K\Omega(\Omega^\top K\Omega)^\dagger\Omega^\top K.
\]
Applying the range theorem to $K^{1/2}$ gives the Nystr\"om trace
guarantee, including fixed-rank truncation
(Corollary~\ref{app:cor:nystrom}). Range containment also improves confidence.
\begin{corollary}[Confidence amplification by range containment]
\label{app:cor:amplification}
Let
\[
 m_0=\left\lceil\max\left\{\mathfrak m_{\rm inj}(\Cfour,r;1/2,1/8),\ \frac{32B_r}{\varepsilon}\right\}\right\rceil,
 \qquad b=\left\lceil\frac{\log(1/\delta)}{\log4}\right\rceil.
\]
A single raw iid sketch with $m=bm_0$ columns satisfies
\eqref{app:eq:rsvd-relative} and \eqref{app:eq:nystrom-relative} with
probability at least $1-\delta$, for their respective fixed inputs. In
particular, the sufficient count is
\[
 O\!\left(\frac{\Cfour r}{\varepsilon}[1+\log(1/\delta)]\right).
\]
\end{corollary}
Its \hyperref[proof:amplification]{proof in Appendix~\ref*{appendix:amplification}} partitions one iid
sketch into independent groups. The full range contains every group
range, so all groups must fail before the combined approximation fails.
This argument does not apply to the minimizer of one stacked regression
objective. Under only a fourth-moment assumption, this distinction is
necessary: Proposition~\ref{app:prop:regression-lower} proves a
$1/(\varepsilon\delta)$ lower bound for the raw solve, already at rank one.
It does not assert this lower bound for a fixed Gaussian tensor law.


\subsection{A separation on the same input}
\label{exact:sec:laws}
For suitable fixed inputs, relative-error approximation and an upper
spectral obstruction occur together for the same draw of the sketch.
Intersecting the two probability bounds gives a high-probability
separation for Gaussian and Rademacher factors, without requiring
independence of the two events.
\begin{figure}[H]
\centering
\begingroup
\definecolor{mechInk}{HTML}{2F3034}
\definecolor{mechBlue}{HTML}{254BE8}
\definecolor{mechCyan}{HTML}{20CDEC}
\definecolor{mechRose}{HTML}{EA5CB5}
\definecolor{mechRoseInk}{HTML}{9B1E65}
\definecolor{mechLavender}{HTML}{7042DF}
\definecolor{mechIce}{HTML}{CBF5FC}
\definecolor{mechPeach}{HTML}{FFB585}
\definecolor{mechLime}{HTML}{DEF58E}
\resizebox{\linewidth}{!}{%
\begin{tikzpicture}[font=\small,text=mechInk,>={Stealth[length=2.1mm]}]
\def\cw{0.37}\def\ch{0.255}
\node[anchor=west,font=\small\bfseries] at (0,5.1) {(a) The restricted sketch};
\node[anchor=west] at (0,4.55) {$V^\top\Omega=m^{-1/2}\,G\,D$};
\foreach \h [count=\i from 0] in {0.30,0.64,0.20,0.45,0.76,0.32,1.25,0.49,0.25,0.60,0.40,0.17}{
  \ifnum\i=6
    \fill[mechRose]
      ({0.95+\i*\cw+0.05},2.86) rectangle ++({\cw-0.10},{\h});
  \else
    \fill[mechBlue]
      ({0.95+\i*\cw+0.05},2.86) rectangle ++({\cw-0.10},{\h});
  \fi
}
\draw[mechInk!55] (0.95,2.86) -- ++({12*\cw},0);
\node[anchor=east,font=\footnotesize] at (0.88,3.22) {$\sqrt{w_i}$};
\pgfmathsetseed{11}
\foreach \i in {0,...,11}{\foreach \j in {0,...,7}{
  \pgfmathsetmacro{\g}{int(8+86*rnd)}
  \ifnum\i=6
    \fill[mechRose!\g!mechPeach] ({0.95+\i*\cw},{2.70-(\j+1)*\ch}) rectangle ++(\cw,\ch);
  \else
    \fill[mechBlue!\g!mechCyan] ({0.95+\i*\cw},{2.70-(\j+1)*\ch}) rectangle ++(\cw,\ch);
  \fi
}}
\draw[mechInk!75] (0.95,2.70) rectangle ++({12*\cw},{-8*\ch});
\draw[mechRoseInk,line width=.7pt] ({0.95+6*\cw},2.70) rectangle ++(\cw,{-8*\ch});
\node[anchor=east,font=\footnotesize] at (0.88,1.68) {$G$};
\node[mechRoseInk,font=\scriptsize,anchor=north] at ({0.95+6.5*\cw},0.62) {sample $J$};
\node[anchor=north west,align=left,font=\scriptsize] at (0,0.16)
 {Columns $g_i$: iid $N(0,I_R)$.\\
  $R=\dim(\range V)$; $D=\diag(\sqrt{w_i})$.\\
  Scales $w_i>0$: independent of $G$, mean one.\\
  Typically $\max_{i\le m}w_i\asymp_d(\log m)^{d-1}$.};
\begin{scope}[shift={(6.55,2.62)}]
\node[anchor=west,font=\small\bfseries] at (0,2.48) {(b) Embedding sees the scales};
\node[anchor=west] at (0,1.93) {$V^\top\Omega\Omega^\top V=\tfrac1m\textstyle\sum_i w_i\,g_ig_i^\top$};
\def\sx{0.97}
\fill[mechLime!85!white] ({0.5*\sx},0.42) rectangle ({1.5*\sx},1.06);
\node[mechBlue,font=\scriptsize] at ({1.0*\sx},1.2) {$1\pm\varepsilon$};
\draw[->,mechInk,line width=.45pt] (0,0.42) -- ({8.55*\sx},0.42) node[right] {$\lambda$};
\foreach \t in {0,1,2,4,6,8}{\draw[mechInk!80] ({\t*\sx},0.42) -- ++(0,-0.07) node[below,font=\scriptsize] {$\t$};}
\foreach \l in {0.22,0.45,0.70,0.95,1.20,1.55,1.95}{\fill[mechBlue] ({\l*\sx},0.68) circle (2pt);}
\fill[mechRose] ({6.5*\sx},0.68) circle (2.5pt);
\node[mechRoseInk,anchor=south,font=\scriptsize,align=center] at ({6.5*\sx},0.82)
 {heaviest sample $J$:\\[-1pt] $\lambda_{\max}\ge w_J\norm{g_J}^2/m$};
\node[anchor=north west,font=\scriptsize,align=left] at (0,-0.17)
 {All eigenvalues must lie in the band.\\
  At fixed accuracy and confidence: $m\gtrsim_d R\log^{d-1}(eR)$.};
\end{scope}
\begin{scope}[shift={(6.55,-1.15)}]
\node[anchor=west,font=\small\bfseries] at (0,2.45) {(c) Approximation forgets them};
\coordinate (O1) at (0.42,0.66); \coordinate (O2) at (2.62,0.66);
\foreach \a in {18,52,95,140}{\draw[->,mechLavender,line width=.85pt] (O1) -- ++(\a:0.86);}
\foreach \a/\l/\c in {18/0.50/mechBlue,52/1.52/mechRose,95/0.64/mechBlue,140/0.33/mechBlue}
 {\draw[->,\c,line width=.85pt] (O2) -- ++(\a:\l);}
\node[font=\scriptsize] at (0.72,0.12) {columns of $G$};
\node[font=\scriptsize] at (2.98,0.12) {columns of $GD$};
\node at (1.72,0.99) {$\equiv$};
\node[anchor=north west,align=left,font=\scriptsize] at (3.92,2.06)
 {Same lines, so $\range(GD)=\range(G)$.\\
  If the input's whole row space lies in\\
  $\range V$, range-based outputs have the\\
  dense-Gaussian law. For target rank $r<R$,\\[2pt]
  $\E\norm{A-P_YA}_{\F}^2\le\bigl(1+\tfrac{r}{m-r-1}\bigr)\tau$,\\[1pt]
  for $m>r+1$. If $m\ge r+1+\lceil r/\varepsilon\rceil$,\\
  this is at most $(1+\varepsilon)\tau$.};
\end{scope}
\draw[->,mechInk!72,line width=.45pt] (5.47,2.3) to[out=25,in=190] (6.45,3.22);
\node[font=\scriptsize,text=mechInk,fill=white,inner sep=2pt] at (5.78,3.49) {Gram};
\draw[->,mechInk!72,line width=.45pt] (5.47,1.05) to[out=-20,in=175] (6.45,-0.49);
\node[font=\scriptsize,text=mechInk,fill=white,inner sep=2pt,anchor=west] at (5.94,0.60) {range};
\end{tikzpicture}%
}
\endgroup
\caption{The Gaussian common-factor mechanism. The Gram retains the heavy
column scales, while the range discards them
(Section~\ref{core:sec:lower}; Proposition~\ref{exact:thm:projective}).
Here $V$ has $R$ orthonormal columns, $r$ is the approximation rank,
$Y=A\Omega$, and $\tau=\norm{A-A_r}_{\F}^2$.
The exact output law requires the input's entire relevant support in
$\range V$, not merely its leading subspace. The weight scale is asymptotic
at fixed $d$; all depicted entries, weights, and eigenvalues are schematic.}
\label{fig:mechanism}
\end{figure}

\begin{corollary}[Relative approximation with an upper spectral obstruction]
\label{separation:cor:same-input}
Fix $d\ge2$ and $0<\varepsilon,\delta<1$. For either independent
Gaussian or independent-coordinate Rademacher factors there are fixed
inputs $A^{(r)}$ and sketch sizes $m_r=O_{d,\varepsilon,\delta}(r)$ such that,
writing $M_r=U_r^\top\Omega\Omega^\top U_r$,
\begin{align*}
 \Pp\bigl\{&\norm{A^{(r)}-\widehat A^{(r)}_r}_\F^2\le1+\varepsilon,\\
 &\lambda_{\min}(M_r)\ge\tfrac12,\quad\lambda_{\max}(M_r)\ge2\bigr\}
 \ge1-\delta-o(1).
\end{align*}
Here $U_r$ spans the leading right singular space, the optimal squared
rank-$r$ tail is one, and the input rank exceeds $m_r$. The same error
bound holds for the untruncated rangefinder on this event. In contrast,
the zero output has squared error $r+1$.
\end{corollary}
\begin{proof}
Both factor laws have directional fourth-moment constant at most $3^d$.
Choose $m_r$ by \eqref{app:eq:relative-budget} and $U_r$ from the
corresponding lower-bound construction. Enlarge unused mode coordinates
so that $N_r>\max\{m_r,r+1\}$, without changing the projected probe law.
With $P_r=U_rU_r^\top$, put
\[
 A^{(r)}=P_r+(N_r-r)^{-1/2}(I_{N_r}-P_r).
\]
Its first $r$ singular values are one, its remaining singular values are
$(N_r-r)^{-1/2}<1$, and its squared tail is one. In the proof of
Theorem~\ref{app:thm:relative}, the lower-injection event and the cross
moment event jointly have probability at least $1-\delta$; both error
bounds and $M_r\succeq I_r/2$ hold on this same event. Thus adding the
lower spectral edge does not spend further confidence. Since
$m_r/[r(\log r)^{d-1}]\to0$ and $m_r\ge r$, the Gaussian and Rademacher
lower-bound proofs each give $\lambda_{\max}(M_r)\ge2$ with probability
$1-o(1)$. Intersect the events by a union bound.
\end{proof}

For regression, the same separation follows by taking design $U_r$
and response $\sqrt r\,U_re_1+v_r$, where $v_r$ is a fixed unit vector
orthogonal to $\range U_r$: the optimal squared residual is one and
the zero-coefficient residual is $r+1$. For PSD trace approximation,
take $K_r=P_r+(N_r-r)^{-1}(I-P_r)$ and apply the range result to
$K_r^{1/2}$. Each is a separate fixed-input probability statement.

A complementary condition-number estimate at linear sketch size
is recorded in Appendix~\ref{appendix:spectral-extensions}.

\subsection{Exact laws from random lines}
Nonzero column scales leave sampled ranges unchanged. This observation,
also used for Gaussian and spherical probes in~\cite[Section~5.4]{CEMT2025},
gives exact output laws when the input's entire relevant support has the
required distribution of random lines. The comparison is with a dense
Gaussian sketch applied to the same input, whose geometry still shapes
the sampled range. In this section the right sketch continues to use
$\Omega=m^{-1/2}[z_1\ \cdots\ z_m]$.
\begin{proposition}[Exact Gaussian range laws]\label{exact:thm:projective}
Let $V:\R^R\to S\subseteq\R^N$ be an isometry. Suppose
$x=V^\top z\ne0$ almost surely and its unoriented line is uniform
on real projective space, meaning the image of the uniform sphere under
$u\sim-u$. Independent probes then admit a coupling
\begin{equation}\label{exact:eq:projective-coupling}
 V^\top\Omega=m^{-1/2}GD,
\end{equation}
where $G$ has iid standard Gaussian entries and $D$ is an almost surely
invertible diagonal matrix, possibly dependent on $G$.
For every fixed input with $\range(A^\top)\subseteq S$, its sampled
range, all fixed powered ranges, and their block Krylov range have the
dense-Gaussian joint law. Rangefinding and optimal fixed-rank outputs
have their Gaussian laws when selected as fixed measurable functions
of the input and the relevant sampled subspace, equivalently its
orthogonal projector. The same selection rule is used for both sketch
laws, including at ties; it does not inspect raw column scales. Equal
ranges always give equal optimal errors. The same convention gives
the output laws for Nystr\"om and its fixed-rank truncation when
$K\succeq0$ and $\range K\subseteq S$.

A sufficient condition, for a random environment $h$, is
\begin{equation}\label{exact:eq:conditional}
 V^\top z\mid h\sim N(0,w(h)I_R),\qquad w(h)>0
 \quad\hbox{almost surely}.
\end{equation}
Then $D=\diag(\sqrt{w(h_i)})$ can be chosen independent of $G$.
The Gaussian supports in Definition~\ref{exact:prop:supports} satisfy
this condition.
\end{proposition}
\begin{proof}
For independent fair signs $\epsilon_i$ and independent Gaussian radii
$\chi_i$, each with the norm distribution of a standard Gaussian in
$\R^R$, put $g_i=\chi_i\epsilon_i x_i/\norm{x_i}$ and
$d_i=\epsilon_i\norm{x_i}/\chi_i$. Uniformity of the unoriented
line makes $g_i$ standard Gaussian and $g_id_i=x_i$.
Independence across augmented probes proves the factorization.
Every deterministic map annihilating $S^\perp$ then has the same
right-hand invertible factor $D$, which preserves each range and the
concatenated block range. Their projectors therefore coincide under
the coupling, so the same measurable function of the input and the
projector gives the same output, including at ties. For PSD outputs use
$\mathcal N=K^{1/2}P_{K^{1/2}\Omega}K^{1/2}$.
Under the conditional condition, $g_i=x_i/\sqrt{w(h_i)}$ has a
constant Gaussian conditional law and hence is independent of the
environment, giving the stronger independence statement.
\end{proof}

\paragraph{Other factor laws.}
The proposition transfers existing Gaussian range guarantees
\cite{HMT2011,MuscoMusco2015}; it gives neither a new Gaussian bound
nor a floating-point conditioning law. The entire relevant input
support must satisfy the hypothesis.
More generally, on a common-factor support with independent factor laws,
\[
 V^\top\Omega=m^{-1/2}X\diag(s_i),\qquad
 X_i=B^\top g_{i1},\quad s_i=\prod_{j=2}^d\ip{b_j}{g_{ij}},
\]
where $B$ collects the orthonormal first-mode vectors. The range equals
that of the first-factor columns with $s_i\ne0$; conditional on these
indices the retained columns are iid with their original law, and
their count is binomial. A coordinate first-mode support with sign
factors gives dense sign columns. A general orthonormal projection of
sign factors need not have independent coordinates. Taking the flat
common-factor length odd in the Rademacher witness makes every $s_i$
nonzero, so no columns are lost, while preserving the logarithmic
obstruction for sufficiently large $r$.


\paragraph{Gaussian relative-error coefficients.}
Under Proposition~\ref{exact:thm:projective} on the entire relevant
support, the Gaussian moment identities~\cite[Propositions~10.1 and A.5]{HMT2011}
and the rank-$r$ competitor in the proof of Theorem~\ref{app:thm:relative}
give, for $m>r+1$,
\begin{align}
 \E\norm{A-P_{A\Omega}A}_\F^2,\quad
 \E\norm{A-\widehat A_r}_\F^2
 &\le\left(1+\frac{r}{m-r-1}\right)\norm{A-A_r}_\F^2,
 \label{exact:eq:relative-rect}\\
 \E\tr(K-\mathcal N),\quad\E\tr(K-\mathcal N_r)
 &\le\left(1+\frac{r}{m-r-1}\right)\tau_r^{\rm tr}(K).
 \label{exact:eq:relative-psd}
\end{align}
Thus $m=r+1+\lceil r/\varepsilon\rceil$ gives expected relative error
at most $1+\varepsilon$ for all four objectives; the PSD cases use the
$K^{1/2}$ correspondence in Corollary~\ref{app:cor:nystrom}.

The common-factor case gives an exact distributional counterpart to
Corollary~\ref{separation:cor:same-input}: the range-based approximation
outputs retain their dense-Gaussian laws while the probability of
violating any fixed upper embedding tolerance tends to one. The
relative-error bounds in the following proposition are expectation
bounds, distinct from the high-probability conclusion of that corollary.
\begin{proposition}[Exact approximation laws despite leading-subspace embedding failure]
\label{exact:prop:separation}
Fix tensor order $d\ge2$, accuracy $\varepsilon\in(0,1)$, and upper
distortion tolerance $\alpha>0$. For each $r$, let
$m=r+1+\lceil r/\varepsilon\rceil$, choose $R\ge m$, and let a fixed input
$A$ have full row support equal to an $R$-dimensional common-factor support
in Definition~\ref{exact:prop:supports}. Choose its leading $r$ right
singular vectors to be the first $r$ vectors of that basis. Gaussian product
probes have exactly the dense-Gaussian approximation laws and satisfy
\eqref{exact:eq:relative-rect}. Nevertheless, for that actual leading
isometry $V_1$, the normalized sketch obeys
\[
 \Pp\{\lambda_{\max}(V_1^\top\Omega\Omega^\top V_1)>1+\alpha\}
 \longrightarrow1\qquad(r\longrightarrow\infty).
\]
The same construction applies to a PSD input with this positive eigenspace,
so \eqref{exact:eq:relative-psd} coexists with the same head obstruction.
One may take $R>m$ to keep the approximation genuinely nonexact.
\end{proposition}
\begin{proof}
Exact laws and the approximation coefficients follow from the preceding
results. On the leading support the columns have form
$m^{-1/2}\sqrt{w_i}g_i$, with iid $g_i\sim N(0,\Id_r)$ independent of
iid $w_i=\prod_{j=2}^d\xi_{ij}^2$, where all $\xi_{ij}$ are standard
normal. Put ${c_{\rm cap}}=3+\varepsilon^{-1}$, so $m/r\le{c_{\rm cap}}$.
The random weight has unbounded support, hence
\[
 p_*=\Pp\{w_i>2{c_{\rm cap}}(1+\alpha)\}>0.
\]
Also $\Pp\{\norm{g_i}^2\ge r/2\}\to1$, by the variance $2r$ of its
squared norm. Therefore the probability $q_r$ that a column satisfies
both inequalities tends to $p_*>0$. Such a column alone gives
\[
 \lambda_{\max}(V_1^\top\Omega\Omega^\top V_1)
 \ge\frac{w_i\norm{g_i}^2}{m}>1+\alpha.
\]
The column events are independent, and $(1-q_r)^m\to0$. This is upper
distortion on the leading rank-$r$ space even though $m>r$; no ambient
rank-deficiency obstruction is involved.
\end{proof}


For comparison, a dense Gaussian sketch with the same
$m/r\to c_\varepsilon=1+\varepsilon^{-1}$ has
\[
 \lambda_{\max}(V_1^\top\Omega\Omega^\top V_1)
 \longrightarrow (1+c_\varepsilon^{-1/2})^2
\]
almost surely~\cite{YBK1988}. Thus failure of the upper tolerance
$1+\alpha$ is not specific to Khatri--Rao sketches when
$\alpha<(1+c_\varepsilon^{-1/2})^2-1$. The proposition distinguishes
the product law from the dense law for every fixed tolerance above
that edge: the product upper eigenvalue diverges, whereas the dense
one stays bounded. At $\varepsilon=1/2$ the dense edge is
$(1+1/\sqrt3)^2\approx2.488$.

\subsection{Regression under Gaussian scale mixtures}
\label{weighted:sec}
Column scales remain relevant to sketch-and-solve and generalized
Nystr\"om, as noted in~\cite[Section~5.4]{CEMT2025}.
The weighted-least-squares interpretation is classical for sampling
sketches~\cite{MaMahoneyYu2015}; related weighted covariance matrices
arise from elliptical data in~\cite{DobribanLiu2019}.
Here the sketch supplies the weights. If $A=V\bar A$ and $b=V\bar b$
both lie in a support satisfying \eqref{exact:eq:conditional}, then
\[
 \norm{\Omega^\top(Ax-b)}^2
 =\frac1m\sum_iw_i\ip{g_i}{\bar A x-\bar b}^2,
\]
with weights independent of the Gaussian $g_i$. Changing the weights
preserves sampled ranges but changes this objective. Consequently the
range-containment confidence argument does not apply to raw regression.

\begin{proposition}[Conditional regression law and a confidence certificate]
\label{weighted:prop}
Draw iid probes with $V^\top z=\sqrt w\,g$, where $g$ is standard Gaussian,
$w>0$ independent of $g$, $\E w=1$, and
$\mu_2=\E w^2<\infty$. Fix $A,b$ with
$\operatorname{span}(\range A,b)\subseteq S$, $\rank A=r$, and
nonzero optimal residual $e=b-Ax_*$.
Let $\widetilde x$ be the sketched least-squares solution and put
\[
 \mathcal E_m^{\rm reg}=\frac{\norm{A\widetilde x-b}^2}{\norm e^2}-1.
\]
For $m\ge r$, there are iid $\xi_i\sim N(0,I_r)$ and
independent $\eta_i\sim N(0,1)$, independent of the weights, such that
\begin{equation}\label{weighted:eq:law}
 \mathcal E_m^{\rm reg}=\norm{a_m}^2,\quad
 a_m=S_m^{-1}\sum_iw_i\xi_i\eta_i,\quad
 S_m=\sum_iw_i\xi_i\xi_i^\top.
\end{equation}
Conditional on $(w_i,\xi_i)_{i\le m}$, $a_m$ is centered Gaussian with
\begin{equation}\label{weighted:eq:covariance}
 C_m=S_m^{-1}\Bigl(\sum_iw_i^2\xi_i\xi_i^\top\Bigr)S_m^{-1},
 \qquad \E[\mathcal E_m^{\rm reg}\mid(w_i,\xi_i)]=\tr C_m.
\end{equation}
For every $t>0$, the conditional Gaussian quadratic-form bound gives
\begin{equation}\label{weighted:eq:conditional-tail}
 \Pp\left\{\mathcal E_m^{\rm reg}>\tr C_m+
       2\sqrt{t\tr(C_m^2)}+2t\norm{C_m}
       \;\middle|\,(w_i,\xi_i)_{i\le m}\right\}\le e^{-t}.
\end{equation}
For fixed $r$ and a fixed weight law, as $m\to\infty$,
\[
 mC_m\longrightarrow\mu_2I_r\quad\hbox{almost surely},\qquad
 m\mathcal E_m^{\rm reg}\ \Rightarrow\ \mu_2\chi_r^2.
\]
In particular, the scaled conditional mean tends to $r\mu_2$.

For $0<\varepsilon,\delta<1$, let $b_m(\delta)$ be deterministic with
$\Pp\{\max_iw_i>b_m(\delta)\}\le\delta/4$, and put $C_w=3\mu_2$.
Then $\Pp\{\mathcal E_m^{\rm reg}>\varepsilon\}\le\delta$ whenever
\begin{equation}\label{weighted:eq:budget}
 m\ge\max\left\{\mathfrak m_{\rm inj}(C_w,r;1/2,\delta/2),
 \frac{2b_m(\delta)}{\varepsilon}
     \bigl(\sqrt r+\sqrt{2\log(4/\delta)}\bigr)^2\right\}.
\end{equation}
\end{proposition}
The normal equations identify $C_m$; the Gaussian quadratic-form
bound retains its full spectrum. Bounding $C_m$ by
$2b_m(\delta)I_r/m$ gives the deterministic certificate above.
Appendix~\ref{appendix:weighted} proves these claims and the limit.

For the common-factor Gaussian law, the certificate uses
\[
 \mu_2=3^{d-1},\qquad
 b_m(\delta)=\bigl[2\log(8(d-1)m/\delta)\bigr]^{d-1}.
\]
For order-two common-Schmidt weights, with
$\rho=\sum_b\lambda_b^2$ and $L=\max_b\lambda_b$, it uses
\[
 \mu_2=1+2\rho,\qquad
 b_m(\delta)=1+2\sqrt{\rho\log(4m/\delta)}+2L\log(4m/\delta).
\]
Both give polylogarithmic confidence cost; for the common-factor law,
a sufficient explicit order is
\[
 m\ge C_d\frac{r+\log(1/\delta)}{\varepsilon}
 \left[\log\frac{e(r+\log(1/\delta))}{\varepsilon\delta}\right]^{d-1}.
\]

The order-two Gaussian common-Schmidt model summarizes the distinction:
\begin{center}\small
\begin{tabularx}{\linewidth}{@{}lY@{}}
\toprule Task & Role of the shared weight\\\midrule
Range-based approximation & Nonzero scales cancel; the Gaussian output law
requires the entire relevant support.\\
Regression, fixed $r$ & The limiting variance factor is $\mu_2=1+2\rho$.\\
Embedding & The width retains an extreme-row contribution governed by $L$.\\
\bottomrule
\end{tabularx}
\end{center}
The regression limit does not assert the finite-sample identity
$\E\mathcal E_m^{\rm reg}=r\mu_2/m$ or convergence of unconditional means.
Its confidence certificate still needs weight-tail control; arbitrary
mixing weights need not give polylogarithmic cost. The fourth-moment
obstruction in Proposition~\ref{app:prop:regression-lower} remains
applicable outside these Gaussian support classes.

\subsection{High confidence under sub-Gaussian factors}
\label{app:sec:subgauss-confidence}

Fourth moments alone give the inverse-confidence cost in
Theorem~\ref{app:thm:relative}. Higher directional moments give a
polylogarithmic confidence certificate for the same raw regression solve.
The underlying Hilbert-space moment principle is classical
\cite{Pinelis1994,Latala1997}; the following formulation keeps the
tensor moment budgets and the orthogonal interaction explicit.

Fix an integer $q\ge1$. Let $A_j$ be any valid directional moment bound
$\E|\ip u z|^{2j}\le A_j\norm u^{2j}$, for $2\le j\le2q$.
One may use Lemma~\ref{core:lem:factor} or the law-specific profiles
in Appendix~\ref{core:sec:law-profiles}. Let $\Gamma_{\perp}>0$ bound
the mixed fourth moment of any two orthogonal unit directions.
Choose an aggregate budget $B_r$ as in
\eqref{app:eq:aggregate-budget}, taking its minimum with $r\Gamma_{\perp}$
if necessary. One may take $\Gamma_{\perp}=\gamma(\Cfour)$ whenever that
number is positive, or use the coordinate or Gaussian constants in Appendix~\ref{appendix:wick}.
If $B_r=0$, every orthogonal interaction vanishes almost surely and
only the lower-injection condition below is needed; omit the formulas
dividing by $B_r$. Otherwise set
\begin{equation}\label{app:eq:hilbert-budgets}
 \nu_H=B_r,\qquad b_j^H=r^{j/2}A_j\ (3\le j\le2q),\qquad
 \Rmom_q^H=\max_{3\le j\le2q}
       \left(\frac{2b_j^H}{j!\nu_H}\right)^{1/(j-2)}.
\end{equation}
For $q=1$ the empty maximum is zero. Write $c_q=\gmom_q^{1/(2q)}$
and $h_q$ as in Theorem~\ref{closed:thm:moment}.

\begin{theorem}[A confidence certificate for the raw solve]
\label{app:thm:subgauss-confidence}
Assume \eqref{app:eq:fourth}, fix $r\ge1$ and $\varepsilon,\delta\in(0,1)$,
and suppose the directional budgets above hold. The rangefinding,
rank-$r$, and regression conclusions of Theorem~\ref{app:thm:relative}
hold with probability at least $1-\delta$ if
\begin{align}\label{app:eq:high-confidence}
 m&\ge\mathfrak m_{\rm inj}(\Cfour,r;1/2,\delta/2),\nonumber\\
 (2/\delta)^{1/(2q)}c_q
 \left(\sqrt{\nu_H/m}+h_q\Rmom_q^H/m\right)
 &\le\sqrt\varepsilon/2.
\end{align}
In particular, in the sub-Gaussian tensor model, with
$t=\log(2/\delta)$, a sufficient sample count is
\begin{equation}\label{app:eq:high-confidence-rate}
 m\ge C_{d,\Ksg}\left[
       \frac{rt}{\varepsilon}
       +\sqrt{\frac r\varepsilon}\,t^d\right].
\end{equation}
This statement concerns each fixed input and has no dependence on
the number of response columns.
\end{theorem}
Appendix~\ref{appendix:hilbert} proves the certificate and gives a
second comparison that keeps the weak Hilbert variance separately.
Its leading term is at most $\sqrt{(B_r+\Gamma_{\perp}(q-1))/m}$ before the
confidence multiplier, compared with
$\gmom_q^{1/(2q)}\sqrt{B_r/m}$ above. The two higher-order radii
need not be ordered, so the smaller valid allowance may be used.

\subsection{An angular-moment bound for the full sampled range}
\label{app:sec:angular-drift}

For Gaussian or spherical product columns, a complementary expected-error
bound uses moments near order two rather than directional fourth moments.
The use of entropy-scale moments for product probes is motivated by
Holden's product-probe note~\cite{Holden2026}; the scalar residual
recursion follows the RPCholesky analysis of
Chen--Epperly--Tropp--Webber~\cite{CETW2025}. We give the fractional-moment
conversion explicitly. The normalization step is the Gaussian/spherical
range invariance recorded in~\cite[Section~5.4]{CEMT2025}.

For positive integer mode dimensions $n_1,\ldots,n_d$, define
\begin{equation}\label{app:eq:angular-kappa}
 \begin{split}
 c_n^{\rm sph}(p)&=\frac{n^p\Gamma(p+1/2)\Gamma(n/2)}
                    {\Gamma(1/2)\Gamma(n/2+p)},\\
 H_{\rm ang}(p)&=\prod_{j=1}^d c_{n_j}^{\rm sph}(p),\\
 \varkappa_{\rm ang}
 &=\exp\!\left[\sum_{j=1}^d
   \{\log n_j+\psi(3/2)-\psi(n_j/2+1)\}\right],
 \end{split}
\end{equation}
where $p\ge1$ is real and $\psi=(\log\Gamma)'$.
In particular, $c_n^{\rm sph}(1)=1$, and a one-dimensional mode contributes one
to $\varkappa_{\rm ang}$.

\begin{theorem}[Expected approximation of the full sampled range]
\label{app:thm:angular-range}
Let $A$ be fixed, let the columns of $\Omega$ be independent real
Gaussian product probes or independent product probes whose mode-$j$
factor is uniform on the sphere of radius $\sqrt{n_j}$, and put
$P=P_{\range(A\Omega)}$.
Fix $1\le r<\rank A$ and $0<\varepsilon\le1$, and set
\[
 T=\norm A_\F^2,\qquad
 \tau=\norm{A-A_r}_\F^2,\qquad
 u_0=T/\tau-1.
\]
If $u_0\le\varepsilon$, the desired bound holds already for the empty
range. Otherwise, the explicit condition
\begin{equation}\label{app:eq:angular-budget-exact}
 m\ge\left\lceil\varkappa_{\rm ang}r
 \left[\frac1\varepsilon-\frac1{u_0}
                     +\log\frac{u_0}{\varepsilon}\right]\right\rceil
\end{equation}
implies
\begin{equation}\label{app:eq:angular-range-error}
 \E\norm{(I-P)A}_\F^2\le(1+\varepsilon)\tau.
\end{equation}
In particular, a simpler sufficient condition is
\begin{equation}\label{app:eq:angular-budget-simple}
 m\ge\left\lceil\varkappa_{\rm ang}r
 \left[\varepsilon^{-1}+\log(\varepsilon^{-1})
                              +\log(T/\tau)\right]\right\rceil.
\end{equation}
For PSD $K$, the same budgets, with $T=\tr K$ and
$\tau=\sum_{i>r}\lambda_i(K)$, give
$\E\tr(K-\mathcal N_K(\Omega))\le(1+\varepsilon)\tau$ for the full
Nystr\"om approximation.
\end{theorem}

The proof is in Appendix~\ref{appendix:angular}.
The conclusion concerns the entire sampled range and expectation.
Its constant $\varkappa_{\rm ang}$ can be smaller than the fourth-moment
constant, but its budget contains the trace-to-tail ratio.
The theorem does not assert the same estimate for a rank-$r$ truncation,
for regression, or for the Gram distortion of an unnormalized sketch.

The angular constant obeys
\[
 \varkappa_{\rm ang}\le(e^{2-\gamma_{\rm E}}/2)^d
 =(2.0743278106\ldots)^d.
\]
It improves the spherical fourth-moment product when some mode has
dimension greater than one. Appendix~\ref{appendix:angular} also gives
the corresponding Frobenius estimate for the full PSD residual.

\section{Numerical illustration}\label{sec:numerics}
The experiments compare the exact range-law mechanism with approximation
on inputs outside those support classes. All approximation inputs have
rank $R$, with $r$ singular values equal to one and the remaining
$R-r$ equal to $(R-r)^{-1/2}$. Thus the optimal squared rank-$r$
tail is one. We plot
$\|A-\widehat A_r\|_\F^2/\tau_r^{\rm F}(A)$, where $\widehat A_r$
is the optimal rank-$r$ output in the sampled range.
Lines are medians and bands are empirical 10--90\% trial quantiles.

\subsection{Exact range-law separation}
Figure~\ref{fig:kr-separation-numerics} uses a common-factor support,
$m=3r+1$, and $R=2m$. The upper Gram edge increases with tensor
order while the coupled Gaussian approximation outputs coincide,
as Proposition~\ref{exact:thm:projective} predicts.
At $r=2048$, median upper eigenvalues are $2.48$ for dense Gaussian
and $6.34,22.77,51.09$ for product orders two, three, and four.
The approximation ratios remain close to $1.30$ over the larger tested
ranks. The sign experiment uses the nonzero common-factor construction
in the paragraph following Proposition~\ref{exact:thm:projective};
its range has the dense-sign law.

\begin{figure}[htbp]
\centering\begingroup
\definecolor{plotNavy}{HTML}{29234F}
\definecolor{plotCyan}{HTML}{00B9D8}
\definecolor{plotBlue}{HTML}{254BE8}
\definecolor{plotPink}{HTML}{EA5CB5}
\definecolor{plotOrange}{HTML}{D07820}
\definecolor{plotReference}{HTML}{7042DF}
\begin{tikzpicture}
\begin{groupplot}[
group style={group size=2 by 1,horizontal sep=1.65cm},
width=.385\linewidth,height=5.1cm,scale only axis,
xmode=log,log basis x=2,
grid=major,grid style={gray!18},
tick label style={font=\scriptsize},
label style={font=\small},title style={font=\small\bfseries},
legend style={font=\scriptsize,draw=none,fill=white},
]
\nextgroupplot[title={(a) Upper spectral edge},
ymode=log,xmin=7,xmax=2280,
xtick={8,32,128,512,2048},
xlabel={Target rank $r$},ylabel={$\lambda_{\max}$},
legend to name=cleanExactLegend,legend columns=4]
\addplot[draw=none,name path=exactEdge0lo,forget plot] coordinates {(8,1.708854013) (16,2.006781246) (32,2.175027096) (64,2.295153031) (128,2.376591117) (256,2.405982145) (512,2.440305992) (1024,2.457388889) (2048,2.469554878)};
\addplot[draw=none,name path=exactEdge0hi,forget plot] coordinates {(8,2.450482509) (16,2.458648422) (32,2.501047732) (64,2.524388492) (128,2.493168474) (256,2.491923983) (512,2.494770527) (1024,2.489234461) (2048,2.485664156)};
\addplot[draw=none,fill=plotNavy,fill opacity=.16,forget plot] fill between[of=exactEdge0lo and exactEdge0hi];
\addplot[color=plotNavy,thick,mark=*,mark size=1.45pt] coordinates {(8,2.054342477) (16,2.226030684) (32,2.345938139) (64,2.386013125) (128,2.432630156) (256,2.455963027) (512,2.463903017) (1024,2.470437436) (2048,2.477092791)};
\addlegendentry{Dense Gaussian}
\addplot[draw=none,name path=exactEdge2lo,forget plot] coordinates {(8,1.833213468) (16,2.424972836) (32,3.223646065) (64,3.598588157) (128,4.084572867) (256,4.552980529) (512,5.013554865) (1024,5.383333205) (2048,5.839738417)};
\addplot[draw=none,name path=exactEdge2hi,forget plot] coordinates {(8,4.653837694) (16,4.971973443) (32,5.152550413) (64,5.768993285) (128,5.720734422) (256,6.207039731) (512,6.566730089) (1024,7.086418433) (2048,7.012639464)};
\addplot[draw=none,fill=plotCyan,fill opacity=.16,forget plot] fill between[of=exactEdge2lo and exactEdge2hi];
\addplot[color=plotCyan,thick,mark=square*,mark size=1.45pt] coordinates {(8,3.015665264) (16,3.485708049) (32,4.387596975) (64,4.438160535) (128,4.840880527) (256,5.197357666) (512,5.560567035) (1024,6.016616348) (2048,6.335125769)};
\addlegendentry{KR, $d=2$}
\addplot[draw=none,name path=exactEdge3lo,forget plot] coordinates {(8,1.242514858) (16,2.842603494) (32,3.921900244) (64,5.155846153) (128,6.727145565) (256,8.623616189) (512,11.5664574) (1024,14.2430184) (2048,16.44258277)};
\addplot[draw=none,name path=exactEdge3hi,forget plot] coordinates {(8,9.533349056) (16,11.80933733) (32,14.21820273) (64,12.61784602) (128,19.14022721) (256,19.42954434) (512,25.61610776) (1024,26.42513265) (2048,31.87465007)};
\addplot[draw=none,fill=plotBlue,fill opacity=.16,forget plot] fill between[of=exactEdge3lo and exactEdge3hi];
\addplot[color=plotBlue,thick,mark=triangle*,mark size=1.45pt] coordinates {(8,3.500292472) (16,4.866101183) (32,6.977388949) (64,8.007479714) (128,9.98027985) (256,12.39873617) (512,15.2856468) (1024,17.86359945) (2048,22.76986472)};
\addlegendentry{KR, $d=3$}
\addplot[draw=none,name path=exactEdge4lo,forget plot] coordinates {(8,0.9767603243) (16,2.351474911) (32,3.91296098) (64,5.777319978) (128,10.77179862) (256,12.06695672) (512,21.67823471) (1024,24.97295701) (2048,39.79708637)};
\addplot[draw=none,name path=exactEdge4hi,forget plot] coordinates {(8,10.78354528) (16,18.82751165) (32,22.55184206) (64,26.46286049) (128,35.3625524) (256,47.47817547) (512,60.82725863) (1024,87.75824618) (2048,80.88388904)};
\addplot[draw=none,fill=plotPink,fill opacity=.16,forget plot] fill between[of=exactEdge4lo and exactEdge4hi];
\addplot[color=plotPink,thick,mark=diamond*,mark size=1.45pt] coordinates {(8,2.795998244) (16,5.531875055) (32,8.281505425) (64,13.21061632) (128,16.44684636) (256,23.34302362) (512,32.03709954) (1024,40.78329986) (2048,51.09065116)};
\addlegendentry{KR, $d=4$}
\addplot[plotNavy,dashed,forget plot] coordinates {(8,2.4880338717)(2048,2.4880338717)};
\nextgroupplot[title={(b) Rank-$r$ approximation},
xmin=7,xmax=570,xtick={8,32,128,512},ymin=1.16,ymax=1.54,
xlabel={Target rank $r$},ylabel={Relative squared error},
legend style={at={(.98,.04)},anchor=south east,font=\scriptsize,draw=none,fill=white}]
\addplot[draw=none,name path=exactError0lo,forget plot] coordinates {(8,1.215139509) (16,1.247372641) (32,1.272342752) (64,1.289871563) (128,1.29507693) (256,1.296772022) (512,1.298035476)};
\addplot[draw=none,name path=exactError0hi,forget plot] coordinates {(8,1.321372943) (16,1.318889351) (32,1.307481329) (64,1.307106619) (128,1.302558382) (256,1.301661358) (512,1.300409622)};
\addplot[draw=none,fill=plotNavy,fill opacity=.16,forget plot] fill between[of=exactError0lo and exactError0hi];
\addplot[color=plotNavy,thick,mark=*,mark size=1.45pt] coordinates {(8,1.271608868) (16,1.28256323) (32,1.293382375) (64,1.29802156) (128,1.298951943) (256,1.299264378) (512,1.299519979)};
\addlegendentry{Gaussian / KR Gaussian}
\addplot[draw=none,name path=exactError1lo,forget plot] coordinates {(8,1.225043302) (16,1.248037724) (32,1.277708057) (64,1.287197993) (128,1.29390286) (256,1.296570258) (512,1.298056452)};
\addplot[draw=none,name path=exactError1hi,forget plot] coordinates {(8,1.294409013) (16,1.300791989) (32,1.302069377) (64,1.301333227) (128,1.301206487) (256,1.300910119) (512,1.300135466)};
\addplot[draw=none,fill=plotOrange,fill opacity=.16,forget plot] fill between[of=exactError1lo and exactError1hi];
\addplot[color=plotOrange,thick,mark=square*,mark size=1.45pt] coordinates {(8,1.255094511) (16,1.277281286) (32,1.287715877) (64,1.293873216) (128,1.296372093) (256,1.298571764) (512,1.299252626)};
\addlegendentry{Rademacher}
\addplot[plotReference,dashed,forget plot] coordinates {(8,1.5)(512,1.5)};
\end{groupplot}
\node[anchor=south] at ([yshift=.75cm]group c1r1.north east)
 {\pgfplotslegendfromname{cleanExactLegend}};
\end{tikzpicture}\endgroup
\caption{Common-factor supports at $m=3r+1$.
The Gaussian product laws share their first-mode Gaussian matrix, so their
range outputs coincide. The Rademacher curve uses nonzero common signs.
Dashed lines mark the asymptotic dense-Gaussian upper edge
\cite{YBK1988} and the Gaussian expected-error bound $1.5$.
Panel (a) uses 96 trials through $r=512$, 64 at $r=1024$, and 32 at
$r=2048$; panel (b) uses 64 through $r=256$ and 32 at $r=512$.}
\label{fig:kr-separation-numerics}
\end{figure}

\subsection{Beyond exact range-law supports}
For Figure~\ref{fig:generic-approximation}, the leading right singular
space is the common-factor witness, and the tail is a uniformly sampled
subspace of its orthogonal complement, fixed before sketching.
All laws share this input, and each sketch is evaluated separately.
The full support therefore need not satisfy the exact Gaussian-law
hypothesis.

Panel (a) increases $r$ at $m=3r+1$ and $R=2m$.
At $r=1024$, median relative errors are $1.300$, $1.522$, $1.805$,
and $1.948$ for dense Gaussian and product orders two, three, and four.
Panel (b) instead fixes the input and increases $m$, always with
$m/R\le1/4$. On its tested grid, the empirical 90th-percentile error
first falls below $1.5$ at $m/r=3,6,6,8$, respectively.
The panels use different mode dimensions and input ranks, specified in
the caption; their error levels are not directly interchangeable.
These are fixed-input observations, not uniform sample-complexity bounds.

\begin{figure}[htbp]
\centering\begingroup
\definecolor{plotNavy}{HTML}{29234F}
\definecolor{plotCyan}{HTML}{00B9D8}
\definecolor{plotBlue}{HTML}{254BE8}
\definecolor{plotPink}{HTML}{EA5CB5}
\definecolor{plotOrange}{HTML}{D07820}
\definecolor{plotReference}{HTML}{7042DF}
\begin{tikzpicture}
\begin{groupplot}[
group style={group size=2 by 1,horizontal sep=1.65cm},
width=.385\linewidth,height=5.1cm,scale only axis,
xmode=log,log basis x=2,
grid=major,grid style={gray!18},
tick label style={font=\scriptsize},
label style={font=\small},title style={font=\small\bfseries},
legend style={font=\scriptsize,draw=none,fill=white},
]
\nextgroupplot[title={(a) Increasing target rank},
xmin=7,xmax=1150,xtick={8,32,128,512,1024},
xlabel={Target rank $r$ ($m=3r+1$)},ylabel={Relative squared error},
legend to name=cleanGeneralLegend,legend columns=4]
\addplot[draw=none,name path=generalRank0lo,forget plot] coordinates {(8,1.208983704) (16,1.253059328) (32,1.273531492) (64,1.284717406) (128,1.29441211) (256,1.296617489) (512,1.298154927) (1024,1.299351516)};
\addplot[draw=none,name path=generalRank0hi,forget plot] coordinates {(8,1.325272665) (16,1.330247469) (32,1.309480332) (64,1.303568437) (128,1.30270462) (256,1.301202123) (512,1.300051924) (1024,1.300333535)};
\addplot[draw=none,fill=plotNavy,fill opacity=.16,forget plot] fill between[of=generalRank0lo and generalRank0hi];
\addplot[color=plotNavy,thick,mark=*,mark size=1.45pt] coordinates {(8,1.258472267) (16,1.287007051) (32,1.292103679) (64,1.294270275) (128,1.298417569) (256,1.299432648) (512,1.299423168) (1024,1.299716507)};
\addlegendentry{Dense Gaussian}
\addplot[draw=none,name path=generalRank2lo,forget plot] coordinates {(8,1.256201931) (16,1.321152601) (32,1.386487686) (64,1.43059529) (128,1.462410632) (256,1.466256433) (512,1.486503659) (1024,1.503807802)};
\addplot[draw=none,name path=generalRank2hi,forget plot] coordinates {(8,1.645645468) (16,1.636739619) (32,1.60664392) (64,1.575640565) (128,1.57239909) (256,1.554045572) (512,1.546628817) (1024,1.539581437)};
\addplot[draw=none,fill=plotCyan,fill opacity=.16,forget plot] fill between[of=generalRank2lo and generalRank2hi];
\addplot[color=plotCyan,thick,mark=square*,mark size=1.45pt] coordinates {(8,1.408974007) (16,1.453673459) (32,1.483921975) (64,1.489651323) (128,1.504168969) (256,1.504615725) (512,1.510060682) (1024,1.521898689)};
\addlegendentry{KR, $d=2$}
\addplot[draw=none,name path=generalRank3lo,forget plot] coordinates {(8,1.301705687) (16,1.431774361) (32,1.549624573) (64,1.598349845) (128,1.661231558) (256,1.709339391) (512,1.741648203) (1024,1.780166874)};
\addplot[draw=none,name path=generalRank3hi,forget plot] coordinates {(8,1.996383272) (16,2.062512655) (32,2.046559511) (64,1.945461573) (128,1.909327364) (256,1.900762894) (512,1.873055415) (1024,1.863101281)};
\addplot[draw=none,fill=plotBlue,fill opacity=.16,forget plot] fill between[of=generalRank3lo and generalRank3hi];
\addplot[color=plotBlue,thick,mark=triangle*,mark size=1.45pt] coordinates {(8,1.57285907) (16,1.682765084) (32,1.731598069) (64,1.754327837) (128,1.786320718) (256,1.798882564) (512,1.790751141) (1024,1.805494834)};
\addlegendentry{KR, $d=3$}
\addplot[draw=none,name path=generalRank4lo,forget plot] coordinates {(8,1.334173086) (16,1.473040992) (32,1.621236767) (64,1.720072116) (128,1.80039841) (256,1.815081239) (512,1.878424627) (1024,1.893656612)};
\addplot[draw=none,name path=generalRank4hi,forget plot] coordinates {(8,2.070471879) (16,2.397736204) (32,2.100747671) (64,2.179078517) (128,2.193842981) (256,2.08711893) (512,2.055499458) (1024,2.000725614)};
\addplot[draw=none,fill=plotPink,fill opacity=.16,forget plot] fill between[of=generalRank4lo and generalRank4hi];
\addplot[color=plotPink,thick,mark=diamond*,mark size=1.45pt] coordinates {(8,1.613089984) (16,1.770945382) (32,1.848814657) (64,1.895841385) (128,1.956630211) (256,1.941483301) (512,1.977436322) (1024,1.947677282)};
\addlegendentry{KR, $d=4$}
\addplot[plotReference,dashed,forget plot] coordinates {(8,1.5)(1024,1.5)};
\nextgroupplot[title={(b) Increasing sketch size},
ymode=log,xmin=1.4,xmax=13,
xtick={1.5,2,3,4,6,8,12},xticklabels={1.5,2,3,4,6,8,12},
ytick={1,1.5,2,3,5,10,15},yticklabels={1,1.5,2,3,5,10,15},
xlabel={Oversampling ratio $m/r$},ylabel={Relative squared error}]
\addplot[draw=none,name path=generalWidth0lo,forget plot] coordinates {(1.5,2.623929591) (2,1.866180375) (3,1.449261013) (4,1.29629855) (6,1.173412443) (8,1.118102388) (12,1.068159213)};
\addplot[draw=none,name path=generalWidth0hi,forget plot] coordinates {(1.5,2.982910627) (2,2.00744497) (3,1.492352253) (4,1.322123771) (6,1.182591272) (8,1.124571605) (12,1.070809551)};
\addplot[draw=none,fill=plotNavy,fill opacity=.16,forget plot] fill between[of=generalWidth0lo and generalWidth0hi];
\addplot[color=plotNavy,thick,mark=*,mark size=1.45pt,forget plot] coordinates {(1.5,2.762077131) (2,1.94601936) (3,1.473445315) (4,1.309562019) (6,1.177623291) (8,1.121080214) (12,1.069591875)};
\addplot[draw=none,name path=generalWidth2lo,forget plot] coordinates {(1.5,4.536497199) (2,2.653259411) (3,1.66644628) (4,1.39690463) (6,1.205592886) (8,1.131224401) (12,1.073308513)};
\addplot[draw=none,name path=generalWidth2hi,forget plot] coordinates {(1.5,6.513207852) (2,3.635624166) (3,2.013685884) (4,1.564807352) (6,1.272843293) (8,1.171741319) (12,1.090096733)};
\addplot[draw=none,fill=plotCyan,fill opacity=.16,forget plot] fill between[of=generalWidth2lo and generalWidth2hi];
\addplot[color=plotCyan,thick,mark=square*,mark size=1.45pt,forget plot] coordinates {(1.5,5.407550004) (2,3.149585811) (3,1.818296343) (4,1.47081239) (6,1.238834803) (8,1.151671321) (12,1.08025989)};
\addplot[draw=none,name path=generalWidth3lo,forget plot] coordinates {(1.5,6.728488136) (2,4.0654756) (3,2.218759032) (4,1.656123772) (6,1.308601244) (8,1.183265818) (12,1.089086126)};
\addplot[draw=none,name path=generalWidth3hi,forget plot] coordinates {(1.5,10.9430301) (2,6.916513005) (3,3.352492365) (4,2.126534262) (6,1.463687585) (8,1.257495603) (12,1.121382764)};
\addplot[draw=none,fill=plotBlue,fill opacity=.16,forget plot] fill between[of=generalWidth3lo and generalWidth3hi];
\addplot[color=plotBlue,thick,mark=triangle*,mark size=1.45pt,forget plot] coordinates {(1.5,8.836193129) (2,5.400659515) (3,2.666329606) (4,1.879623661) (6,1.373368541) (8,1.219007537) (12,1.104009683)};
\addplot[draw=none,name path=generalWidth4lo,forget plot] coordinates {(1.5,9.797748979) (2,6.316358348) (3,3.017538979) (4,2.006116488) (6,1.430779492) (8,1.240023283) (12,1.108460751)};
\addplot[draw=none,name path=generalWidth4hi,forget plot] coordinates {(1.5,14.49485132) (2,10.14045913) (3,5.082408976) (4,2.947042026) (6,1.724720575) (8,1.389492762) (12,1.161067556)};
\addplot[draw=none,fill=plotPink,fill opacity=.16,forget plot] fill between[of=generalWidth4lo and generalWidth4hi];
\addplot[color=plotPink,thick,mark=diamond*,mark size=1.45pt,forget plot] coordinates {(1.5,11.6522936) (2,7.737847362) (3,3.860489956) (4,2.397451059) (6,1.57575852) (8,1.311476109) (12,1.136964619)};
\addplot[plotReference,dashed,forget plot] coordinates {(1.5,1.5)(12,1.5)};
\end{groupplot}
\node[anchor=south] at ([yshift=.75cm]group c1r1.north east)
 {\pgfplotslegendfromname{cleanGeneralLegend}};
\end{tikzpicture}\endgroup
\caption{Approximation with a common-factor head
and a generic orthogonal tail. Panel (a) uses $N=16r$, $R=2m$, and
common modes $(16)$, $(4,4)$, and $(2,2,4)$ for orders two, three,
and four. There are 128 trials through $r=64$, 96 at $128$, 64 at
$256$, and 32 at $512$ and $1024$.
Panel (b) fixes $r=32$, $N=16384$, and $R=1536$, with common modes
$(512)$, $(16,32)$, and $(8,8,8)$; 96 trials use nested sketch widths.
The dashed level $1.5$ is an error target; it also bounds Gaussian
expected error in panel (a).}
\label{fig:generic-approximation}
\end{figure}

\paragraph{Finite certificates.}
Table~\ref{tab:current-certificates} compares sufficient embedding widths
at $r=n_j=8$, $\varepsilon=1/2$, and $\delta=1/10$.
The closed bound uses Proposition~\ref{core:prop:certificate} with
$\Ksg^2=8/3$, selecting $1\le q\le200$.
The finite-profile calculation uses the law-specific moments in
Appendix~\ref{core:sec:law-profiles} and the centered recurrence
\eqref{core:eq:refined-discrete}, selecting $1\le q\le6$.
Rational budgets and upward rounding certify the latter counts, which
are rounded up to multiples of ten. These widths guarantee embedding
of every fixed subspace; they are not empirical thresholds.

\begin{table}[htbp]
\centering\small
\renewcommand{\arraystretch}{1.15}
\begin{tabular}{@{}lrr@{}}
\toprule
Certificate / factor law & $d=2$ & $d=3$\\
\midrule
Closed, $\Ksg^2=8/3$ & $14\,715$ & $149\,890$\\
Finite profile, Gaussian & $1\,540$ & $5\,930$\\
Finite profile, spherical & $910$ & $2\,570$\\
Finite profile, independent signs & $1\,360$ & $4\,090$\\
\bottomrule
\end{tabular}
\caption{Sufficient embedding widths at fixed accuracy and confidence.
Moment orders are selected deterministically.}
\label{tab:current-certificates}
\end{table}

\paragraph{Code and data availability.}
Code, raw trials, and certificate calculations are available at
\url{https://github.com/DiarHaidary/khatri-rao-numerics}.

\FloatBarrier


\section{Discussion and open questions}\label{sec:discussion}
The sharp embedding cost and the linear-size approximation guarantee
reflect different requirements on the same sampled matrix. The
remaining questions concern the parameters and subspace geometries
for which these guarantees can be sharpened.

\paragraph{Accuracy and confidence.}
The matching bounds determine rank dependence at fixed parameters;
they do not determine the jointly optimal dependence on
$r,\varepsilon,\delta,d$, and $\Ksg$.
Theorem~\ref{variance:thm:refinement} removes the variance rank
logarithm when $\delta\ge2/r$. At smaller failure probabilities,
the two upper bounds trade a variance logarithm against an additional
confidence logarithm in the large-row term. Closing this gap is open.
The Gaussian obstruction in Proposition~\ref{lower:prop:accuracy},
general OSE lower bounds~\cite{NelsonNguyen2014}, and sign-product
fixed-vector lower bounds~\cite{MSW2019} constrain different regimes
and do not yet give a matched joint theorem. For raw regression,
Proposition~\ref{app:prop:regression-lower} gives the obstruction
under fourth moments, while Theorem~\ref{app:thm:subgauss-confidence}
gives polylogarithmic confidence for sub-Gaussian product laws on every
fixed input. The jointly sharp dependence on confidence, accuracy,
rank, and tensor order remains open.

\paragraph{Subspace geometry.}
Theorem~\ref{geom:thm:schmidt} identifies the largest Schmidt weight
as the parameter governing the order-two common-Schmidt transition.
The all-confidence bound retains this geometric parameter, while
Corollary~\ref{app:cor:input} isolates the mixed moment relevant to
approximation. A broader characterization of subspaces admitting $O(r)$ embedding
width, with constants independent of large mode dimensions, remains
open. Useful geometric criteria should be computable from the input
and explain when the extreme-row obstruction disappears.

\paragraph{Weighted regression.}
The conditional bound \eqref{weighted:eq:conditional-tail} suggests
controlling $\tr C_m$, $\tr(C_m^2)$, and $\norm{C_m}$ jointly with
the weighted Gram inverse. A large weight affects both factors in
$C_m$; replacing it by its maximum can lose this compensation.
The fixed-rank limit identifies $\E w^2$ as the limiting variance
factor, but does not establish a uniform finite-sample bound as rank
grows. Removing unnecessary logarithms from the confidence certificate
requires such an estimate. The general Hilbert comparison in
Theorem~\ref{app:thm:subgauss-confidence} provides another sufficient
certificate; its minimum with the weighted certificate is valid without
spending additional failure probability.

\phantomsection
\addcontentsline{toc}{section}{References}
\begin{thebibliography}{99}\small
\setlength{\itemsep}{2pt}\setlength{\parskip}{0pt}

\bibitem{Ahle2020}
T. D. Ahle, M. Kapralov, J. B. T. Knudsen, R. Pagh, A. Velingker,
D. P. Woodruff, and A. Zandieh,
\emph{Oblivious sketching of high-degree polynomial kernels},
Proceedings of the 2020 ACM--SIAM Symposium on Discrete Algorithms,
2020, 141--160.
\href{https://doi.org/10.1137/1.9781611975994.9}{doi:10.1137/1.9781611975994.9}.

\bibitem{Simons2026}
N. Amsel et al.,
\emph{Linear systems and eigenvalue problems: Open questions from a Simons workshop},
2026, version 3 (21 August 2026).
\href{https://arxiv.org/abs/2602.05394v3}{arXiv:2602.05394v3}.

\bibitem{BambergerKrahmerWard2022}
S. Bamberger, F. Krahmer, and R. Ward,
\emph{Johnson--Lindenstrauss embeddings with Kronecker structure},
SIAM Journal on Matrix Analysis and Applications \textbf{43}(4) (2022),
1806--1850.
\href{https://doi.org/10.1137/21M1432491}{doi:10.1137/21M1432491}.

\bibitem{MatrixChaos2025}
A. S. Bandeira, K. Lucca, P. Nizi\'c-Nikolac, and R. van Handel,
\emph{Matrix chaos inequalities and chaos of combinatorial type},
2025, version 2.
\href{https://arxiv.org/abs/2412.18468v2}{arXiv:2412.18468v2}.

\bibitem{BM2026}
L. Beretta and C. Musco,
\emph{A tight analysis of Khatri--Rao oblivious subspace embeddings},
2026, version 3.
\href{https://arxiv.org/abs/2608.28094v3}{arXiv:2608.28094v3}.

\bibitem{Bousquet2002}
O. Bousquet,
\emph{A Bennett concentration inequality and its application to
suprema of empirical processes},
Comptes Rendus Math\'ematique \textbf{334}(6) (2002), 495--500.
\href{https://doi.org/10.1016/S1631-073X(02)02292-6}
{doi:10.1016/S1631-073X(02)02292-6}.

\bibitem{BGKL2026}
Z. Bujanovi\'c, L. Grubi\v{s}i\'c, D. Kressner, and H. Y. Lam,
\emph{Subspace embedding with random Khatri--Rao products and its application
to eigensolvers}, IMA Journal of Numerical Analysis \textbf{46}(3) (2026),
1328--1355.
\href{https://doi.org/10.1093/imanum/draf043}{doi:10.1093/imanum/draf043}.

\bibitem{BK2021}
Z. Bujanovi\'c and D. Kressner,
\emph{Norm and trace estimation with random rank-one vectors},
SIAM Journal on Matrix Analysis and Applications \textbf{42}(1) (2021),
202--223.
\href{https://doi.org/10.1137/20M1331718}{doi:10.1137/20M1331718}.

\bibitem{CHZ2022}
T. T. Cai, R. Han, and A. R. Zhang,
\emph{On the non-asymptotic concentration of heteroskedastic Wishart-type matrix},
Electronic Journal of Probability \textbf{27} (2022), 1--40.
\href{https://doi.org/10.1214/22-EJP758}{doi:10.1214/22-EJP758}.

\bibitem{CEMT2025}
C. Cama\~no, E. N. Epperly, R. A. Meyer, and J. A. Tropp,
\emph{Faster linear algebra algorithms with structured random matrices},
2025. \href{https://arxiv.org/abs/2508.21189}{arXiv:2508.21189}.

\bibitem{TTStack2026}
P. Cazeaux, M.-S. Dupuy, and R. Figueroa Justiniano,
\emph{Linear-scaling tensor train sketching}, 2026, version 3.
\href{https://arxiv.org/abs/2603.11009v3}{arXiv:2603.11009v3}.

\bibitem{ChenJin2021}
K. Chen and R. Jin,
\emph{Tensor-structured sketching for constrained least squares},
SIAM Journal on Matrix Analysis and Applications \textbf{42}(4) (2021),
1703--1731.
\href{https://doi.org/10.1137/20M1374596}{doi:10.1137/20M1374596}.

\bibitem{CW2009}
K. L. Clarkson and D. P. Woodruff,
\emph{Numerical linear algebra in the streaming model},
Proceedings of the 41st ACM Symposium on Theory of Computing (2009), 205--214.
\href{https://doi.org/10.1145/1536414.1536445}{doi:10.1145/1536414.1536445}.

\bibitem{CNW2016}
M. B. Cohen, J. Nelson, and D. P. Woodruff,
\emph{Optimal approximate matrix product in terms of stable rank},
43rd International Colloquium on Automata, Languages, and Programming
(ICALP 2016), LIPIcs \textbf{55}, 11:1--11:14.
\href{https://doi.org/10.4230/LIPIcs.ICALP.2016.11}{doi:10.4230/LIPIcs.ICALP.2016.11}.

\bibitem{DHWZ2026}
G. Dai, Y. He, K. Wang, and Y. Zhu,
\emph{Concentration of bounded sparse chaoses and sparse Khatri--Rao
embeddings}, 2026, version 1.
\href{https://arxiv.org/abs/2609.28875v1}{arXiv:2609.28875v1}.

\bibitem{DobribanLiu2019}
E. Dobriban and S. Liu,
\emph{Asymptotics for sketching in least squares regression},
Advances in Neural Information Processing Systems \textbf{32} (2019).
\href{https://arxiv.org/abs/1810.06089v2}{arXiv:1810.06089v2}.

\bibitem{DMMW2011}
P. Drineas, M. W. Mahoney, S. Muthukrishnan, and T. Sarl\'os,
\emph{Faster least squares approximation},
Numerische Mathematik \textbf{117}(2) (2011), 219--249.
\href{https://doi.org/10.1007/s00211-010-0331-6}{doi:10.1007/s00211-010-0331-6}.

\bibitem{HMT2011}
N. Halko, P.-G. Martinsson, and J. A. Tropp,
\emph{Finding structure with randomness: Probabilistic algorithms for
constructing approximate matrix decompositions},
SIAM Review \textbf{53}(2) (2011), 217--288.
\href{https://doi.org/10.1137/090771806}{doi:10.1137/090771806}.

\bibitem{Hoeffding1963}
W. Hoeffding,
\emph{Probability inequalities for sums of bounded random variables},
Journal of the American Statistical Association \textbf{58}(301) (1963),
13--30.
\href{https://doi.org/10.1080/01621459.1963.10500830}
{doi:10.1080/01621459.1963.10500830}.

\bibitem{Isserlis}
L. Isserlis,
\emph{On a formula for the product-moment coefficient of any order of a
normal frequency distribution in any number of variables},
Biometrika \textbf{12}(1--2) (1918), 134--139.
\href{https://doi.org/10.1093/biomet/12.1-2.134}{doi:10.1093/biomet/12.1-2.134}.

\bibitem{Janson}
S. Janson, \emph{Gaussian Hilbert Spaces}, Cambridge Tracts in Mathematics
\textbf{129}, Cambridge University Press, 1997.
\href{https://doi.org/10.1017/CBO9780511526169}{doi:10.1017/CBO9780511526169}.

\bibitem{JinKoldaWard2021}
R. Jin, T. G. Kolda, and R. Ward,
\emph{Faster Johnson--Lindenstrauss transforms via Kronecker products},
Information and Inference: A Journal of the IMA \textbf{10}(4) (2021),
1533--1562.
\href{https://doi.org/10.1093/imaiai/iaaa028}{doi:10.1093/imaiai/iaaa028}.

\bibitem{KM2015}
V. Koltchinskii and S. Mendelson,
\emph{Bounding the smallest singular value of a random matrix without
concentration}, International Mathematics Research Notices
\textbf{2015}(23), 12991--13008.
\href{https://doi.org/10.1093/imrn/rnv096}{doi:10.1093/imrn/rnv096}.

\bibitem{Latala1997}
R. Lata\l{}a,
\emph{Estimation of moments of sums of independent real random variables},
Annals of Probability \textbf{25}(3) (1997), 1502--1513.
\href{https://doi.org/10.1214/aop/1024404522}{doi:10.1214/aop/1024404522}.

\bibitem{MaMahoneyYu2015}
P. Ma, M. W. Mahoney, and B. Yu,
\emph{A statistical perspective on algorithmic leveraging},
Journal of Machine Learning Research \textbf{16}(27) (2015), 861--911.
\url{https://www.jmlr.org/papers/v16/ma15a.html}.

\bibitem{MSW2019}
M. Meister, T. Sarl\'os, and D. P. Woodruff,
\emph{Tight dimensionality reduction for sketching low degree polynomial kernels},
Advances in Neural Information Processing Systems \textbf{32} (2019).
\url{https://proceedings.neurips.cc/paper/2019/hash/ba6d843eb4251a4526ce65d1807a9309-Abstract.html}.

\bibitem{MuscoMusco2015}
C. Musco and C. Musco,
\emph{Randomized block Krylov methods for stronger and faster approximate singular value decomposition},
Advances in Neural Information Processing Systems \textbf{28} (2015).
\href{https://arxiv.org/abs/1504.05477}{arXiv:1504.05477}.

\bibitem{NelsonNguyen2014}
J. Nelson and H. L. Nguyen,
\emph{Lower bounds for oblivious subspace embeddings},
Automata, Languages, and Programming (ICALP 2014),
Lecture Notes in Computer Science \textbf{8572}, 883--894.
\href{https://doi.org/10.1007/978-3-662-43948-7_73}{doi:10.1007/978-3-662-43948-7\_73}.

\bibitem{Oliveira2010}
R. I. Oliveira,
\emph{Sums of random Hermitian matrices and an inequality by Rudelson},
Electronic Communications in Probability \textbf{15} (2010), 203--212.
\href{https://doi.org/10.1214/ECP.v15-1544}{doi:10.1214/ECP.v15-1544}.

\bibitem{Oliveira2016}
R. I. Oliveira,
\emph{The lower tail of random quadratic forms with applications to
ordinary least squares}, Probability Theory and Related Fields
\textbf{166}(3--4) (2016), 1175--1194.
\href{https://doi.org/10.1007/s00440-016-0738-9}{doi:10.1007/s00440-016-0738-9}.

\bibitem{RakhshanRabusseau2020}
B. Rakhshan and G. Rabusseau,
\emph{Tensorized random projections}, Proceedings of the Twenty Third
International Conference on Artificial Intelligence and Statistics,
Proceedings of Machine Learning Research \textbf{108} (2020), 3306--3316.
\href{https://proceedings.mlr.press/v108/rakhshan20a.html}{Proceedings of Machine Learning Research}.

\bibitem{Rudelson1999}
M. Rudelson, \emph{Random vectors in the isotropic position},
Journal of Functional Analysis \textbf{164}(1) (1999), 60--72.
\href{https://doi.org/10.1006/jfan.1998.3384}{doi:10.1006/jfan.1998.3384}.

\bibitem{SVB2026}
A. K. Saibaba, B. D. Verma, and G. Ballard,
\emph{Improved analysis of Khatri--Rao random projections and applications},
2026, version 2 (1 April 2026).
\href{https://arxiv.org/abs/2507.23207v2}{arXiv:2507.23207v2}.

\bibitem{Sarlos2006}
T. Sarl\'os,
\emph{Improved approximation algorithms for large matrices via random projections},
Proceedings of the 47th IEEE Symposium on Foundations of Computer Science (2006), 143--152.
\href{https://doi.org/10.1109/FOCS.2006.37}{doi:10.1109/FOCS.2006.37}.

\bibitem{Tikhomirov2018}
K. Tikhomirov,
\emph{Sample covariance matrices of heavy-tailed distributions},
International Mathematics Research Notices
\textbf{2018}(20), 6254--6289.
\href{https://doi.org/10.1093/imrn/rnx067}{doi:10.1093/imrn/rnx067}.
\href{https://arxiv.org/abs/1606.03557}{arXiv:1606.03557}.

\bibitem{TW2026}
A. Townsend and C. Wang,
\emph{Oblivious subspace injection is not enough for relative error},
2026, version 2.
\href{https://arxiv.org/abs/2604.10215v2}{arXiv:2604.10215v2}.

\bibitem{TroppElementary}
J. A. Tropp,
\emph{The expected norm of a sum of independent random matrices:
an elementary approach}, in \emph{High Dimensional Probability VII:
The Carg\`ese Volume}, Progress in Probability \textbf{71}, Birkh\"auser,
2016, 173--202.
\href{https://doi.org/10.1007/978-3-319-40519-3_8}
{doi:10.1007/978-3-319-40519-3\_8}.

\bibitem{TYUCPSD2017}
J. A. Tropp, A. Yurtsever, M. Udell, and V. Cevher,
\emph{Fixed-rank approximation of a positive-semidefinite matrix from
streaming data}, Advances in Neural Information Processing Systems
\textbf{30} (2017).
\href{https://arxiv.org/abs/1706.05736}{arXiv:1706.05736}.

\bibitem{YBK1988}
Y. Q. Yin, Z. D. Bai, and P. R. Krishnaiah,
\emph{On the limit of the largest eigenvalue of the large dimensional sample covariance matrix},
Probability Theory and Related Fields \textbf{78}(4) (1988), 509--521.
\href{https://doi.org/10.1007/BF00353874}{doi:10.1007/BF00353874}.

\bibitem{Rungta2001}
P. Rungta, V. Bu\v{z}ek, C. M. Caves, M. Hillery, and G. J. Milburn,
\emph{Universal state inversion and concurrence in arbitrary dimensions},
Physical Review A \textbf{64} (2001), 042315.
\href{https://doi.org/10.1103/PhysRevA.64.042315}
{doi:10.1103/PhysRevA.64.042315}.
\href{https://arxiv.org/abs/quant-ph/0102040}{arXiv:quant-ph/0102040}.

\bibitem{MNPP2015}
D. Malicet, I. Nourdin, G. Peccati, and G. Poly,
\emph{Squared chaotic random variables: New moment inequalities with
applications}, 2015.
\href{https://arxiv.org/abs/1503.02154}{arXiv:1503.02154}.

\bibitem{Catoni2012}
O. Catoni,
\emph{Challenging the empirical mean and empirical variance: A deviation study},
Annales de l'Institut Henri Poincar\'e, Probabilit\'es et Statistiques
\textbf{48} (2012), no.~4, 1148--1185.
\href{https://doi.org/10.1214/11-AIHP454}{doi:10.1214/11-AIHP454}.

\bibitem{Eaton1970}
M. L. Eaton,
\emph{A note on symmetric Bernoulli random variables},
Annals of Mathematical Statistics \textbf{41}(4) (1970), 1223--1226.
\href{https://doi.org/10.1214/aoms/1177696897}{doi:10.1214/aoms/1177696897}.

\bibitem{Komorowski1988}
R. Komorowski,
\emph{On the best possible constants in the Khintchine inequality for $p\ge3$},
Bulletin of the London Mathematical Society \textbf{20}(1) (1988), 73--75.

\bibitem{Holden2026}
S. Holden,
\emph{RA-11: probability-sensitive product probes and diagonal correction},
public working note, 14 September 2026, Sections 3, 4, 6, and 7.
Repository contribution with disclosed AI assistance.
\href{https://github.com/ajt60gaibb/OpenProblemsInNLA/blob/main/references/holden-further-2026-09-14/RA-11/manuscript.md}{Full proof text}.

\bibitem{CETW2025}
Y. Chen, E. N. Epperly, J. A. Tropp, and R. J. Webber,
\emph{Randomly pivoted Cholesky: Practical approximation of a kernel
matrix with few entry evaluations},
Communications on Pure and Applied Mathematics (2025).
\href{https://arxiv.org/abs/2207.06503}{arXiv:2207.06503}.

\bibitem{Pinelis1994}
I. Pinelis,
\emph{Optimum bounds for the distributions of martingales in Banach spaces},
Annals of Probability \textbf{22}(4) (1994), 1679--1706.
\href{https://arxiv.org/abs/1208.2200}{arXiv:1208.2200}.

\end{thebibliography}

\clearpage\appendix
\section{Embedding proofs and quantitative certificates}
\label{appendix:covariance}
\subsection{Accuracy refinement by sample covariance}
\label{appendix:variance-refinement}
\begin{proof}[Proof of Theorem~\ref{variance:thm:refinement}]
Write $p=d-1$ and $\Phi(s)=rs^p+s^{p+1}$.
Constants denoted $C_{d,\Ksg}$ depend only on $d,\Ksg$ and may change
between displays. We first assume $m\ge2r$, as both final budgets ensure
after increasing their constants.

\paragraph{Row norms.}
Let $G\sim N(0,I_r)$ be independent of $y$. Integrating
$|\ip Gy|^{2q}$ in the two orders and using
\eqref{core:eq:directional}--\eqref{core:eq:gaussian-moments} gives
$\gmom_q\E\norm y^{2q}\le(q!)^d\Ksg^{2dq}D_{r,q}$ for integers $q\ge2$,
where we dropped $\theta^d\le1$. Since $\gmom_q\ge q!$, $q!\le q^q$,
and $D_{r,q}\le(r+2q)^q$,
\begin{equation}\label{variance:eq:row-moment}
 (\E\norm y^{2q})^{1/q}\le\Ksg^{2d}q^p(r+2q).
\end{equation}
Markov's inequality with $q=\lceil\log(m/\beta)\rceil$, increased to two
if necessary, and a union bound give, with failure at most $\beta$,
\begin{equation}\label{variance:eq:row-maximum}
 \max_{i\le m}\norm{y_i}^2
 \le C_{d,\Ksg}\Phi(\log(2m/\beta)).
\end{equation}
Taking $q=\lceil\log(em)\rceil$ also gives
\begin{equation}\label{variance:eq:expected-maximum}
 \E\max_{i\le m}\norm{y_i}^2\le C_{d,\Ksg}\Phi(\log(em)),
\end{equation}
because the $L_q$ norm of the maximum is at most $m^{1/q}$ times
the $L_q$ norm of one squared row norm.

\paragraph{Published covariance input.}
The $y_i$ need not be centered. Introduce independent Rademacher signs
$\eta_i$ and apply \cite[Corollary~2]{Tikhomirov2018} to $\eta_i y_i$
at moment order six. These vectors are centered and isotropic;
\eqref{core:eq:directional} bounds their directional sixth moments by a
constant depending only on $d,\Ksg$. Their Gram matrices and norms are
unchanged. For $r\ge2$, outside an event $\mathcal B$ of probability at most $1/r$,
\begin{equation}\label{variance:eq:tikhomirov}
 \norm Z\le C_{d,\Ksg}\left[\sqrt{\frac rm}
                   +\frac{\max_i\norm{y_i}^2}{m}\right].
\end{equation}
Use \eqref{variance:eq:row-maximum} with $\beta=\delta/2$. If $\delta\ge2/r$,
then, with failure at most $\delta$,
\begin{equation}\label{variance:eq:moderate-implicit}
 \norm Z\le C_{d,\Ksg}\left[\sqrt{\frac rm}
                   +\frac{\Phi(\log(4m/\delta))}{m}\right].
\end{equation}
Independence and \eqref{core:eq:variance} also imply
$\E\norm Z_{\F}^2=m^{-1}\tr\E(H-I_r)^2\le C_{d,\Ksg}r^2/m$.
Cauchy--Schwarz bounds the contribution of $\mathcal B$ to $\E\norm Z$
by $C_{d,\Ksg}\sqrt{r/m}$. Integrate
\eqref{variance:eq:tikhomirov} on its complement and use
\eqref{variance:eq:expected-maximum} to obtain
\begin{equation}\label{variance:eq:expectation}
 \E\norm Z\le C_{d,\Ksg}\left[\sqrt{\frac rm}
                   +\frac{\Phi(\log(em))}{m}\right].
\end{equation}
For $r=1$, the scalar variance bound gives this directly.

\paragraph{Arbitrary confidence.}
Set $q=\lceil\log(8m/\delta)\rceil$ and
$R=e\Ksg^{2d}q^p(r+2q)$. Equation~\eqref{variance:eq:row-moment} gives
\[
 \Pp\{\norm y^2>R\}\le e^{-q}\le\frac\delta{8m},\qquad
 b:=\E[\norm y^2\mathbf1_{\{\norm y^2>R\}}]\le Re^{-q}
                                                    \le\frac{\delta R}{8m}.
\]
Define $\widetilde H=yy^\top\mathbf1_{\{\norm y^2\le R\}}$ and
$\widetilde Z=m^{-1}\sum_i\widetilde H_i-\E\widetilde H$.
With failure at most $\delta/8$, no sample is removed, and
$\norm Z\le\norm{\widetilde Z}+b$. Also,
\[
 \E\norm{\widetilde Z}\le\E\norm Z+2b
             \le C_{d,\Ksg}\left[\sqrt{r/m}+R/m\right].
\]
Truncation occurs only in this proof; the sketch remains unchanged.
Apply \cite[Theorem~2.3]{Bousquet2002} to the centered quadratic forms
$a^\top\widetilde Ha-\E(a^\top\widetilde Ha)$ and their negatives,
on a countable dense subset of the unit sphere. Their absolute values
are at most $R$, and their variances are bounded by
$\E|\ip ay|^4\le C_{d,\Ksg}$. With $u=\log(8/\delta)$, the inequality
and $2\sqrt{xy}\le x+y$ give, with failure at most $\delta/8$,
\[
 \norm{\widetilde Z}\le2\E\norm{\widetilde Z}
             +C_{d,\Ksg}\sqrt{u/m}+\frac{4Ru}{3m}.
\]
Combining the bounds yields, with failure at most $\delta/4$,
\begin{equation}\label{variance:eq:general-implicit}
 \norm Z\le C_{d,\Ksg}\left[\sqrt{\frac{r+u}{m}}
                   +\frac{u\Phi(\log(8m/\delta))}{m}\right].
\end{equation}

\paragraph{Closing the sample-size logarithm.}
Use $(a,w)=(r,1)$ for \eqref{variance:eq:moderate-implicit} and
$(a,w)=(r+u,u)$ for \eqref{variance:eq:general-implicit}. Put
\[
 t=\log(4r/\delta),\qquad F=\frac a{\varepsilon^2}
       +\frac{w\Phi(t)}{\varepsilon},\qquad M=AF,
       \qquad\ell=\log(1/\varepsilon),
\]
where $A$ is a sufficiently large constant depending only on $d,\Ksg$.
In both cases $a\ge r$, $1\le w\le2t$, and
$\log(8M/\delta)\le C_d[t+\ell+\log(eA)]$.
For each fixed integer $k\ge1$, splitting into $\ell\le t$ and $\ell>t$ gives
\[
 w\ell^k\le wt^k+2\ell^{k+1}\le wt^k+C_k/\varepsilon.
\]
It follows that
\[
 w\Phi(\log(8M/\delta))\le C_d[1+\log^d(eA)]
       \left[w\Phi(t)+r/\varepsilon\right]
       \le C_d[1+\log^d(eA)]\varepsilon F.
\]
The square-root term at $M$ is at most $\varepsilon/\sqrt A$ and the
linear term at most $C_{d,\Ksg}[1+\log^d(eA)]\varepsilon/A$.
Choose $A$ so their sum is at most $\varepsilon$.
Finally, $\log^k(8m/\delta)/m$ decreases once $\log(8m/\delta)\ge k$;
increasing the same constant ensures this for $k\le d$ and every $m\ge M$.
Both conclusions therefore hold for every integer sample size above the stated budgets.
\end{proof}


\subsection{Radialization and the factorial envelope}\label{appendix:factorial}
\begin{lemma}[Direction-preserving radialization]\label{core:lem:radial}
Let $y\in\R^r$ satisfy
$\E|\ip u y|^{2p}\le A_p\norm u^{2p}$ for all $u\in\R^r$, where
$p\ge2$ is an integer. Then
\begin{equation}\label{core:eq:radial}
 \E(yy^{\mathsf T})^p\preceq
 \frac{c_p^{\rm rad}A_pD_{r,p}}{r\gmom_p}I_r,\qquad
 c_p^{\rm rad}=3\sqrt3\left(\frac{p-1}{p}\right)^{p-1}.
\end{equation}
\end{lemma}
\begin{proof}
Fix a unit vector $x$ and, independently of $y$, take
$z\sim N(0,I_r+\lambda xx^{\mathsf T})$, with $\lambda>0$.
Conditional Gaussian integration gives
\[
 \E_z\ip z y^{2p}
 =\gmom_p(\norm y^2+\lambda\ip x y^2)^p
 \ge\gmom_p\frac{p^p}{(p-1)^{p-1}}\lambda
                  \norm y^{2p-2}\ip x y^2.
\]
The last step is weighted AM--GM:
$(a+b)^p\ge p^p a^{p-1}b/(p-1)^{p-1}$ for $a,b\ge0$.
Integrating in the opposite order therefore yields
\begin{equation}\label{core:eq:tilt}
 x^{\mathsf T}\E(yy^{\mathsf T})^px
 \le\frac{(p-1)^{p-1}A_p}{\gmom_pp^p\lambda}\E\norm z^{2p}.
\end{equation}
For $G\sim N(0,I_r)$, $\norm z^2$ has the law
$\norm G^2+\lambda\ip xG^2$. Gaussian polar decomposition gives
\[
 \E\norm z^{2p}
 =D_{r,p}\sum_{k=0}^p\binom pk\lambda^k
                  \frac{\gmom_k}{D_{r,k}}.
\]
Here $\E\norm G^{2p}=D_{r,p}$ follows by integrating the chi-square
density, and the spherical moment $\gmom_k/D_{r,k}$ follows by
factoring $\E\ip xG^{2k}=\gmom_k$ into independent radius and direction.
Using $D_{r,k}\ge r^k$ and $\binom pk\le p^k/k!$ gives
\[
 \E\norm z^{2p}\le
 D_{r,p}\sum_{k\ge0}\frac{\gmom_k}{k!}
                \left(\frac{p\lambda}{r}\right)^k
 =\frac{D_{r,p}}{\sqrt{1-2p\lambda/r}}
 \quad\left(\frac{p\lambda}{r}<\frac12\right).
\]
The series identity is the binomial expansion of $(1-2a)^{-1/2}$.
Set $\lambda=r/(3p)$ in \eqref{core:eq:tilt} to prove the result.
\end{proof}

\begin{lemma}[Sharp fourth moment from the exponential normalization]
\label{core:lem:sharp-sg-fourth}
Let $\xi_{\rm sg}>0$ be the unique solution of
\begin{equation}\label{core:eq:sg-fourth-root}
 \frac{e^{\xi_{\rm sg}}-1}{\xi_{\rm sg}}=\Ksg^2,
 \qquad c_{\rm sg}=\Ksg^2\xi_{\rm sg}.
\end{equation}
Every factor in Lemma~\ref{core:lem:factor} satisfies
\[
 \E|\ip ug|^4\le c_{\rm sg}\norm u^4.
\]
Consequently its tensor product has directional fourth-moment constant
$c_{\rm sg}^d$. The scalar constant is sharp under isotropy and
\eqref{core:eq:model}, and
$1\le c_{\rm sg}<2\theta\Ksg^4$.
\end{lemma}
\begin{proof}
For unit $u$, put $T=|\ip ug|^2/\Ksg^2$ and
$\mu=\E T=\Ksg^{-2}$. Define the probability measure
$d\mathbb Q=T\,d\Pp/\mu$. The function
\[
 f(t)=\frac{e^t-1}{t}=\int_0^1e^{st}\,ds,
 \qquad f(0)=1,
\]
is strictly increasing and convex on $[0,\infty)$. Jensen under
$\mathbb Q$ gives
\[
 f\!\left(\frac{\E T^2}{\mu}\right)
 \le\E_{\mathbb Q}f(T)
 =\frac{\E e^T-1}{\mu}\le\Ksg^2.
\]
Thus $\E T^2\le\mu\xi_{\rm sg}$, proving the fourth-moment bound.
Since $\Ksg^2\ge1/\log2>1$, the defining positive root exists uniquely;
also $f(\log2)=1/\log2$ gives $\xi_{\rm sg}\ge\log2$ and
$c_{\rm sg}\ge1$.

For sharpness, let a scalar $X$ equal $\sqrt{c_{\rm sg}}$ and
$-\sqrt{c_{\rm sg}}$ with probability $1/(2c_{\rm sg})$ each,
and zero otherwise. Then $\E X=0$, $\E X^2=1$,
$\E e^{X^2/\Ksg^2}=2$, and $\E X^4=c_{\rm sg}$.
The strict inequality
$e^b-1-b>b^2/2$ at $b=\xi_{\rm sg}>0$ gives
$\xi_{\rm sg}<2(\Ksg^2-1)$ and hence
$c_{\rm sg}<2\theta\Ksg^4$.
Convexity over normalized eigenvalues also gives
$\E(g^\top Bg)^2\le c_{\rm sg}(\tr B)^2$ for $B\succeq0$.
The factorwise induction in Lemma~\ref{core:lem:factor} therefore
gives the product constant.
\end{proof}

At $\Ksg^2=8/3$, the scalar fourth-moment constant is
$c_{\rm sg}=4.588619425\ldots$, compared with $80/9$ from the
exponential remainder alone. This improves the finite variance budgets
and fourth-moment approximation results. The universal factorial
normalization below continues to use $[2\theta\Ksg^4]^d$.
This distinction is necessary for the stated all-order envelope.
For the three-point extremizer with $\xi_{\rm sg}=3$, take $r=1$
and $d=2$. Then $\E H^3=c_{\rm sg}^4$, whereas replacing the
factorial normalization by $c_{\rm sg}^2$ while retaining radius coefficient $1/2$
would require $c_{\rm sg}^4\le6c_{\rm sg}^2\Ksg^4$.
The ratio of the two sides is $\xi_{\rm sg}^2/6=3/2$.

Apply Lemma~\ref{core:lem:radial}
to \eqref{core:eq:directional} with $a=Uu$. This retains a separate
explicit matrix budget at every order:
\begin{equation}\label{core:eq:direct-budget}
 \E H^p\preceq a_p^{\mathrm{dir}}I_r,\qquad
 a_p^{\mathrm{dir}}=
 \frac{c_p^{\rm rad}[\theta\beta_p\Ksg^{2p}]^dD_{r,p}}{r\gmom_p}
 \quad(p\ge2).
\end{equation}

At second order the finite tilted Gaussian calculation can be minimized
explicitly. Fix a unit test vector $x$ and let
$z\sim N(0,I_r+\lambda xx^\top)$. The inequality
$(a+b)^2\ge4ab$ gives
\[
 x^\top\E H^2x
 \le\frac{A_2}{12\lambda}
   \left[r(r+2)+2\lambda(r+2)+3\lambda^2\right],
 \qquad A_2=[2\theta\Ksg^4]^d.
\]
The choice $\lambda=\sqrt{r(r+2)/3}$ yields
\[
 \E H^2\preceq vI_r,\qquad
 V:=\E(H-I_r)^2\preceq(v-1)I_r.
\]
Moreover $s_r\ge1$ and $2\theta\Ksg^4>1$, so $v>1$.
The latter strict inequality follows from isotropy and
$e^x-1-x>x^2/2$ for $x>0$, by taking expectations in any unit factor direction.
The covariance comparison is easiest to apply when the higher moments
share a factorial envelope. We first record two scalar estimates for
the full exponential remainder, then derive the envelope while keeping
the same second-moment normalization.

For $p\ge3$, the supremum defining $\beta_p$ is attained at the unique
positive root $x_p$ of
\[
 (p-x)e^x=p+(p-1)x,
 \qquad p-1<x_p<p.
\]
Indeed, the derivative has the sign of
$D_p(x)=(p-x)e^x-p-(p-1)x$, whose second derivative is
$(p-2-x)e^x$. Starting from $D_p(0)=D_p'(0)=0$, this shows that
$D_p$ increases before decreasing through a unique positive zero.
Moreover, $D_p(p)<0$ and $D_p(p-1)>0$, since
$e^n>1+n+n^2$ for $n\ge2$. The last inequality follows from
\[
 \frac{n^2}{2}+\frac{n^3}{6}+\frac{n^4}{24}-n^2
 =\frac{n^2(n-2)(n+6)}{24}\ge0
\]
and the remaining positive exponential terms.
Consequently,
\begin{equation}\label{core:eq:beta-estimates}
 \beta_3\le3,\qquad
 \beta_p\le\frac{(p/e)^p}{1-pe^{1-p}}
           <\frac54(p/e)^p\quad(p\ge4).
\end{equation}
For the first inequality, Taylor expansion and AM--GM give
$e^x-1-x\ge x^3[1/6+1/(2\sqrt3)]\ge x^3/3$.
For the second, use $x_p>p-1$, maximize $x^pe^{-x}$ at $p$, and
note $pe^{1-p}\le4e^{-3}<1/5$.
For reference, $\beta_3\approx1.82864$ and $\beta_4\approx5.24911$,
compared with $3!=6$ and $4!=24$.

\begin{proposition}[Factorial matrix envelope]\label{core:prop:raw}
For every integer $p\ge2$,
\begin{equation}\label{core:eq:raw}
 \E H^p\preceq\frac{p!}{2}v
 \bigl[\tfrac12\Ksg^{2d}p^{d-2}(r+p)\bigr]^{p-2}I_r.
\end{equation}
\end{proposition}
\begin{proof}[Proof of Proposition~\ref{core:prop:raw}]
The case $p=2$ is \eqref{core:eq:variance}. For $p\ge3$, the full
exponential remainder and weighted radialization give
\[
 \E H^p\preceq
 c_p^{\rm rad}[\theta\beta_p\Ksg^{2p}]^d\frac{D_{r,p}}{r\gmom_p}I_r,
 \qquad
 c_p^{\rm rad}=3\sqrt3\left(\frac{p-1}{p}\right)^{p-1}.
\]
The elementary inequality $p!\le2p^{p-2}$ follows inductively, since
the ratio of successive normalized quantities is
$(p/(p+1))^{p-2}\le1$. Since $\beta_p\le p!$, the ratio of this raw
bound to the asserted factorial envelope at order $d$ is the
corresponding ratio at $d=2$ multiplied by
\[
 \left(\frac{\beta_p}{2p^{p-2}}\right)^{d-2}\le1.
\]
It is therefore enough to prove the claim for $d=2$.

Write
\[
 F_p(r)=\frac{D_{r,p}}{r^2(r+p)^{p-2}}.
\]
Retain the outer factors $r+2$ and $r+2p-2$ and apply AM--GM to
the remaining factors. Their arithmetic mean is $r+p$, so
\[
 F_p(r)\le\frac{(r+2)(r+2p-2)}{r(r+p)}
 =1+\frac p{r+p}+\frac{4(p-1)}{r(r+p)}
 \le\frac{3(2p-1)}{p+1}<6.
\]
The ratio at $d=2$, with radius coefficient $1/2$, is
\[
 R_p=\frac{c_p^{\rm rad}\beta_p^2}{2p!\gmom_p}
          F_p(r)\frac r{s_r}\,2^{p-2}.
\]
For $p=3$, one has
\[
 F_3(r)\frac r{s_r}
 =\frac{6(r+2)(r+4)}{(r+3)(r+2+\sqrt{3r(r+2)})}\le4.
\]
After rearrangement and squaring, this inequality follows from
\[
 12r(r+3)^2-(r+6)^2(r+2)
 =11r^3+58r^2+48r-72\ge0\qquad(r\ge1).
\]
The scalar estimate $\beta_3\le3$ and $c_3^{\rm rad}=4\sqrt3/3$ now give
$R_3\le8\sqrt3/15<1$.

For $p\ge4$, the scalar remainder estimate gives
$\beta_p\le(5/4)(p/e)^p$. Also $c_p^{\rm rad}\le c_3^{\rm rad}$; indeed, the derivative
of $(x-1)\log(1-1/x)$ is $\log(1-1/x)+1/x<0$.
Using $p!\gmom_p=(2p)!/2^p$, Stirling's lower bound, and
$r/s_r\le6/(1+\sqrt3)$ gives
\[
 R_p\le
 \frac{25c_p^{\rm rad}F_p(r)}{128\sqrt{4\pi p}}\frac r{s_r}
 \le\frac{75\sqrt3}{32\sqrt\pi(1+\sqrt3)}<1.
\]
This proves the envelope for all $p\ge2$ and $d\ge2$.
\end{proof}
\subsection{The centered step and closed moment comparison}\label{appendix:closed}
Rank one makes each cyclic word a product of scalar matrix coefficients.
Spectral H\"older controls higher insertions; the two-insertion estimate
is the classical trace AM--GM inequality~\cite{TroppElementary}.
\begin{lemma}[Centered one-step inequality]\label{closed:lem:step}
Let $H\succeq0$ have rank at most one, $\E H=I_r$, and finite moments
through order $2p$. Put $X=H-I_r$. For every deterministic symmetric $A$,
\begin{align}\label{closed:eq:one-step}
 \E\tr(A+X)^{2p}\le{}&\tr A^{2p}
 +\binom{2p}{2}\tr(A^{2p-2}\E X^2)\nonumber\\
 &+\sum_{j=3}^{2p}\binom{2p}{j}
 \tr\bigl(|A|^{2p-j}\E(H+I_r)^j\bigr).
\end{align}
\end{lemma}
\begin{proof}[Proof of Lemma~\ref{closed:lem:step}]
Indeed, the one-insertion terms vanish since $\E X=0$. For two
insertions, if $a+b={n_A}$ is even, diagonalizing $A$ and applying weighted
AM--GM to each matrix entry gives
\[
 \tr(A^aXA^bX)\le\tr(|A|^{n_A}X^2).
\]
For higher insertions, expand each $X$ as $H-I_r$. Write
$H=\tau uu^{\mathsf T}$ with $\tau\ge0$ and $\|u\|=1$.
A cyclic trace containing $\ell\ge1$ copies of $H$ and ${n_A}>0$
copies of $A$ is
\[
 \tau^\ell\prod_{a=1}^{\ell}\langle u,A^{k_a}u\rangle,
 \qquad\sum_a k_a={n_A}.
\]
Spectral H\"older bounds its absolute value by
$\tau^\ell\langle u,|A|^{n_A}u\rangle
=\tr(|A|^{n_A}H^\ell)$. The same bound holds when ${n_A}=0$;
the term with no $H$ is bounded by $\tr|A|^{n_A}$.
Summing over choices of the $\ell$ copies of $H$ supplies
$\tr(|A|^{n_A}(H+I_r)^j)$ for a word with $j$ insertions.
This proves \eqref{closed:eq:one-step}.
\end{proof}

\begin{theorem}[Closed finite-moment covariance bound]\label{closed:thm:moment}
Let $H_1,\ldots,H_m$ be iid positive semidefinite matrices of rank
at most one, with $\E H_i=I_r$. Write $X_i=H_i-I_r$,
$Z=m^{-1}\sum_iX_i$, and let $q\ge1$ be an integer. Suppose
\[
 \E X_i^2\preceq\nu I_r,\qquad
 \E(H_i+I_r)^j\preceq\frac{j!}{2}\nu\Rmom^{j-2}I_r
 \quad(3\le j\le2q),
\]
where $\nu,\Rmom>0$. With the Gaussian moments from
\eqref{core:eq:gaussian-moments}, define
\[
 h_1=0,\qquad h_q=\frac{2(q-1)}{\sqrt{2q-3}}\quad(q\ge2).
\]
Then
\begin{equation}\label{closed:eq:bound}
 (\E\tr Z^{2q})^{1/(2q)}
 \le r^{1/(2q)}\gmom_q^{1/(2q)}
 \left(\sqrt{\frac\nu m}+\frac{h_q\Rmom}{m}\right).
\end{equation}
In particular, the $2q$-th power of the right side divided by
$\varepsilon^{2q}$ bounds $\Pp\{\|Z\|>\varepsilon\}$.
\end{theorem}

\begin{proof}[Proof of Theorem~\ref{closed:thm:moment}]
Apply Lemma~\ref{closed:lem:step}.
Let $S_n=\sum_{i=1}^nX_i$, and put
$M_{n,p}=r^{-1}\E\tr S_n^{2p}$, with $M_{n,0}=1$.
Cauchy--Schwarz in the joint probability and normalized spectral
measure yields
\[
 r^{-1}\E\tr|S_n|^{2k+1}\le\sqrt{M_{n,k}M_{n,k+1}}.
\]
For $u=(u_0,\ldots,u_q)\ge0$ define
\begin{equation}\label{core:eq:interpolation}
 \mathcal I_{2k}(u)=u_k,\qquad
 \mathcal I_{2k+1}(u)=\sqrt{u_ku_{k+1}},\qquad
 \Phi_p(u)=\sum_{j=2}^{2p}\binom{2p}{j}b_j\mathcal I_{2p-j}(u),
\end{equation}
where $b_2=\nu$ and $b_j=(j!/2)\nu\Rmom^{j-2}$ for $j\ge3$;
put $\Phi_0=0$. Conditioning \eqref{closed:eq:one-step} on
$S_{n-1}$ gives
\begin{equation}\label{closed:eq:update}
 M_{n,p}\le M_{n-1,p}+\Phi_p(M_{n-1}).
\end{equation}
Every $\Phi_p$ is componentwise nondecreasing.

For the pairing comparison, set
\begin{equation}\label{core:eq:pairing}
 \omega_{2k}=\gmom_k,\quad
 \omega_{2k+1}=\sqrt{\gmom_k\gmom_{k+1}},\quad
 C_{p,j}=\frac{(2p)!\,\omega_{2p-j}}
                   {2p(2p-j)!\gmom_p}.
\end{equation}
For $q\ge2$, these coefficients satisfy
\begin{equation}\label{closed:eq:pairing-bound}
 C_{p,j}\le h_q^{j-2}\qquad(1\le p\le q,\ 2\le j\le2p).
\end{equation}
Here $C_{p,2}=1$. For $j=2\ell$,
\[
 C_{p,2\ell}=\frac{2^{\ell-1}(p-1)!}{(p-\ell)!},
\]
so its successive ratios are at most $2(p-1)$. For $j=2\ell+1$,
writing $k=p-\ell\ge1$ gives
\[
 C_{p,2\ell+1}=\frac{2^\ell(p-1)!}{(k-1)!\sqrt{2k-1}}.
\]
Thus $C_{p,3}=h_p\le h_q$ for $p\ge2$, and the successive odd
ratios, when $k\ge2$, are
$2(k-1)\sqrt{(2k-1)/(2k-3)}\le2k\le2(p-1)$.
The inequality $h_p\le h_q$ follows from the increasing sequence
$h_p^2=2p-1+(2p-3)^{-1}$ for $p\ge2$.
Finally, $2(p-1)\le2(q-1)\le h_q^2$. Induction on each parity
proves \eqref{closed:eq:pairing-bound}.

For $q\ge2$, write $a=h_q\Rmom$, $s(t)=\sqrt{\nu t}+a$, and
\[
 B_0(t)=1,\qquad B_p(t)=\gmom_ps(t)^{2p}\quad(1\le p\le q).
\]
For $t>0$, $\mathcal I_k(B(t))=\omega_ks(t)^k$, and the pairing
bound gives
\begin{align*}
 \frac{\Phi_p(B(t))}{2p\gmom_ps(t)^{2p-1}}
 &=\sum_{j=2}^{2p}\frac{C_{p,j}b_j}{j!}s(t)^{1-j}\\
 &\le\frac\nu{2s(t)}\sum_{j=2}^{2p}
                       \left(\frac a{s(t)}\right)^{j-2}
 \le\frac\nu{2(s(t)-a)}=s'(t).
\end{align*}
Hence $B_p'(t)\ge\Phi_p(B(t))$. Since $B$ is componentwise
nondecreasing, integration on $[n,n+1]$ yields
\[
 B_p(n+1)\ge B_p(n)+\Phi_p(B(n))\quad(n\ge0).
\]
At $n=0$ this integration remains valid: $B_p$ is absolutely
continuous on every compact interval and its derivative has only an
integrable $t^{-1/2}$ singularity. The initial data satisfy
$M_{0,p}\le B_p(0)$. Induction in $n$ using
\eqref{closed:eq:update} and monotonicity of $\Phi$ now proves
$M_{m,q}\le B_q(m)$, which is \eqref{closed:eq:bound}.
For $q=1$, independence and centering give directly
$\E\tr S_m^2=m\tr\E X_1^2\le mr\nu$, yielding the asserted
bound with $h_1=0$. Finally,
$\|Z\|^{2q}\le\tr Z^{2q}$ and Markov's inequality give the
tail statement.
\end{proof}
\subsection{Shifting raw factorial budgets}\label{appendix:shift}

\begin{lemma}[Shifted factorial budgets]\label{core:lem:shift}
Suppose $H\succeq0$ has rank at most one, $\E H=I_r$, and
$\E H^j\preceq(j!/2)v\mathcal R^{j-2}I_r$ for $2\le j\le2q$, with $v,\mathcal R>0$.
Then the hypotheses of Theorem~\ref{closed:thm:moment} hold with
\begin{equation}\label{core:eq:shift-parameters}
 \nu=v,\qquad\Rmom=\mathcal R+1+\frac4{3v}.
\end{equation}
\end{lemma}

\begin{proof}[Proof of Lemma~\ref{core:lem:shift}]
The variance is $V=\E H^2-I_r\preceq vI_r$. For $j\ge3$, binomial
expansion gives
\[
 \E(H+I_r)^j\preceq
 \left[j+1+\frac{j!}{2}v
 \sum_{h=0}^{j-2}\frac{\mathcal R^{j-2-h}}{h!}\right]I_r.
\]
Set $n=j-2$ and $c=1+4/(3v)$. Every coefficient of $(\mathcal R+c)^n$
dominates the corresponding $1/h!$, since $c\ge1$.
Its excess constant coefficient is at least
$c^n-1/n!\ge c-1=4/(3v)\ge2(j+1)/(j!v)$.
This also absorbs the first term and proves the asserted budgets.
\end{proof}
\begin{proof}[Proof of Proposition~\ref{core:prop:certificate}]
\phantomsection\label{proof:certificate}
The scale $\tfrac12\Ksg^{2d}p^{d-2}(r+p)$ increases with $p$. Apply
Proposition~\ref{core:prop:raw} through order $2q$, followed by
Lemma~\ref{core:lem:shift}. Theorem~\ref{closed:thm:moment} then applies
with $\nu=v$ and $\Rmom=\Rmom_q$. Markov's inequality proves
\eqref{core:eq:explicit-closed}; solving the resulting quadratic in
$m^{-1/2}$ gives \eqref{core:eq:explicit-m}.
\end{proof}

\subsection{Sharper certificates}\label{sec:sharper}
The closed bound is sufficient for the sharp dimension order. For a fixed
moment order or a structured covariance, one can retain the individual
budgets, zero initial moments, and exact terminal insertions instead.
The resulting computation is finite and monotone in every input budget.

The following comparison is independent of the tensor model. Let
$H_1,\ldots,H_m$ be independent identically distributed positive
semidefinite matrices of rank at most one, with $\E H_i=I_r$ and finite
moments through order $2q$. Write $X_i=H_i-I_r$,
$S_n=\sum_{i=1}^nX_i$, and $Z=S_m/m$. Assume finite scalar budgets
\begin{equation}\label{core:eq:budgets}
 V=\E X_1^2\preceq\nu I_r,\qquad b_2=\nu>0,\qquad
 \E(H_1+I_r)^j\preceq b_jI_r\quad(3\le j\le2q).
\end{equation}
The budgets need not be moments of a scalar probability law.
We compare two ways to propagate them: a direct update of even trace
moments after each sample, and a single continuous scale that yields
an integral certificate. Both begin at a zero partial sum. The discrete
bound is tighter; the integral form avoids updating all moments after
each of the $m$ samples.

Use the interpolation maps and insertion operators in
\eqref{core:eq:interpolation}, now with the individual budgets $b_j$,
and the pairing weights in \eqref{core:eq:pairing}.
Starting at $u^{(0)}=(1,0,\ldots,0)$, the discrete certificate is
\begin{equation}\label{core:eq:discrete}
 u^{(n)}=u^{(n-1)}+\Phi(u^{(n-1)}).
\end{equation}
At order $2q$ these coefficients define the growth rate $\phi_q$ of
the common scale. Integrating its reciprocal gives the time needed
to reach a scale, and inversion evaluates the scale after $m$ samples.
For $s>0$, define
\begin{equation}\label{core:eq:integral}
 \phi_q(s)=\sum_{j=2}^{2q}\frac{C_{q,j}b_j}{j!}s^{1-j},\qquad
 \Gamma_q(a)=\int_0^a\frac{ds}{\phi_q(s)},\qquad
 T_q(m)=\Gamma_q^{-1}(m).
\end{equation}

\begin{theorem}[Finite-profile covariance certificate]\label{core:thm:finite}
Under \eqref{core:eq:budgets}, the inverse in
\eqref{core:eq:integral} is uniquely defined for every $m>0$, and
\begin{equation}\label{core:eq:finite-moment}
 \E\tr Z^{2q}\le\frac r{m^{2q}}u_q^{(m)}
 \le\frac{r\gmom_q}{m^{2q}}T_q(m)^{2q}.
\end{equation}
Consequently,
\begin{equation}\label{core:eq:finite-tail}
 \Pp\{\norm Z>\varepsilon\}\le
 \frac{r u_q^{(m)}}{m^{2q}\varepsilon^{2q}}
 \le\frac{r\gmom_qT_q(m)^{2q}}{m^{2q}\varepsilon^{2q}}.
\end{equation}
If upper bounds $\widehat\tau_{2p}$ on
$\tau_{2p}=r^{-1}\E\tr X_1^{2p}$ are available, with
$\widehat\tau_2\le\nu$ and $\widehat\tau_{2p}\le b_{2p}$, the stronger
discrete certificate replaces $u^{(n)}$ by
\begin{align}\label{core:eq:refined-discrete}
 \widehat u_0^{(n)}&=1,\qquad
 \widehat u_1^{(n)}=n\widehat\tau_2,\nonumber\\
 \widehat u_p^{(n)}&=\widehat u_p^{(n-1)}
 +\binom{2p}{2}\nu\widehat u_{p-1}^{(n-1)}+\widehat\tau_{2p}
 +\sum_{j=3}^{2p-2}\binom{2p}{j}b_j
                 \mathcal I_{2p-j}(\widehat u^{(n-1)})
 \quad(2\le p\le q),
\end{align}
where $\widehat u_p^{(0)}=0$ for $p\ge1$.
\end{theorem}
\begin{proof}
Apply Lemma~\ref{closed:lem:step} to obtain \eqref{closed:eq:one-step}.
Set $M_{n,p}=r^{-1}\E\tr S_n^{2p}$ and $M_{n,0}=1$.
Cauchy--Schwarz over the joint sample and normalized spectral measure
gives
\[
 r^{-1}\E\tr|S_n|^{2k+1}
 \le\sqrt{M_{n,k}M_{n,k+1}}.
\]
Conditioning \eqref{closed:eq:one-step} on $S_{n-1}$ now yields
$M_{n,p}\le M_{n-1,p}+\Phi_p(M_{n-1})$.
Each $\Phi_p$ is nonnegative, componentwise nondecreasing, and depends
only on coordinates $0,\ldots,p-1$. Induction proves
$M_{n,p}\le u_p^{(n)}$.

For the refined update, the all-$X$ term is exactly $\tau_{2p}$.
When $p\ge2$, a word with $2p-1$ copies of $X$ has one copy of
$S_{n-1}$ and cyclic trace $\tr(S_{n-1}X^{2p-1})$. Its expectation is
zero because the two matrices are independent and $\E S_{n-1}=0$.
The second moment is $M_{n,1}=n\tau_2$.
These observations prove \eqref{core:eq:refined-discrete} by the same
induction, and show $\widehat u_p^{(n)}\le u_p^{(n)}$.

To prove the integral comparison, first record
\begin{equation}\label{core:eq:coefficient-inequalities}
 C_{p,j}\le C_{p+1,j},\qquad
 C_{p,j}\le(2p)^{(j-2)/2}.
\end{equation}
The coefficient formulas in the proof of
Theorem~\ref{closed:thm:moment} give these statements directly.
Even coefficients increase with $p$. For an odd index $j=2\ell+1$,
write $k=p-\ell\ge1$; the ratio on increasing $p$ is
$[(k+\ell)/k]\sqrt{(2k-1)/(2k+1)}\ge1$.
The smallest case $\ell=1$ reduces after squaring to $2k^2-1\ge0$.
The second bound follows from \eqref{closed:eq:pairing-bound} and
$h_p\le\sqrt{2p}$, with $p=1$ immediate.

Since $b_2=\nu>0$, $\phi_q(s)>0$ and $1/\phi_q(s)$ is integrable at
zero. At infinity $\phi_q(s)\sim\nu/(2s)$.
Thus $\Gamma_q$ increases continuously from zero to infinity and has a
unique inverse. With $T_q(0)=0$, it satisfies
$T_q'(t)=\phi_q(T_q(t))$ for $t>0$.

Construct continuous functions in increasing degree by
\[
 H_0(t)=1,\qquad H_p(0)=0,\qquad H_p'(t)=\Phi_p(H(t)).
\]
The triangular structure makes this a recursive integration of known
continuous functions. All components are nondecreasing, so
$H(n+1)\ge H(n)+\Phi(H(n))$ and $u_p^{(n)}\le H_p(n)$.
Put $B_0=1$ and $B_p(t)=\gmom_pT_q(t)^{2p}$ for $1\le p\le q$.
The common scale gives $\mathcal I_k(B)=\omega_kT_q^k$. Hence
\[
 \frac{\Phi_p(B)}{2p\gmom_pT_q^{2p-1}}
 =\sum_{j=2}^{2p}\frac{C_{p,j}b_j}{j!}T_q^{1-j}
 \le\phi_q(T_q)=T_q'.
\]
Thus $B_p'\ge\Phi_p(B)$. Induction in $p$ and integration from zero
give $H_p\le B_p$. Integration is justified even if derivatives are
unbounded at zero: on $[a,t]$ it is ordinary integration, and
$B_p(a)\to0$ as $a\downarrow0$; the positive derivative has integral
$B_p(t)$. This proves both bounds in \eqref{core:eq:finite-moment}.
Finally, $\norm{Z}^{2q}\le\tr Z^{2q}$ and Markov's inequality prove
\eqref{core:eq:finite-tail}.
\end{proof}
\begin{lemma}[Centered trace budgets from raw moments]
\label{core:lem:centered-profile}
Suppose $H=yy^{\mathsf T}$, $\E H=I_r$, and
$\E H^h\preceq a_hI_r$ for $2\le h\le2q$, with $a_0=a_1=1$.
Then valid centered trace budgets for
\eqref{core:eq:refined-discrete} are
\begin{equation}\label{core:eq:centered-profile}
 \widehat\tau_{2p}
 =1+\sum_{\substack{2\le h\le2p\\h\ \mathrm{even}}}
        \binom{2p}{h}a_h
   -\sum_{\substack{1\le h\le2p-1\\h\ \mathrm{odd}}}
        \binom{2p}{h}r^{h-1}\qquad(1\le p\le q).
\end{equation}
These satisfy $0\le\tau_{2p}\le\widehat\tau_{2p}\le b_{2p}$
for the positive shifted budgets
$b_j=\sum_{h=0}^j\binom jh a_h$ at $j\ge3$ and
$b_2=a_2-1$. In particular,
$\widehat\tau_2=a_2-1$ and
$\widehat\tau_4=a_4+6a_2-4r^2-3$.
\end{lemma}
\begin{proof}
Put $t=\tr H=\norm y^2$. The rank-one identity and the even power give
\[
 \tau_{2p}=1+\sum_{h=1}^{2p}(-1)^h\binom{2p}{h}\frac{\E t^h}{r}.
\]
For even $h$, the raw budget gives $\E t^h/r\le a_h$.
For odd $h$, Jensen and $\E t=r$ give
$\E t^h/r\ge r^{h-1}$. Applying these bounds with their respective
signs proves \eqref{core:eq:centered-profile}. The resulting expression
is nonnegative because it bounds the nonnegative actual moment.
It is no larger than the positive binomial budget at even orders
at least four; at second order it is exactly $a_2-1$.
\end{proof}

\subsubsection{Finite radial choices}
\label{core:sec:finite-radial}
Let $A_p$ bound all directional $2p$ moments of $y$. Define
\[
 P_{r,p}(\lambda)=\sum_{k=0}^p\binom pk\lambda^k
                     \frac{\gmom_k}{D_{r,k}},\qquad
 L_p(\lambda)=
 \begin{cases}
  (1+\lambda)^p,&0\le\lambda\le1/(p-1),\\
  p^p\lambda/(p-1)^{p-1},&\lambda\ge1/(p-1).
 \end{cases}
\]
The exact finite Gaussian calculation in Lemma~\ref{core:lem:radial}
gives, for every deterministic $\lambda\ge0$,
\begin{equation}\label{core:eq:finite-tilt}
 \E H^p\preceq a_p^{\rm tilt}(\lambda)I_r,\qquad
 a_p^{\rm tilt}(\lambda)=
 \frac{A_pD_{r,p}}{\gmom_p}
       \frac{P_{r,p}(\lambda)}{L_p(\lambda)}.
\end{equation}
Indeed, $t=\ip xy^2/\norm y^2\in[0,1]$ and
$L_p(\lambda)=\min_{0<t\le1}(1+\lambda t)^p/t$.
This retains all terms in the finite Gaussian polynomial.
Any evaluated tilt is a valid upper bound; finding its global minimum
is unnecessary for validity. Two further explicit candidates are
\begin{equation}\label{core:eq:holder-quad}
 a_p^{\rm H\ddot{o}l}=A_p(D_{r,p}/\gmom_p)^{(p-1)/p},\qquad
 a_2^{\rm quad}=A_2\frac{r+2+\sqrt{3r(r+2)}}6.
\end{equation}
For the first, H\"older bounds
$\E\norm y^{2p-2}\ip xy^2$ using
$\E\norm y^{2p}\le A_pD_{r,p}/\gmom_p$ and the directional moment.
For the second, minimize the quadratic tilted-Gaussian polynomial
after the weighted AM--GM step. At $r=1$ it equals $A_2$;
for every $p$ the scalar bound $a_p=A_p$ is valid at that rank.

\subsubsection{Explicit tensor budgets}

Let $A_p=[\theta\beta_p\Ksg^{2p}]^d$ for the sub-Gaussian
model. Lemma~\ref{core:lem:orthogonal-fourth} gives
\begin{equation}\label{core:eq:orthogonal-v}
 v_{\rm orth}=A_2+(r-1)\gamma(A_2)\le A_2r.
\end{equation}
Lemma~\ref{core:lem:sharp-sg-fourth} supplies the additional candidate
\begin{equation}\label{core:eq:sharp-sg-variance}
 a_2^{\rm sg}=\min\left\{c_{\rm sg}^d s_r,
      c_{\rm sg}^d+(r-1)\gamma(c_{\rm sg}^d)\right\}.
\end{equation}
The two entries follow, respectively, from the quadratic radial bound
and the orthogonal fourth-moment bound using the sharper scalar
constant. Retain the separate raw and variance budgets by choosing
\begin{align}\label{core:eq:tensor-profile}
 a_0=a_1&=1,\qquad
 a_2=\min\{v,v_{\rm orth},a_2^{\rm sg},a_2^{\rm dir},a_2^{\rm quad},
                   \inf_{\lambda\ge0}a_2^{\rm tilt}(\lambda)\},\nonumber\\
 a_k&=\min\left\{a_k^{\rm dir},a_k^{\rm H\ddot{o}l},
     \inf_{\lambda\ge0}a_k^{\rm tilt}(\lambda),
     \frac{k!}{2}v[\tfrac12\Ksg^{2d}k^{d-2}(r+k)]^{k-2}\right\},
        \quad3\le k\le2q,\nonumber\\
 \nu&=a_2-1,\qquad b_2=\nu,\qquad
 b_j=\sum_{k=0}^j\binom jk a_k\quad(3\le j\le2q).
\end{align}
At $r=1$, also take the minimum with $A_k$. Any finite list of
chosen tilts may replace the infima. Each candidate is an upper bound,
so minimizing and positive binomial expansion preserve validity.
If $\nu=0$, then $\E\norm{H-I}_\F^2=0$ and $Z=0$ almost surely;
otherwise Theorem~\ref{core:thm:finite} applies directly.

All these raw budgets are no larger than the factorial envelope
with radius coefficient $1/2$. Its positive shift has parameters $(v,\Rmom_q)$.
Decreasing any budget decreases $\phi_q$, increases $\Gamma_q$, and
decreases $T_q(m)$. Therefore the finite-profile allowance
\begin{equation}\label{core:eq:profile-allowance}
 \eta_qc_q\frac{T_q(m)}m
\end{equation}
is no larger than \eqref{core:eq:explicit-closed}.
Indeed, the dominating factorial profile satisfies
$\phi_q^{\rm fact}(s)\le v/[2(s-h_q\Rmom_q)]$ above $h_q\Rmom_q$,
so integration gives $T_q(m)\le\sqrt{vm}+h_q\Rmom_q$.
For $q=1$, the respective scales are $\sqrt{\nu m}$ and $\sqrt{vm}$.
The discrete certificate and centered terminal budgets in
Lemma~\ref{core:lem:centered-profile} can improve this further.

\subsubsection{Exact scalar profiles for standard factor laws}
\label{core:sec:law-profiles}

The sub-Gaussian parameter is unnecessary when the factor law has an
explicit directional moment profile.
The Gaussian and spherical formulas are classical; the Rademacher
flat extremizer is due to Eaton and Komorowski
\cite{Eaton1970,Komorowski1988}. We include an elementary proof at even
integer moments, the orders needed here. For an integer $p\ge1$, set
\begin{equation}\label{core:eq:standard-local-profiles}
 \begin{aligned}
 c^{\rm G}_{n,p}&=\gmom_p,\\
 c^{\rm sph}_{n,p}&=\gmom_p\frac{n^p}{D_{n,p}},\\
 c^{\rm Rad}_{n,p}&=\frac1{2^n n^p}
   \sum_{a=0}^n\binom na(n-2a)^{2p}.
 \end{aligned}
\end{equation}
These are the sharp directional $2p$-moment constants for, respectively,
a standard Gaussian factor, a factor uniform on the sphere of radius
$\sqrt n$, and a factor with independent Rademacher coordinates.
In each case $c_{n,1}=1$.
In particular,
\[
 c^{\rm Rad}_{n,2}=3-\frac2n,\qquad
 c^{\rm Rad}_{n,3}=15-\frac{30}n+\frac{16}{n^2},\qquad
 c^{\rm sph}_{n,2}=\frac{3n}{n+2}.
\]
Thus dimension-one sign and spherical modes contribute no moment loss.

\begin{lemma}[Sharp product scalar budgets]
\label{core:lem:standard-product-profiles}
Choose any of the three laws independently in each mode, let
$A_p=\prod_{j=1}^d c_{n_j,p}^{(j)}$, and form the product probe $z$.
Then
\[
 \E|\ip a z|^{2p}\le A_p\norm a^{2p}.
\]
The constant $A_p$ is attained by an appropriate rank-one unit tensor.
Consequently the direction-preserving radialization lemma can be
applied with this $A_p$ in place of
$[\theta\beta_p\Ksg^{2p}]^d$ at every required integer order.
\end{lemma}
\begin{proof}
The Gaussian identity is classical. The spherical identity follows by
factoring a Gaussian into its independent radius and direction:
$\E R^{2p}=D_{n,p}$ and $\E G_1^{2p}=\gmom_p$.

For the sign law, signs of deterministic coefficients can be absorbed
into the independent Rademacher variables. Fix two nonnegative
coefficients $a,b$ with $c=a^2+b^2$, and condition on the other
coordinates, whose sum is $s$. With independent signs $\epsilon,\eta$,
\[
 \E(a\epsilon+b\eta)^{2k}
 =\frac{(c+2ab)^k+(c-2ab)^k}{2}.
\]
This is a polynomial in $a^2b^2$ with nonnegative coefficients, and
is strictly increasing when $k\ge2$ and $c>0$.
Expanding $\E(s+a\epsilon+b\eta)^{2p}$ retains only even local
moments, with nonnegative coefficients
$\binom{2p}{2k}s^{2p-2k}$.
For $p\ge2$, averaging unequal squared coefficients therefore strictly
increases the moment while preserving the squared norm. A maximizer
exists by compactness; every squared coefficient at a maximizer must
equal $1/n$. The binomial law of $\sum_{a=1}^n\epsilon_a$ gives the
last formula in \eqref{core:eq:standard-local-profiles}.
The case $p=1$ is isotropy.

For tensorization, slice a coefficient tensor in its last factor.
The local directional bound leaves a power of the sum of squared
partially contracted coefficients. Minkowski in $L^p$, followed by
induction, bounds this by $A_p\norm a^{2p}$.
A rank-one tensor made of local maximizing directions attains the
product bound by independence.
\end{proof}

\begin{proposition}[Gaussian and sign factorial radius]
\label{core:prop:gaussian-factorial}
Suppose $d\ge2$ and each factor is standard Gaussian or has independent
Rademacher coordinates, independently across modes. With
$v_{\rm G}=3^dr$, for every integer $p\ge2$,
\begin{equation}\label{core:eq:gaussian-factorial}
 \E H^p\preceq\frac{p!}{2}v_{\rm G}
 \left[\frac{11}{5}\,2^{d-2}p^{d-2}(r+p)\right]^{p-2}I_r.
\end{equation}
For every integer $q\ge1$, the closed certificates
\eqref{core:eq:explicit-closed}--\eqref{core:eq:explicit-m} therefore
apply with $(v,\Rmom_q)$ replaced by $(v_{\rm G},\Rmom_q^{\rm G})$, where
\begin{equation}\label{core:eq:gaussian-radius}
 \Rmom_q^{\rm G}=
 \frac{11}{5}\,2^{d-2}(2q)^{d-2}(r+2q)
       +1+\frac4{3v_{\rm G}}.
\end{equation}
The raw right side of \eqref{core:eq:gaussian-factorial} is also an
additional candidate in the finite profile. Spherical factors satisfy
the same envelope; their exact radial multipliers below can improve it.
\end{proposition}
\begin{proof}
All the stated laws have directional $2p$ moments bounded by
$A_p=\gmom_p^d$. For the sign comparison, expand an even moment of a
linear form. Only even coordinate powers survive, and every
corresponding Gaussian coordinate moment is at least one. The
factorwise induction then gives the product bound. This argument at
arbitrary orders requires full coordinate independence in the sign
case; four-wise independence alone supplies the fourth-moment bounds.
For spheres, $n^p/D_{n,p}\le1$ gives the same comparison.

The case $p=2$ follows from the directional fourth moment and
$\E H^2\preceq3^drI_r$. For $p\ge3$, use the H\"older candidate
in \eqref{core:eq:holder-quad} and
$D_{r,p}\le r(r+p)^{p-1}$, which is AM--GM on the final $p-1$
factors of $D_{r,p}$. Relative to the right side of
\eqref{core:eq:gaussian-factorial}, with $11/5$ replaced temporarily
by one, the resulting ratio is at most
\[
 \left[\frac{\gmom_p}{3(2p)^{p-2}}\right]^{d-2}
 \frac{2\gmom_p^{1+1/p}}{9p!}
             \left(\frac{r+p}{r}\right)^{1/p}.
\]
The first factor is at most one: retain $1\cdot3$ in $\gmom_p$ and
bound each of its remaining $p-2$ factors by $2p$. Since $r\ge1$,
the second factor is bounded by
\[
 G_p=\frac29\frac{\gmom_p}{p!}
                         [\gmom_p(p+1)]^{1/p}.
\]
At $p=3,4$, the needed inequalities are the exact rational comparisons
\[
 G_3^3=\frac{2500}{243}\le(11/5)^3,
 \qquad G_4^4=525(35/36)^4\le(11/5)^8.
\]
For $p\ge4$, we have $[\gmom_p(p+1)]^{1/p}\le6p/5$.
Indeed, $\gmom_p(p+1)/p^p$ decreases for $p\ge2$, because its
consecutive ratio is
\[
 \frac{(2p+1)(p+2)p^p}{(p+1)^{p+2}}\le1.
\]
To see this, the first three binomial terms give
$(1+1/p)^p\ge2+(p-1)/(2p)$, whereas
$(2p+1)(p+2)/(p+1)^2=2+p/(p+1)^2$
and $(p-1)/(2p)\ge p/(p+1)^2$ for $p\ge2$.
The bound at $p=4$ is $525\le(24/5)^4$; it persists since
$(6/5)^p$ increases.

Consequently $G_p\le(4/15)p\gmom_p/p!$ for $p\ge4$.
At $p=5$ this upper bound is $21/2\le(11/5)^3$, and its
consecutive ratio is $2+1/p\le11/5$ for $p\ge5$.
Thus $G_p\le(11/5)^{p-2}$ for all $p\ge3$.
This proves \eqref{core:eq:gaussian-factorial};
Lemma~\ref{core:lem:shift} gives \eqref{core:eq:gaussian-radius}.
\end{proof}

For a purely spherical product, the exact Gaussian coupling also
improves the whole uncentered matrix profile. With
\begin{equation}\label{core:eq:spherical-beta-profile}
 \varsigma_p=\prod_{j=1}^d\frac{n_j^p}{D_{n_j,p}},
\end{equation}
and the same fixed isometry $U$, let $H_{\rm sph}=y_{\rm sph}
y_{\rm sph}^\top$ and $H_{\rm G}=y_{\rm G}y_{\rm G}^\top$.
Independent Gaussian radii give the exact matrix identity
\begin{equation}\label{core:eq:spherical-matrix-profile}
 \E H_{\rm sph}^{p}=\varsigma_p\E H_{\rm G}^{p}.
\end{equation}
Indeed, $y_{\rm G}=a y_{\rm sph}$ with an independent scalar
$a=\prod_jR_j/\sqrt{n_j}$,
and $\E a^{2p}=\varsigma_p^{-1}$.
In particular, a Gaussian raw matrix budget $a_p^{\rm G}$ gives the
spherical raw budget $\varsigma_pa_p^{\rm G}$, while at second order
\[
 \E(H_{\rm sph}-I_r)^2
 =\varsigma_2\E H_{\rm G}^2-I_r.
\]
The positive binomial budgets and the centered trace-profile formulas
can then be rebuilt using these improved raw inputs.
There is no claim that the same scalar factor multiplies every
centered moment or the centered sum; those contain several different
homogeneous degrees.

For Gaussian/sign factors, Lemma~\ref{sharp:lem:gaussian-bilinear}
permits the additional raw second-moment candidate
$3^d+(r-1)(3^d-1)/2$, whose centered version is
$[(3^d-1)/2](r+1)$.
At tensor order two, Lemma~\ref{core:lem:order-two-wick} gives the
additional raw candidate $3r+4+2\sqrt r$. Thus one may use
\[
 a_2\le\min\{4r+5,\,3r+4+2\sqrt r\},\qquad
 \nu=a_2-1\le\min\{4r+4,\,3r+3+2\sqrt r\}.
\]
Here $a_2$ denotes the minimum of the applicable raw candidates,
including any smaller radial or exact geometric profile already
available. For a known Gaussian subspace,
\eqref{core:eq:wick-exact} gives the raw matrix exactly.
For spherical factors multiply each applicable Gaussian raw
candidate by $\varsigma_2$ before subtracting the identity.
The coordinate-class fourth-moment conclusions need only four-wise
independence; arbitrary-order sign profiles still require full
coordinate independence.
All higher radial formulas hold with these law-specific $A_p$;
the universal sub-Gaussian factorial envelope remains an optional
additional candidate when a valid $\Ksg$ is supplied.


\subsection{Spectral consequences and accuracy-sensitive lower bounds}
\label{appendix:spectral-extensions}
\begin{proposition}[Spectrum at linear sketch size]\label{core:prop:conditioning}
Assume the factor model of Theorem~\ref{core:thm:ose}, with $d\ge2$.
Let $\Cfour=c_{\rm sg}^d$, with $c_{\rm sg}$ from
Lemma~\ref{core:lem:sharp-sg-fourth}, a valid product directional
fourth-moment constant. For a fixed isometry $U\in\R^{N\times r}$ put
$M=U^\top\Omega\Omega^\top U$, $m=\lceil Cr\rceil$, $C\ge1$, and
$t=\log(8r/\delta)$, where $0<\delta<1$.
There is a constant $B_{d,\Ksg}$ such that, whenever
\[
 m\ge\mathfrak m_{\rm inj}(\Cfour,r;1/2,\delta/2),
\]
with probability at least $1-\delta$,
\[
 \lambda_{\min}(M)\ge\tfrac12,
 \qquad
 \lambda_{\max}(M)\le1+B_{d,\Ksg}
 \left[\sqrt{\frac tC}+\frac{t^{d-1}}C+\frac{t^d}{Cr}\right].
\]
Writing $\operatorname{cond}(B)=\sigma_{\max}(B)/\sigma_{\min}(B)$ for a full-column-rank matrix,
for fixed $d,\Ksg$ and sufficiently large fixed $C$, this implies
$\operatorname{cond}(\Omega^\top U)=O_{d,\Ksg,C}((\log r)^{(d-1)/2})$
with probability tending to one as $r\to\infty$.
For Gaussian factors (with $\Ksg^2\ge8/3$), the common-factor witness
attains this order:
\[
 \operatorname{cond}(\Omega^\top U_r)
 =\Theta_{d,\Ksg,C}((\log r)^{(d-1)/2})
 \quad\text{with probability }1-o(1).
\]
\end{proposition}
\begin{proof}
Apply \eqref{core:eq:explicit-closed} with failure $\delta/2$ and
$q=\lceil\log(8r/\delta)\rceil$. The proof of the upper OSE theorem gives
the displayed spectral allowance without imposing that it be below one.
The lower-injection bound spends at most $\delta/2$, so their
intersection has probability at least $1-\delta$.
For the asymptotic statement use failure $r^{-1}$ in the upper estimate.
For the lower edge, use \eqref{app:eq:injection-simple} with
$\alpha=1/2$, $\eta=e^{-r}$, and $s=r$. Its sufficient count is
at most $120\Cfour r$ for $r\ge1$. Thus a fixed
$C\ge120\Cfour$ makes $M\succeq I_r/2$ with failure at most $e^{-r}$.
The resulting upper edge is $O((\log r)^{d-1})$ for every $d\ge2$.
On the Gaussian common-factor witness, the existing one-heavy-column
argument gives $\lambda_{\max}(M)\ge c_d(\log r)^{d-1}/C$ with
probability $1-o(1)$. Also $r^{-1}\tr M\to1$ in probability, since its
variance is $[3^{d-1}(1+2/r)-1]/m$; hence
$\lambda_{\min}(M)\le2$ with probability $1-o(1)$. Take square roots
of the ratio of spectral edges.
\end{proof}

\begin{proposition}[Accuracy-sensitive Gaussian upper obstruction]\label{lower:prop:accuracy}
Fix $d\ge2$ and put $p=d-1$. On the Gaussian common-factor subspace
let $M=m^{-1}\sum_{i=1}^m w_i\xi_i\xi_i^\top$, with independent
$\xi_i\sim N(0,I_r)$ and $w_i=\prod_{j=1}^p\eta_{ij}^2$, where all
$\eta_{ij}$ are standard normal. For $r\ge4$, $m\ge r$ sufficiently
large, and $L_m=(\log m/p)^p$, there is a constant $a_d>0$ such that
\[
 \Pp\left\{\lambda_{\max}(M)-1\ge\frac{rL_m}{8m}\right\}
 \ge1-\exp\left[-a_d\frac{\sqrt m}{(\log m)^{p/2}}\right]
       -e^{-r/16}-\frac{64(3^d-1)m}{r^2L_m^2}.
\]
Consequently, for $\varepsilon_r\in(0,1)$ with
$\varepsilon_r r(\log r)^p\to\infty$, uniformly over integer sizes
\[
 r\le m\le\frac{r(\log r)^p}{16p^p\varepsilon_r},
 \qquad
 \Pp\{\norm{M-I_r}>\varepsilon_r\}\longrightarrow1.
\]
Sizes $m<r$ fail deterministically. Thus in this accuracy regime any
fixed-confidence OSE on this witness needs
$m=\Omega_d(r(\log r)^{d-1}/\varepsilon_r)$.
\end{proposition}
\begin{proof}
Choose $J$ maximizing $w_i$ using only the weights and set
$x=\xi_J/\norm{\xi_J}$. An exact identity is
\[
 m[x^\top Mx-1]=w_J(\norm{\xi_J}^2-1)+Q,
 \quad
 Q=\sum_iw_i-m+
   \sum_{i\ne J}w_i(\ip{\xi_i}{x}^2-1).
\]
Conditional on all weights and $\xi_J$, the second sum defining $Q$
has mean zero and variance $2\sum_{i\ne J}w_i^2$. The first sum has
variance $m(3^p-1)$, and the cross covariance vanishes by conditional
expectation. Therefore $\E Q=0$ and $\E Q^2\le(3^d-1)m$.
Chebyshev gives
$\Pp\{Q<-rL_m/8\}\le64(3^d-1)m/(r^2L_m^2)$.
The existing Gaussian extreme-weight argument gives $w_J\ge L_m$
outside the stated first failure term. Selection by the weights leaves
$\xi_J$ standard Gaussian, so its norm is at least $\sqrt{r/2}$ outside
probability $e^{-r/16}$. Since $r/2-1\ge r/4$, intersecting these three
events proves the finite bound. In the specified size range the spectral
increase is at least $2\varepsilon_r$, and the last failure term is
$O_d([\varepsilon_r r(\log r)^p]^{-1})$. The other two terms tend to
zero uniformly: $m\ge r$, and the accuracy condition bounds $m$ above
by $o(r^2(\log r)^{2p})$, so $\log m=O_d(\log r)$.
\end{proof}

\subsection{Common-Schmidt transition and spherical comparison}
\label{geom:appendix}
The arguments here use classical Gaussian moments, scalar exponential
estimates, an Euclidean net, and Gaussian radius--direction independence.
Weighted Gaussian covariance estimates have a broader theory in
\cite{CHZ2022}. We give the elementary special case needed for
Section~\ref{geom:subsec:geometry}; the Gaussian/spherical range
invariance in~\cite[Section~5.4]{CEMT2025} is consistent with the polar
comparison below.

\paragraph{Common-Schmidt row law and exact variance.}
For an order-two support with orthonormal basis
\[
 v_a=\sum_{b=1}^s\sqrt{\lambda_b}\,e_{a,b}\otimes f_b,
 \quad a=1,\ldots,r,
\]
the $rs$ vectors $e_{a,b}$ are jointly orthonormal and the $f_b$ are
orthonormal. Zero weights may be omitted. Conditional on the second
Gaussian factor, the projections have independent centered Gaussian
coordinates with common variance $w=\sum_b\lambda_bh_b^2$. Dividing by
$\sqrt w$ shows that the resulting $g$ is independent of $w$, since its
conditional law is constant. Here $w>0$ almost surely and
\[
 \E w=1,\qquad \E w^2=1+2\rho,
 \qquad \rho=\sum_b\lambda_b^2,\quad L=\max_b\lambda_b.
\]
For a standard $r$-dimensional Gaussian,
$\E(gg^\top)^2=(r+2)\Id_r$: its diagonal entries are
$\E[g_a^2\sum_bg_b^2]=3+(r-1)$ and its off-diagonal entries vanish.
Expansion of $(wgg^\top-\Id_r)^2$ proves
\eqref{geom:eq:variance}. In particular,
\begin{equation}\label{geom:eq:frob-variance}
 \E\left\|m^{-1}\sum_i y_iy_i^\top-\Id_r\right\|_\F^2
 =\frac r m\bigl[r+1+2(r+2)\rho\bigr],
 \qquad \E\ip uy^4=3(1+2\rho)\norm u^4.
\end{equation}
These are exact formulas, including $r=1$ and all Schmidt spectra.

\paragraph{A conditional weighted-Gaussian certificate.}
\begin{lemma}\label{geom:lem:weighted-gaussian}
Let $a_1,\ldots,a_m\ge0$ be deterministic, and let $g_i$ be independent
standard $r$-dimensional Gaussians. Define
$A_2=\sum_i a_i^2$, $B=\max_i a_i$, $\bar a=m^{-1}\sum_i a_i$, and
$t=r\log9+\log(2/\alpha)$ for $0<\alpha<1$. Then, with probability at
least $1-\alpha$,
\begin{equation}\label{geom:eq:weighted-certificate}
 \left\|m^{-1}\sum_i a_ig_ig_i^\top-\bar a\Id_r\right\|
 \le\frac{4\sqrt{A_2t}+4Bt}{m}.
\end{equation}
\end{lemma}
\begin{proof}
For any fixed unit $u$, $\ip u{g_i}$ are independent standard real
Gaussians. The Gaussian density gives
$\E e^{zG^2}=(1-2z)^{-1/2}$ for $z<1/2$. Expansion of the logarithm
therefore gives, for $|z|<(2B)^{-1}$,
\[
 \log\E\exp\left[z\sum_i a_i(\ip u{g_i}^2-1)\right]
 \le\frac{A_2z^2}{1-2B|z|}.
\]
If $B=0$ the conclusion is deterministic. Otherwise the exponential
Markov inequality, using
$z=\sqrt t/(\sqrt{A_2}+2B\sqrt t)$ and its negative, gives
\[
 \Pp\left\{\left|\sum_i a_i(\ip u{g_i}^2-1)\right|
 >2\sqrt{A_2t}+2Bt\right\}\le2e^{-t}.
\]
A $1/4$-net of the unit sphere has at most $9^r$ points, by disjoint
radius-$1/8$ balls and volume comparison. For any symmetric matrix $D$,
approximating a maximizing unit vector by a net point gives
$\norm D\le2\max_{u\text{ in net}}|u^\top Du|$. A union bound proves
\eqref{geom:eq:weighted-certificate}. This proof does not require a
matrix exponential bound.
\end{proof}

\begin{proof}[Proof of Theorem~\ref{geom:thm:schmidt}]
\phantomsection\label{proof:schmidt}
\emph{Upper bound.}
Put
\[
 u=\log(8/\delta),\quad v=\log(4m/\delta),\quad
 t=r\log9+u,\quad b=1+2\sqrt{\rho v}+2Lv.
\]
The Gaussian exponential moment, applied to the weights $\lambda_b$,
gives
\begin{align*}
 \Pp\{\max_iw_i>b\}&\le\delta/4,\\
 \Pp\{|m^{-1}\textstyle\sum_iw_i-1|>
             2\sqrt{\rho u/m}+2Lu/m\}&\le\delta/4.
\end{align*}
Let $\mu_2=1+2\rho$. Convexity gives
$\mu_4:=\E w^4\le\sum_b\lambda_b\E h_b^8=105$.
The independent variables $w_i^2\mathbf1_{\{w_i\le b\}}$ are bounded
by $b^2$, have means at most $\mu_2$, and variances at most $\mu_4$.
Scalar Bernstein therefore gives, except with probability
$e^{-u}=\delta/8$, and on the maximum-weight event,
\begin{equation}\label{geom:eq:weight-squares}
 \sum_iw_i^2\le A_m:=m\mu_2+\sqrt{2m\mu_4u}+\frac23b^2u.
\end{equation}
Condition on the weights and apply Lemma~\ref{geom:lem:weighted-gaussian}
with $\alpha=\delta/4$. A union bound gives, with probability at least
$1-\delta$,
\begin{equation}\label{geom:eq:full-certificate}
 \left\|m^{-1}\sum_i y_iy_i^\top-I_r\right\|
 \le2\sqrt{\rho u/m}+\frac{2Lu}{m}
           +\frac{4\sqrt{tA_m}+4bt}{m}.
\end{equation}
This is a numerical certificate using the actual second and fourth
moments of $w$. With $s_k=\sum_b\lambda_b^k$, the latter is
\[
 \mu_4=1+12\rho+12\rho^2+32s_3+48s_4.
\]

To obtain \eqref{geom:eq:all-confidence}, assume $m\ge u$ and use
$\mu_2\le3$, $\mu_4\le105$, $u\le t$, and
$b\le2+3L\log(4m/\delta)$. The right side of
\eqref{geom:eq:full-certificate} is at most
\[
 C_0\left[\sqrt{\frac{\tau_\delta}{m}}+
              \frac{\tau_\delta[1+L\log(4m/\delta)]}{m}\right],
 \qquad \tau_\delta=r+\log(8/\delta),
\]
for a universal $C_0$. Let $a=\log(e\tau_\delta/(\varepsilon\delta))$ and
$M=C\tau_\delta(\varepsilon^{-2}+L\varepsilon^{-1}a)$.
Since $a\ge1$ and $L\varepsilon\le1$,
$\log(4M/\delta)\le C_1a+\log C$ for a universal $C_1$.
At $m=M$ the last display is bounded by
$\varepsilon[C_2/\sqrt C+C_3(1+\log C)/C]$.
Choose $C$ sufficiently large. Both terms decrease for $m\ge M$,
which proves the claimed all-confidence bound. Its fixed-parameter
specialization gives the upper transition uniformly over the weights.

\emph{Lower bound.}
Take $c=1/16$. If $m<r$, rank deficiency gives error at least one.
Otherwise choose a largest Schmidt coefficient and denote its Gaussian
coordinate in sample $i$ by $h_{i,*}$. Then $w_i\ge Lh_{i,*}^2$.
For $m\ge r$ sufficiently large, integration of the Gaussian density on
$[\sqrt{\log m},\sqrt{\log m}+1/\sqrt{\log m}]$ yields
\[
 \Pp\{h_{i,*}^2\ge\log m\}\ge
       c_0m^{-1/2}(\log m)^{-1/2}.
\]
Consequently, except with probability
$\exp[-c_0\sqrt m/\sqrt{\log m}]$, at least one such coordinate occurs.
Choose an index $J$ attaining $\max_ih_{i,*}^2$, using only the weights'
underlying $h$ variables. The Gaussian vector $g_J$ is still independent
and standard; exponential Markov gives
$\Pp\{\norm{g_J}^2<r/2\}\le e^{-r/16}$. Positivity then implies
\[
 \lambda_{\max}\left(m^{-1}\sum_i y_iy_i^\top\right)
 \ge\frac{rL\log m}{2m}.
\]
In the nonempty range $r\le m\le cr[1+L\log(er)]$, necessarily
$L\log(er)\ge1/c-1=15$. For large $r$,
$\log r\ge\tfrac12\log(er)$, and hence
\[
 \frac{rL\log m}{2m}
 \ge\frac{L\log(er)}{4c[1+L\log(er)]}
 \ge\frac{15}{4c\cdot16}>2.
\]
The error exceeds one. Both exceptional probabilities tend to zero
uniformly in this range, since $m\ge r$ and
$m\le cr[1+\log(er)]$. This proves the lower part, uniformly even when
the weights and their number vary with $r$. For flat weights,
$L=\rho=1/s$, giving the claimed transition.
\end{proof}

\paragraph{Spherical moment transfer.}
\phantomsection\label{proof:spherical-transfer}
Let $x_i^{(j)}=\sqrt{n_j}u_i^{(j)}$, with all $u_i^{(j)}$ independent
and uniform on their real unit spheres. Draw independent
$R_i^{(j)}\sim\chi_{n_j}$, independent of all directions, and put
$g_i^{(j)}=R_i^{(j)}u_i^{(j)}$. Gaussian polar decomposition makes these
independent standard Gaussian factors. For a fixed isometry
$U\in\R^{(\prod_jn_j)\times r}$, define
\[
 y_i^{\rm sph}=U^\top\bigotimes_j x_i^{(j)},\qquad
 y_i^{\rm G}=U^\top\bigotimes_j g_i^{(j)},\qquad
 a_i=\prod_j\frac{R_i^{(j)}}{\sqrt{n_j}}.
\]
Then $y_i^{\rm G}=a_iy_i^{\rm sph}$ and $\E[a_i^2\mid(u)]=1$.
For the centered sample covariance errors $Z_{\rm sph},Z_{\rm G}$,
\begin{equation}\label{geom:eq:conditional-jensen}
 \E[Z_{\rm G}\mid(u)]=Z_{\rm sph},\qquad
 \E\tr Z_{\rm sph}^{2p}\le\E\tr Z_{\rm G}^{2p}
 \quad(p\in\mathbb N).
\end{equation}
The second assertion is conditional Jensen for the convex function
$D\mapsto\tr|D|^{2p}=\norm D_{S_{2p}}^{2p}$ on symmetric matrices,
followed by integration. All moments exist. Thus the Gaussian centered
trace-moment proof and Markov step underlying
Theorem~\ref{core:thm:ose} apply without changing constants. Independence
of coordinates within a spherical factor is neither true nor needed.

For a single factor $x=\sqrt n u$, polar independence and
$\E R^{2p}=\prod_{j=0}^{p-1}(n+2j)$ give
\begin{equation}\label{geom:eq:spherical-polar}
 \E\ip ax^{2p}
 =\left(\prod_{j=0}^{p-1}\frac n{n+2j}\right)
       \E\ip ag^{2p}.
\end{equation}
In particular $\E xx^\top=\Id_n$ and
$\E\ip ax^4=\frac{3n}{n+2}\norm a^4$. For an individual tensor row,
the same homogeneous calculation gives an exact multiplicative correction
$\prod_j\prod_{k=0}^{p-1}n_j/(n_j+2k)$ for its $p$th uncentered
covariance moment. For the centered sum the comparison is Jensen's
inequality \eqref{geom:eq:conditional-jensen}; there is no assertion of
that same exact correction for the entire centered sum.

\paragraph{Spherical factors retain the worst-case logarithm.}
\phantomsection\label{proof:spherical-lower}
Fix $d\ge2$ and, for large $r$, put
$h=\lfloor\log r/(d-1)\rfloor$. Choose mode dimensions
$n_1=r$, $n_2=\cdots=n_d=h$ and the common-factor support
$Ue_a=e_a\otimes e_1\otimes\cdots\otimes e_1$.
The projected row is
\[
 y_i=\left(\prod_{j=2}^d\sqrt h\,u_{i,1}^{(j)}\right)\xi_i,
 \qquad \xi_i=\sqrt r\,u_i^{(1)},\qquad \norm{\xi_i}^2=r.
\]
We give an elementary lower bound for a large spherical coordinate.
Write $u=G/\norm G$ in dimension $h$. The independent events
$G_1^2\ge h$ and $\sum_{a=2}^hG_a^2\le2h$ imply $u_1^2\ge1/3$.
The second has probability at least $1/2$, by Markov's inequality.
Integrating the first on $[\sqrt h,\sqrt h+1/\sqrt h]$ gives
\[
 \Pp\{u_1^2\ge1/3\}\ge c_1h^{-1/2}e^{-h/2}.
\]
With independence across common modes, a row has all $d-1$ such
coordinates with probability at least
$c_d h^{-(d-1)/2}e^{-(d-1)h/2}
\ge c_d h^{-(d-1)/2}r^{-1/2}$. Across $m\ge r$ independent rows, the
probability none occurs is at most
\[
 \exp[-c_d\sqrt r/h^{(d-1)/2}],
\]
which tends to zero. For an occurring row,
$\norm{y_i}^2\ge r(h/3)^{d-1}$. As
$h\ge\log r/[2(d-1)]$ for large $r$, every
\[
 r\le m\le\frac{r(\log r)^{d-1}}{2[6(d-1)]^{d-1}}
\]
then has sample-covariance top eigenvalue at least two. For $m<r$ the
rank argument applies. The failure probability for distortion $1/2$
tends to one uniformly throughout this entire range. Together with
\eqref{geom:eq:conditional-jensen}, this proves the spherical worst-case
order $\Theta_{d,\varepsilon,\delta}(r\log^{d-1}(er))$ as $r\to\infty$,
with tensor order, distortion, and confidence fixed.


\section{Approximation details}\label{appendix:approximation}
\subsection{Lower injection under fourth moments}\label{appendix:injection}

\begin{lemma}[Product fourth moments]\label{app:lem:product}
If independent isotropic factors $x_j$ obey directional fourth-moment
bounds $c_j$, their product $z=x_1\otimes\cdots\otimes x_d$ obeys
\eqref{app:eq:fourth} with $\Cfour=\prod_j c_j$. Independent centered, unit-variance coordinates with fourth
moment $\mu_4$ have $c_j=\max\{3,\mu_4\}$, by expanding the fourth power.
Thus Gaussian and sign factors admit $c_j=3$, including four-wise
independent sign coordinates, while standardized sparse signs of density
$\zeta_j$ admit $c_j=\max\{3,\zeta_j^{-1}\}$.
Under \eqref{core:eq:model}, one may instead use $c_j=c_{\rm sg}$
from Lemma~\ref{core:lem:sharp-sg-fourth}.
For fully independent signs and normalized spherical factors,
the respective exact constants are $3-2/n_j$ and $3n_j/(n_j+2)$;
their higher moments are given in Appendix~\ref{core:sec:law-profiles}.
\end{lemma}
\begin{proof}
Use the Hilbert-valued extension: for coefficient vectors $a_b$,
\[
 \E\left(\sum_b\ip{a_b}{x_j}^2\right)^2
 \le c_j\left(\sum_b\norm{a_b}^2\right)^2,
\]
by Cauchy--Schwarz on each pair of components. Successive contractions
of the coefficient tensor, conditioning on one independent factor at a
time, give the claimed product constant; the same induction gives
isotropy.
For centered independent coordinates $\xi_b$ with unit variance and
$\E\xi_b^4\le\mu_4$, expansion gives
\[
 \E\Bigl(\sum_b a_b\xi_b\Bigr)^4
 =3\Bigl(\sum_ba_b^2\Bigr)^2+
   \sum_b(\E\xi_b^4-3)a_b^4
 \le\max\{3,\mu_4\}\norm a^4.
\]
The same expansion uses only fourth-order joint moments, which also
justifies the four-wise-independent sign example.
\end{proof}


The choice $\beta=2(s+2)$ gives
\begin{equation}\label{app:eq:injection-simple}
 \mathfrak m_{\rm inj}(C,s;\alpha,\eta)
 \le\frac{14C}{3(1-\alpha)^2}
       \left(s+2+2\log\frac2\eta\right).
\end{equation}

At fixed $\eta$,
\[
 \mathfrak m_{\rm inj}(C,r;1/2,\eta)
 =\left[8\left(1+\frac2{\sqrt3}\right)+o(1)\right]Cr
 \qquad(r\to\infty).
\]
The lower-covariance method and rate are classical; the profile supplies
finite constants for the raw iid sketch.

\paragraph{Optimizing the injection level.}
For a deterministic failure allocation $0<\eta<\delta$, set
\[
 h_\eta=\tfrac14\mathfrak m_{\rm inj}(\Cfour,r;1/2,\eta),\qquad
 k_\eta=\frac{B_r}{\varepsilon(\delta-\eta)}.
\]
The proof of Theorem~\ref{app:thm:relative}, with lower level $\alpha$,
requires $m\ge\max\{h_\eta/(1-\alpha)^2,k_\eta/\alpha^2\}$.
For $B_r>0$, the minimizing level
$\alpha=\sqrt{k_\eta}/(\sqrt{h_\eta}+\sqrt{k_\eta})$ gives
\begin{equation}\label{app:eq:relative-balanced}
 m\ge(\sqrt{h_\eta}+\sqrt{k_\eta})^2.
\end{equation}
The allocation $\eta$ can also be optimized before sampling. The same
formula applies to Corollary~\ref{app:cor:input} with
$(\Cfour,B_r)$ replaced by $(C_Q,\kappa(Q,W))$; zero interaction
requires only injection.


\begin{lemma}[Gaussian-smoothed lower covariance]
\label{app:lem:lower-smoothing}
Let $v_1,\ldots,v_m\in\R^k$ be iid, with
$V=\E vv^\top$, $0\preceq V\preceq I_k$, and
\[
 \E\ip\vartheta v^4\le C(\vartheta^\top V\vartheta)^2
       \quad\text{for every }\vartheta\in\R^k,
 \qquad C\ge1.
\]
Put $s=\tr V$. For $0<\alpha,\eta<1$, if
$m\ge\mathfrak m_{\rm inj}(C,s;\alpha,\eta)$, then
\[
 \Pp\left\{\frac1m\sum_i v_iv_i^\top+I_k-V
                       \succeq\alpha I_k\right\}\ge1-\eta.
\]
\end{lemma}
\begin{proof}
We use Gaussian smoothing for lower covariance tails
\cite{Oliveira2016}, with Catoni's logarithmic influence function
\cite{Catoni2012}. If $s=0$, then $v=0$ almost surely and the assertion
is deterministic. Otherwise fix $\beta,\lambda>0$ before sampling and set
\[
 \varphi(t)=\log(1+t+t^2/2),\qquad t\ge0.
\]
Its derivative is $(1+t)/(1+t+t^2/2)$, while
$\varphi''(t)=-(t+t^2/2)/(1+t+t^2/2)^2$. Thus $\varphi$ is concave,
$0\le\varphi'\le1$, and
\[
 e^{-\varphi(t)}\le1-t+t^2/2,
 \qquad e^{\varphi(t)}=1+t+t^2/2.
\]
The first inequality follows because the two quadratic polynomials
have product $1+t^4/4$. Write $\mu_\vartheta=\vartheta^\top V\vartheta$.
Independence, the fourth-moment hypothesis, and $1+x\le e^x$ imply
\[
 \E\exp Z_\vartheta\le1,\qquad
 Z_\vartheta=m\lambda\mu_\vartheta
  -\frac{m\lambda^2C}{2}\mu_\vartheta^2
  -\sum_i\varphi(\lambda\ip\vartheta{v_i}^2).
\]
Use the Gaussian prior $\pi=N(0,\beta^{-1}I_k)$. Fubini and Markov give
$\int e^{Z_\vartheta}\,d\pi(\vartheta)\le2/\eta$ with probability
at least $1-\eta/2$. On this event, the change-of-measure inequality
\[
 \int Z_\vartheta\,d\nu(\vartheta)
 \le\operatorname{KL}(\nu\Vert\pi)
                   +\log\int e^{Z_\vartheta}\,d\pi(\vartheta)
\]
holds simultaneously for all probability measures $\nu\ll\pi$ for
which the displayed integrals are finite. It follows from Jensen's
inequality after inserting $d\pi/d\nu$.
Take $\pi_u=N(u,\beta^{-1}I_k)$ for any unit $u$; its relative entropy
is $\beta/2$. Gaussian moment identities give
\begin{align*}
 \E_{\pi_u}\mu_\vartheta
  &=u^\top Vu+s/\beta,\\
 \E_{\pi_u}\mu_\vartheta^2
  &=(u^\top Vu)^2+
       \frac{2s\,u^\top Vu+4u^\top V^2u}{\beta}
       +\frac{s^2+2\tr V^2}{\beta^2}\\
  &\le1+\frac{2(s+2)}\beta+\frac{s(s+2)}{\beta^2}.
\end{align*}
The last step uses $V\preceq I_k$, $u^\top Vu\le1$,
$u^\top V^2u\le1$, and $\tr V^2\le s$. Consequently, simultaneously
for all unit $u$,
\begin{align*}
 \frac1{m\lambda}\sum_i\E_{\pi_u}
            \varphi(\lambda\ip\vartheta{v_i}^2)
 &\ge u^\top Vu+\frac s\beta
 -\frac{\lambda C}{2}
       \left(1+\frac{2(s+2)}\beta+\frac{s(s+2)}{\beta^2}\right)
 -\frac{\beta/2+\log(2/\eta)}{m\lambda}.
\end{align*}

Concavity and $\varphi'\le1$ remove the Gaussian smoothing:
\[
 \E_{\pi_u}\varphi(\lambda\ip\vartheta{v_i}^2)
 \le\varphi\!\left(\lambda\ip u{v_i}^2+
                       \lambda\norm{v_i}^2/\beta\right)
 \le\lambda\ip u{v_i}^2+
                       \varphi(\lambda\norm{v_i}^2/\beta).
\]
The common radial term admits an upper exponential bound under only
fourth moments. For an independent $g\sim N(0,I_k)$,
\[
 3\E\norm v^4=\E_g\E\ip gv^4
 \le C\E_g(g^\top Vg)^2
 =C(s^2+2\tr V^2)\le Cs(s+2).
\]
For $R_i=\varphi(\lambda\norm{v_i}^2/\beta)/\lambda$, the identity for
$e^\varphi$, independence, and Markov therefore give, with probability
at least $1-\eta/2$,
\[
 \frac1m\sum_iR_i
 \le\frac s\beta+\frac{\lambda C s(s+2)}{6\beta^2}
                          +\frac{\log(2/\eta)}{m\lambda}.
\]
Intersect the two events and subtract the common radial term:
\[
 \frac1m\sum_i\ip u{v_i}^2
 \ge u^\top Vu-\frac{\lambda C}{2}A_s^{\rm inj}(\beta)
                  -\frac{B_\eta^{\rm inj}(\beta)}{m\lambda}
 \quad\text{for every unit }u.
\]
Here the $\beta^{-2}$ coefficient is $1+1/3=4/3$, incorporating the
radial upper bound. Choosing
$\lambda=\sqrt{2B_\eta^{\rm inj}(\beta)/(mCA_s^{\rm inj}(\beta))}$ gives the uniform
deficit $\sqrt{2CA_s^{\rm inj}(\beta)B_\eta^{\rm inj}(\beta)/m}$.
The minimizing $\beta$ is deterministic. The stated budget makes this
deficit at most $1-\alpha$, so adding $I_k-V$ proves the conclusion.

All expectations are justified under the stated fourth moments.
The radial exponential is a quadratic polynomial in $\norm v^2$,
and Fubini applies to the nonnegative prior integrand. Gaussian posterior
integrals are finite by their polynomial moment bounds. Thus no
exponential moment of the raw covariance sample is assumed.
\end{proof}

\begin{proof}[Proof of Lemma~\ref{app:lem:injection}]
\phantomsection\label{proof:injection}
Apply Lemma~\ref{app:lem:lower-smoothing} with
$v_i=Q^\top z_i$ and $V=I_r$.

For completeness, the minimizing $\beta$ is the unique positive root of
\begin{equation}\label{app:eq:injection-cubic}
 \beta^3-(b+2aT)\beta-4bT=0,\qquad
 a=2(s+2),\quad b=\frac43s(s+2),\quad T=2\log(2/\eta).
\end{equation}
Indeed,
$A_s^{\rm inj}B_\eta^{\rm inj}=\beta/2+T+a/2+(aT+b/2)/\beta+bT/\beta^2$
is strictly convex on $(0,\infty)$, and its derivative equation is
the displayed cubic. At fixed $\eta$ as $s\to\infty$, its minimizer
satisfies $\beta/s\to2/\sqrt3$, and the objective divided by $s$
tends to $1+2/\sqrt3$. This proves the stated leading coefficient.
For each fixed $\beta$, $A_s^{\rm inj}(\beta)$ is nondecreasing in $s$.
Minimizing over $\beta$ preserves this order, so an upper bound
on the effective dimension is sufficient.
\end{proof}

\subsection{Bilinear moments and order-two Wick geometry}\label{appendix:wick}

\begin{lemma}[Orthogonal fourth moments]\label{core:lem:orthogonal-fourth}
Let $y\in\R^s$ be isotropic and satisfy
$\E\ip uy^4\le c\norm u^4$ for all $u$, with $c\ge1$. Define
\begin{equation}\label{core:eq:orthogonal-gamma}
 \gamma(c)=\min\{c/2,c-1\}.
\end{equation}
For orthogonal unit vectors $u,v$,
\begin{equation}\label{core:eq:orthogonal-fourth}
 \E\ip uy^2\ip vy^2\le\gamma(c).
\end{equation}
Consequently,
\begin{equation}\label{core:eq:orthogonal-variance}
 \E(yy^{\mathsf T})^2\preceq
 [c+(s-1)\gamma(c)]I_s.
\end{equation}
\end{lemma}
\begin{proof}
Put $A=\ip uy$, $B=\ip vy$, and $h=\E A^2B^2$.
The fourth-moment bounds in the directions $(u+v)/\sqrt2$ and
$(u-v)/\sqrt2$ give
\[
 \E A^4+\E B^4+6h\le4c.
\]
Cauchy--Schwarz gives $\E A^4+\E B^4\ge2h$, hence $h\le c/2$.
Isotropy and Jensen also give
$\E A^4+\E B^4+2h=\E(A^2+B^2)^2\ge4$, hence $h\le c-1$.
To prove the matrix assertion, extend a unit test vector to an
orthonormal basis. Its own fourth moment contributes at most $c$ to
$\E\norm y^2\ip uy^2$, and each of the other $s-1$ coordinates
contributes at most $\gamma(c)$.
\end{proof}
The following choices for $B_r$ in \eqref{app:eq:aggregate-budget}
follow from the bilinear and Wick calculations below; take the minimum
of applicable bounds.
\begin{center}\small
\begin{tabularx}{\linewidth}{@{}Yl@{}}
\toprule Probe law & Aggregate bound $B_r$\\\midrule
Directional fourth moment $\Cfour$ & $r\gamma(\Cfour)$\\
Gaussian or four-wise-independent sign factors & $r(3^d-1)/2$\\
Independent factors with four-wise-independent coordinates & $r(C_{\rm coord}-1)/2$\\
Order two, coordinate fourth moments at most three & $\Xi_r=\min\{4r,3r+2\sqrt r\}$\\
Two normalized spherical factors & $\varsigma_2\Xi_r$\\
\bottomrule
\end{tabularx}
\end{center}
In the coordinate row, coordinates are centered with variance one and
$C_{\rm coord}=\prod_j\max\{3,\mu_{4,j}\}$, where $\mu_{4,j}$ bounds
their fourth moments. Factors are mutually independent throughout.
For spheres, $\varsigma_2=\prod_{j=1}^2n_j/(n_j+2)$.
The leading coefficient three in $\Xi_r$ is sharp for Gaussian factors.
These are aggregate bounds; dividing $B_r$ by $r$ need not give a
valid bound for each individual pair of directions.


The bilinear calculations also retain subspace geometry at second order.
For two Gaussian factors, Lemma~\ref{core:lem:order-two-wick} gives
an exact formula using the partial traces and partial transpose of
the head projector, and the uniform centered bound
\[
 \E(H-I_r)^2\preceq
 \min\{4r+4,3r+3+2\sqrt r\}I_r.
\]
These smaller variance inputs enter the finite-profile recurrence
without changing its higher-order budgets.

\begin{lemma}[Centered bilinear moment]
\label{sharp:lem:bilinear}
For an isotropic random vector $z$ satisfying the directional fourth
moment bound with constant $C\ge1$, put
\[
 \gamma(C)=\min\{C/2,C-1\},\qquad
 \gamma_{\parallel}(C)=C-1-\gamma(C)\ge0.
\]
Then for all deterministic vectors $u,v$,
\begin{equation}\label{sharp:eq:bilinear}
 \mathbb E\bigl(\langle u,z\rangle\langle v,z\rangle
                   -\langle u,v\rangle\bigr)^2
 \le\gamma(C)\|u\|^2\|v\|^2
       +\gamma_{\parallel}(C)\langle u,v\rangle^2.
\end{equation}
Consequently, for matrices $U,W$ with the same number of rows and a
normalized iid sketch with $m$ columns,
\[
 \mathbb E\|U^\top(\Omega\Omega^\top-I)W\|_F^2
 \le\frac{\gamma(C)\|U\|_F^2\|W\|_F^2
                   +\gamma_{\parallel}(C)\|U^\top W\|_F^2}{m}.
\]
\end{lemma}
\begin{proof}
By homogeneity it is enough to take $\|u\|=\|v\|=1$. Put
$A=\langle u,z\rangle$, $B=\langle v,z\rangle$,
$\chi=\langle u,v\rangle$, and $h=\mathbb E A^2B^2$.
The bounds in directions $u+v$ and $u-v$ give
\[
 \mathbb E A^4+\mathbb E B^4+6h\le4C(1+\chi^2).
\]
Since $\mathbb E A^4+\mathbb E B^4\ge2h$,
\[
 \operatorname{Var}(AB)=h-\chi^2
 \le\frac C2+\left(\frac C2-1\right)\chi^2.
\]
There is also the bound $\operatorname{Var}(AB)\le C-1$.
To see this when $|\chi|<1$, let
\[
 X=\frac{A+B}{\sqrt{2+2\chi}},\quad
 Y=\frac{A-B}{\sqrt{2-2\chi}},\quad
 a=\frac{1+\chi}{2},\quad b=\frac{1-\chi}{2}.
\]
Both $X$ and $Y$ have unit second moment and fourth moment at most $C$,
and $AB-\chi=a(X^2-1)-b(Y^2-1)$. The triangle inequality in $L^2$
gives $\|AB-\chi\|_2\le(a+b)\sqrt{C-1}=\sqrt{C-1}$.
The endpoint $|\chi|=1$ is the directional fourth-moment bound itself.
For $C\ge2$ use the preceding affine bound in $\chi^2$; for $C\le2$
use $C-1$. These are exactly the displayed coefficients $\gamma,\gamma_{\parallel}$.
Finally, expand the Frobenius norm into entries, remove distinct-sample
terms using centering and independence, and sum this scalar inequality.
\end{proof}

\begin{lemma}[Gaussian and sign product bilinear moments]
\label{sharp:lem:gaussian-bilinear}
Let $z=g^{(1)}\otimes\cdots\otimes g^{(d)}$, where the factors are
independent standard Gaussian vectors, and let
$\gamma_{\mathrm G}(d)=(3^d-1)/2$. Then for all deterministic $u,v$,
\[
 \mathbb E\bigl(\langle u,z\rangle\langle v,z\rangle
                   -\langle u,v\rangle\bigr)^2
 \le\gamma_{\mathrm G}(d)\bigl(\|u\|^2\|v\|^2+\langle u,v\rangle^2\bigr).
\]
The same upper bound holds when any or all Gaussian factors are replaced
by isotropic four-wise-independent sign vectors. In particular,
for orthogonal unit vectors $u,v$,
\[
 \mathbb E\langle u,z\rangle^2\langle v,z\rangle^2\le\gamma_{\mathrm G}(d).
\]
This orthogonal constant is optimal for Gaussian factors when every mode
dimension is at least two. Consequently, for every fixed isometry
$U\in\mathbb R^{N\times r}$ and $y=U^\top z$,
\begin{equation}\label{sharp:eq:gaussian-variance}
 \mathbb E(yy^\top-I_r)^2\preceq\gamma_{\mathrm G}(d)(r+1)I_r.
\end{equation}
\end{lemma}
\begin{proof}
The needed two-form Gaussian square-product inequality is a special
case of chaotic square-product inequalities~\cite{MNPP2015}; the present
decoupled case has an elementary induction.
If $P,Q$ are scalar forms, linear in each of the $d$ independent
Gaussian factors, then
\[
 \mathbb E P^2Q^2\ge\mathbb E P^2\,\mathbb E Q^2.
\]
The statement is immediate at $d=0$. For the induction step, condition
on all but the first factor and write $P=\sum_i g_i^{(1)}a_i$,
$Q=\sum_i g_i^{(1)}b_i$, where $a_i,b_i$ are multilinear forms in the
remaining factors. The Gaussian fourth-moment formula gives
\[
 \mathbb E_1[P^2Q^2]
 =\left(\sum_i a_i^2\right)\left(\sum_j b_j^2\right)
   +2\left(\sum_i a_ib_i\right)^2.
\]
Drop the last nonnegative term and apply the induction hypothesis to
each $\mathbb E a_i^2b_j^2$. The resulting sum is
$\mathbb E P^2\,\mathbb E Q^2$.

Normalize $u,v$ to be unit vectors and put $A=\langle u,z\rangle$,
$B=\langle v,z\rangle$, $\chi=\langle u,v\rangle$, and
$h=\mathbb E A^2B^2$. Apply the square-product inequality to $A+B,A-B$:
\[
 \mathbb E A^4+\mathbb E B^4-2h
 =\mathbb E(A+B)^2(A-B)^2\ge4(1-\chi^2).
\]
The product directional fourth-moment bound gives
\[
 \mathbb E A^4+\mathbb E B^4+6h\le4\cdot3^d(1+\chi^2).
\]
Subtracting yields
$h\le(3^d-1)/2+(3^d+1)\chi^2/2$; subtracting $\chi^2$ proves the
centered inequality.

For a four-wise-independent sign vector $\epsilon$, direct expansion gives
\[
 \mathbb E\langle a,\epsilon\rangle^2\langle b,\epsilon\rangle^2
 =\|a\|^2\|b\|^2+2\langle a,b\rangle^2-2\sum_i a_i^2b_i^2
 \le\|a\|^2\|b\|^2+2\langle a,b\rangle^2.
\]
Conditioning on all other factors therefore allows one sign factor to
be replaced by a Gaussian factor without decreasing the raw mixed
fourth moment. Repeat over all sign factors; the second moment and
$\chi$ remain unchanged.

For the matrix conclusion, expand
$\mathbb E\|y\|^2\langle v,y\rangle^2$ in an orthonormal basis beginning
with any fixed unit vector $v$. Its own coordinate contributes at most
$3^d$, and each other coordinate contributes at most
$\gamma_{\mathrm G}(d)$. Subtract $I_r$ and use
$3^d-1=2\gamma_{\mathrm G}(d)$.

For optimality, let $P=\prod_{j=1}^d g_1^{(j)}$ and
$Q=\prod_{j=1}^d g_2^{(j)}$, which are independent, have unit second
moment, and have fourth moment $3^d$. The forms
$A=(P+Q)/\sqrt2$, $B=(P-Q)/\sqrt2$ have orthogonal unit coefficient
tensors and satisfy
\[
 \mathbb E A^2B^2
 =\frac14\mathbb E(P^2-Q^2)^2
 =\frac{3^d-1}{2}.
\]
\end{proof}

\begin{lemma}[Product-coordinate bilinear moments]
\label{app:lem:coordinate-cross}
Suppose the independent factors have centered, unit-variance coordinates,
four-wise independent within each factor, with coordinate fourth moments
at most $\mu_{4,j}$. Put
\[
 c_j=\max\{3,\mu_{4,j}\},\qquad C_{\rm coord}=\prod_{j=1}^d c_j,
 \qquad \Gamma_{\rm coord}=(C_{\rm coord}-1)/2.
\]
Then for all deterministic vectors $u,v$,
\begin{equation}\label{app:eq:coordinate-bilinear}
 \E\bigl(\ip uz\ip vz-\ip uv\bigr)^2
 \le\Gamma_{\rm coord}
       \bigl(\norm u^2\norm v^2+\ip uv^2\bigr).
\end{equation}
In particular, for orthogonal unit $u,v$,
\begin{equation}\label{app:eq:coordinate-cross}
 \E\ip uz^2\ip vz^2\le\Gamma_{\rm coord}.
\end{equation}
This extends Lemma~\ref{sharp:lem:gaussian-bilinear} beyond Gaussian
and sign coordinates. Its Gaussian orthogonal constant is sharp as
shown there.
\end{lemma}
\begin{proof}
For a centered, unit-variance coordinate vector $\xi$,
four-wise independence gives
\[
 \E(a^\top\xi)^2(b^\top\xi)^2
 =\norm a^2\norm b^2+2\ip ab^2
   +\sum_i(\E\xi_i^4-3)a_i^2b_i^2.
\]
The expression is nondecreasing in every coordinate fourth moment.
For factor $j$, replace its coordinates by independent copies of
$\sigma G$, where $G$ is standard Gaussian, $\sigma$ is independent
of $G$, and $\sigma^2=c_j/3$ with probability $3/c_j$ and zero
otherwise. This comparison coordinate has variance one and fourth
moment $c_j$. Condition on the other factors and replace factors
one at a time. The raw mixed square cannot decrease, while the
second moments and $\ip uv$ remain unchanged.

For the comparison product law, all excess coordinate fourth moments
are nonnegative. Consequently the induction used in
Lemma~\ref{sharp:lem:gaussian-bilinear} remains valid: for any two
multilinear forms $F,G$,
$\E F^2G^2\ge\E F^2\,\E G^2$. Indeed, after conditioning on one
factor, discard the nonnegative inner-product square and the
nonnegative excess-fourth-moment terms; apply the induction
hypothesis to every coefficient pair in the remaining sum.
The comparison directional fourth moment is at most $C_{\rm coord}$
by Lemma~\ref{app:lem:product}.
For unit $u,v$, put $A=\ip uz$, $B=\ip vz$, $\chi=\ip uv$, and
$h=\E A^2B^2$, now under that comparison law. The two bounds give
\[
 \E A^4+\E B^4-2h\ge4(1-\chi^2),\qquad
 \E A^4+\E B^4+6h\le4C_{\rm coord}(1+\chi^2).
\]
Subtract and then subtract $\chi^2$ to obtain
$h-\chi^2\le(C_{\rm coord}-1)(1+\chi^2)/2$.
Comparison and homogeneity prove the statement for the original law.
\end{proof}

\begin{lemma}[Order-two Wick geometry and variance]
\label{core:lem:order-two-wick}
Let $z=g\otimes h$ for independent standard real Gaussian vectors.
For an orthogonal rank-$r$ projector $P$, write
$P_A=\operatorname{Tr}_2P$, $P_B=\operatorname{Tr}_1P$, and let
$P^\Gamma$ denote partial transpose in the second factor. Then
\begin{equation}\label{core:eq:wick-ambient}
 M_P:=\E[(z^\top Pz)zz^\top]
 =rI+2(P_A\otimes I+I\otimes P_B)+2P+2P^\Gamma.
\end{equation}
In particular, if $P=UU^\top$, $U^\top U=I_r$, and
$H=(U^\top z)(U^\top z)^\top$, then
\begin{equation}\label{core:eq:wick-exact}
 \E H^2=(r+2)I_r+
       2U^\top(P_A\otimes I+I\otimes P_B+P^\Gamma)U.
\end{equation}
Consequently,
\begin{equation}\label{core:eq:wick-variance}
 \E H^2\preceq(3r+4+2\sqrt r)I_r,\qquad
 \E(H-I_r)^2\preceq(3r+3+2\sqrt r)I_r.
\end{equation}
The two inequalities, but not necessarily the Wick identities, remain
valid for independent factors with centered, unit-variance coordinates,
four-wise independent within each factor, whose fourth moments are at
most three. For two normalized spherical factors, multiply the
Gaussian raw fourth-moment bounds by
$\varsigma_2=\prod_{j=1}^2 n_j/(n_j+2)$ before subtracting $I_r$.
\end{lemma}
\begin{proof}
The nine real Gaussian Wick contractions prove
\eqref{core:eq:wick-ambient}. The pairing internal to each factor
contributes $rI$; the four pairings internal in one factor contribute
twice each partial trace; the remaining four contribute
$2P+2P^\Gamma$. Compression by $U$ proves
\eqref{core:eq:wick-exact}.

The double-reduction construction of universal state inversion
\cite[Eq.~(2.11)]{Rungta2001} has the positivity property
\begin{equation}\label{core:eq:double-reduction}
 rI-P_A\otimes I-I\otimes P_B+P\succeq0.
\end{equation}
Here is a direct proof of the matrix inequality. Let
$J_{ab}=e_ae_b^\top-e_be_a^\top$. Entrywise expansion gives
\[
 \mathscr D_n(X):=(\tr X)I-X
       =\sum_{a<b}J_{ab}X^\top J_{ab}^\top.
\]
Thus $(\mathscr D_{n_1}\otimes\mathscr D_{n_2})(P)$ equals
\[
 \sum_{a<b}\sum_{c<e}
 (J_{ab}\otimes J_{ce})P^\top
                    (J_{ab}\otimes J_{ce})^\top\succeq0.
\]
It is a sum of positive congruences of $P^\top=P$, and its expansion
is the left side of \eqref{core:eq:double-reduction}.
The formula also covers one-dimensional factors, when the appropriate
sum is empty.

Compression of \eqref{core:eq:double-reduction} bounds the sum of the
two partial-trace terms by $(r+1)I_r$.
Partial transpose preserves the Frobenius norm, so
$\norm{P^\Gamma}\le\norm{P^\Gamma}_\F=\sqrt r$.
Substitute in \eqref{core:eq:wick-exact}, then use
$\E(H-I_r)^2=\E H^2-I_r$.

For the coordinate class, test $\E H^2$ in any direction and expand
the remaining squared norm in an orthonormal basis. Each summand is
a mixed square. The conditional fourth-moment formula in
Lemma~\ref{app:lem:coordinate-cross} shows that replacement by
Gaussian factors cannot decrease it. This proves Loewner domination.
For spheres, Gaussian polar decomposition gives the exact
homogeneous relation $\E H_{\rm sph}^2=\varsigma_2\E H_{\rm G}^2$.
\end{proof}

\begin{corollary}[Aggregate cross moment for order-two probes]
\label{app:cor:order-two-cross}
Under the order-two coordinate assumptions of
Lemma~\ref{core:lem:order-two-wick}, for every isometry
$Q\in\R^{N\times r}$ and every unit vector $v\perp\range Q$,
\begin{equation}\label{app:eq:order-two-cross}
 \E\norm{Q^\top z}^2\ip vz^2\le\Xi_r,\qquad
 \Xi_r=\min\{4r,\,3r+2\sqrt r\}.
\end{equation}
The leading coefficient three is sharp for Gaussian factors.
For two normalized spherical factors the bound is
$\varsigma_2\Xi_r$, with $\varsigma_2$ as above.
\end{corollary}
\begin{proof}
Gaussian comparison in Lemma~\ref{core:lem:order-two-wick} reduces
the coordinate class to Gaussian factors. Set $P=QQ^\top$.
Since $Pv=0$, double reduction gives
$v^\top(P_A\otimes I+I\otimes P_B)v\le r$.
Testing \eqref{core:eq:wick-ambient} therefore gives
$v^\top M_Pv\le3r+2\sqrt r$.
The separate orthogonal mixed-square bound
$\Gamma_{\rm coord}=4$ gives $v^\top M_Pv\le4r$.
The spherical conclusion follows from the same homogeneous relation.

For sharpness of the leading coefficient, take the head space
spanned by $e_1\otimes e_1,\ldots,e_r\otimes e_1$ and the orthogonal
unit tail direction $v=e_{r+1}\otimes e_1$. Independence gives
\[
 \E\norm{Q^\top z}^2\ip vz^2
 =\E\left(\sum_{a=1}^r g_a^2\right)g_{r+1}^2\,\E h_1^4
 =3r.
\]
Thus no uniform leading coefficient smaller than three is possible.
\end{proof}

\subsection{PSD trace and confidence amplification}
\label{appendix:amplification}
\begin{corollary}[PSD Nystr\"om trace-relative error]
\label{app:cor:nystrom}
Let $K\in\R^{N\times N}$ be positive semidefinite, let
$\tau_r^{\rm tr}(K)=\sum_{j>r}\lambda_j(K)$, and define
\[
 \mathcal N=\mathcal N_K(\Omega),\qquad \mathcal N_r=(\mathcal N)_r.
\]
Under \eqref{app:eq:relative-budget}, with probability at least $1-\delta$,
\begin{equation}\label{app:eq:nystrom-relative}
 \tr(K-\mathcal N)\le(1+\varepsilon)\tau_r^{\rm tr}(K),\qquad
 \tr(K-\mathcal N_r)\le(1+\varepsilon)\tau_r^{\rm tr}(K).
\end{equation}
\end{corollary}
\begin{proof}
For $A=K^{1/2}$ and $P=P_{A\Omega}$, the standard Nystr\"om
representation $\mathcal N=APA$~\cite[Proposition~A.2]{TYUCPSD2017}
identifies the two trace errors with the projection and best rank-$r$
squared Frobenius errors in $\range P$. Apply Theorem~\ref{app:thm:relative}
to $A$, whose optimal squared tail is $\tau_r^{\rm tr}(K)$.
\end{proof}

\begin{proof}[Proof of Corollary~\ref{app:cor:amplification}]
\phantomsection\label{proof:amplification}
Partition the columns into $b$ independent groups of size $m_0$. Theorem
\ref{app:thm:relative} with confidence parameter $1/4$ makes each group
successful with probability at least $3/4$. Overall normalization changes
no range. The concatenated range contains every group range, so both the
projection error and the optimal rank-$r$ error in it are no larger than
the corresponding group errors. The probability that all groups fail is
at most $4^{-b}\le\delta$. Apply the same argument to $K^{1/2}$ for the
Nystr\"om statements.
\end{proof}
Independent regression solves
can instead be amplified by selecting the candidate with smallest original
residual, when evaluating that residual is permitted.


\subsection{High-confidence Hilbert comparisons}\label{appendix:hilbert}
\begin{proof}[Proof of Theorem~\ref{app:thm:subgauss-confidence}]
Let $Q$ be an isometry with $r$ columns and let $Q^\top W=0$.
If $F=\norm W_\F>0$, regard
$Y=(Q^\top z)(W^\top z)^\top/F$ as a vector in the Frobenius
Hilbert space. Isotropy gives $\E Y=0$, and
Lemma~\ref{app:lem:cross} gives $\E\norm Y_\F^2\le B_r$.
Minkowski's inequality applied to a sum of squares shows that
\[
 \E\norm{Tz}^{2j}\le A_j\norm T_\F^{2j}.
\]
Cauchy--Schwarz therefore gives
$\E\norm Y_\F^j\le r^{j/2}A_j=b_j^H$.

For a centered Hilbert vector $Y$ and a fixed vector $x$, expand
$\norm{x+sY}^{2p}=(\norm x^2+2s\ip xY+s^2\norm Y^2)^p$.
The linear coefficient has expectation zero. The absolute value of
every other coefficient is at most the corresponding coefficient of
$(\norm x+s\norm Y)^{2p}$. Thus the even moments of
$\sum_{i\le n}Y_i$ satisfy the scalar recurrence
\eqref{closed:eq:update}, with $b_2=\nu_H$ and the budgets $b_j^H$.
The pairing comparison in Appendix~\ref{appendix:closed} applies
without a trace or a rank factor and yields
\[
 \left(\E\left\|m^{-1}\sum_{i=1}^mY_i\right\|_\F^{2q}\right)^{1/(2q)}
 \le c_q\left(\sqrt{\nu_H/m}+h_q\Rmom_q^H/m\right).
\]
Markov's inequality and \eqref{app:eq:high-confidence} bound failure
of the interaction event by $\delta/2$. Lower injection fails with
probability at most $\delta/2$. The same deterministic competitors
and normal equations as in Theorem~\ref{app:thm:relative} now apply.
If $W=0$, the interaction vanishes.

For the displayed rate, use the coarser valid budget
$\nu_H=Cr$, where $C=[2\theta\Ksg^4]^d$, and the original
factorial scalar bound $A_j\le\theta^d(j!)^d\Ksg^{2dj}$.
The elementary inequality $j!\le2j^{j-2}$ gives
\[
 b_j^H\le\frac{j!}{2}Cr
   \bigl[\Ksg^{2d}\sqrt r\,(2q)^{d-1}\bigr]^{j-2}.
\]
Choose $q=\lceil\log(2/\delta)\rceil$. The resulting interaction
allowance is at most
$C_{d,\Ksg}[\sqrt{rq/m}+\sqrt r\,q^d/m]$.
Making its two terms at most $\sqrt\varepsilon/4$ proves
\eqref{app:eq:high-confidence-rate} and also covers the injection
budget after increasing $C_{d,\Ksg}$.
\end{proof}

\begin{proposition}[Keeping the weak Hilbert variance]
\label{app:prop:radial-hilbert}
Assume $B_r>0$ and retain the orthogonal mixed-square constant
$\Gamma_{\perp}>0$ above. Put $r_H=B_r/\Gamma_{\perp}>0$ and, for every real $s>0$,
define
\[
 D_{s,p}=\prod_{j=0}^{p-1}(s+2j),\qquad D_{s,0}=1.
\]
The interaction allowance in \eqref{app:eq:high-confidence} can be
decreased by taking its minimum with
\begin{equation}\label{app:eq:radial-hilbert}
 (2/\delta)^{1/(2q)}D_{r_H,q}^{1/(2q)}
       \left(\sqrt{\Gamma_{\perp}/m}+a_{r_H,q}^H/m\right).
\end{equation}
For $s>0$, set
\[
 \omega^s_{2k}=D_{s,k},\qquad
 \omega^s_{2k+1}=\sqrt{D_{s,k}D_{s,k+1}},\qquad
 C^s_{p,j}=\frac{(2p)!\,\omega^s_{2p-j}}
                    {2p(2p-j)!D_{s,p}},
\]
and take
\[
 a_{s,q}^H=\max_{\substack{2\le p\le q\\3\le j\le2p}}
    \left(\frac{2C^s_{p,j}b_j^H}{j!\Gamma_{\perp}}\right)^{1/(j-2)},
 \qquad a_{s,1}^H=0.
\]
The higher budgets remain $b_j^H=r^{j/2}A_j$, with the actual head
rank $r$. The real parameter $r_H$ enters only the comparison
pairing moments; $\Gamma_{\perp}$ remains a separate bound for the weak
Hilbert covariance. In particular, the aggregate bound $\Xi_r$
does not justify replacing $\Gamma_{\perp}$ by $\Xi_r/r$.
Both allowances bound the same even moment, so their deterministic
minimum needs only one application of Markov's inequality.
\end{proposition}
\begin{proof}
For a test matrix $M$ with $\norm M_\F=1$, take a singular value
decomposition $M=\sum_a\sigma_a u_av_a^\top$. The orthogonal
mixed-square bound and Cauchy--Schwarz give
\[
 \bigl(\E\ip M Y_\F^2\bigr)^{1/2}
 \le\frac{\sqrt{\Gamma_{\perp}}}{F}\sum_a\sigma_a\norm{Wv_a}
 \le\sqrt{\Gamma_{\perp}}.
\]
The covariance of $Y$ as a Hilbert vector therefore has operator
norm at most $\Gamma_{\perp}$, while its trace is
$\E\norm Y_\F^2\le B_r=\Gamma_{\perp} r_H$.
The quadratic coefficient in the Hilbert expansion is at most
\[
 p\norm x^{2p-2}\E\norm Y^2
 +2p(p-1)\norm x^{2p-4}\E\ip xY^2
 \le p\Gamma_{\perp}(r_H+2p-2)\norm x^{2p-2}.
\]
At $x=0$, use the polynomial expansion itself, so no negative
power is required. All higher coefficients use $b_j^H$.

Put $s(t)=\sqrt{\Gamma_{\perp} t}+a_{r_H,q}^H$ and
$\mathcal B_p(t)=D_{r_H,p}s(t)^{2p}$, with $\mathcal B_0=1$.
The identity $D_{r_H,p}/D_{r_H,p-1}=r_H+2p-2$ makes the normalized
quadratic term exactly $\Gamma_{\perp}/(2s)$.
By the definition of $a_{r_H,q}^H$, every higher term, divided by
$2pD_{r_H,p}s^{2p-1}$, is at most
$\Gamma_{\perp}(a_{r_H,q}^H/s)^{j-2}/(2s)$.
The resulting geometric sum is bounded by
$\Gamma_{\perp}/[2(s-a_{r_H,q}^H)]=s'(t)$.
The same monotone induction and integrable initial singularity as
in Appendix~\ref{appendix:closed} prove the comparison with
$\mathcal B_q(m)$. Only the displayed product identity for
$D_{r_H,p}$ is used, so $r_H$ need not be an integer.
For $q=1$, this is the direct second-moment bound.
\end{proof}

The leading variance coefficient in this comparison is
$\sqrt{\Gamma_{\perp}} D_{r_H,q}^{1/(2q)}$, at most
$\sqrt{B_r+\Gamma_{\perp}(q-1)}$ by AM--GM. The scalar comparison uses
$\sqrt{B_r}\,\gmom_q^{1/(2q)}$.
Take their minimum because the higher-insertion radii need not be
ordered. This refinement concerns the residual interaction; it
does not remove the variance logarithm from the two-sided embedding
theorem.

\subsection{Angular moments and full-range error}\label{appendix:angular}
\begin{proof}[Proof of Theorem~\ref{app:thm:angular-range}]
Let $z=\bigotimes_j x_j$ have the specified spherical factor law.
For one factor, the beta distribution of its first squared coordinate
gives
\[
 \E|\ip u{x_j}|^{2p}=c_{n_j}^{\rm sph}(p)\norm u^{2p}.
\]
This formula also holds at $n_j=1$. To tensorize, slice a coefficient
tensor as $a=\sum_b a_b\otimes e_b$ and condition on all earlier
factors. The last-factor moment formula and Minkowski's inequality in
$L^p$ yield
\[
 \left\|\sum_b|\ip{a_b}{z_{<d}}|^2\right\|_{L^p}
 \le\sum_b\bigl(\E|\ip{a_b}{z_{<d}}|^{2p}\bigr)^{1/p}.
\]
Induction therefore proves
\[
 \E|\ip a z|^{2p}\le H_{\rm ang}(p)\norm a^{2p},
 \qquad
 \E(z^\top Bz)^p\le H_{\rm ang}(p)(\tr B)^p
 \quad(B\succeq0).
\]
The second inequality follows by diagonalizing $B$ and applying
Minkowski once more. Differentiating \eqref{app:eq:angular-kappa}
at $p=1$ gives
\begin{equation}\label{app:eq:angular-limit}
 \lim_{\beta\downarrow0}
 H_{\rm ang}(1+\beta)^{1/\beta}=\varkappa_{\rm ang}.
\end{equation}

For a nonzero PSD residual $R$, set $a=z^\top R^2z$ and $b=z^\top Rz$.
The numerator vanishes whenever the denominator does.
Since $\E a=\tr(R^2)$, the probability measure
$dQ=a\,dP/\tr(R^2)$ is well defined and has $b>0$ almost surely.
For $\beta>0$ and $p=1+\beta$, convexity and H\"older give
\[
 \E\frac ab
 =\tr(R^2)\E_Qb^{-1}
 \ge\tr(R^2)(\E_Qb^\beta)^{-1/\beta},
 \qquad
 \E_Qb^\beta\le H_{\rm ang}(p)(\tr R)^\beta.
\]
For the last estimate, diagonalize $R^2=\sum_i\lambda_i v_iv_i^\top$
with unit $v_i$ and apply
\[
 \E[|\ip{v_i}z|^2b^\beta]
 \le(\E|\ip{v_i}z|^{2p})^{1/p}
       (\E b^p)^{\beta/p}
 \le H_{\rm ang}(p)(\tr R)^\beta
\]
before summing. Letting $\beta\downarrow0$ proves
\begin{equation}\label{app:eq:angular-drift}
 \E\frac{z^\top R^2z}{z^\top Rz}
 \ge\frac{\tr(R^2)}{\varkappa_{\rm ang}\tr R}.
\end{equation}
Zero quotients are interpreted as zero.

Write $K=A^\top A$ and let $Z_t$ contain the first $t$ probe columns.
The full Nystr\"om residual is
\[
 R_t=K-KZ_t(Z_t^\top KZ_t)^\dagger Z_t^\top K.
\]
It satisfies $0\preceq R_{t+1}\preceq R_t\preceq K$, and adding a
column gives
\[
 R_{t+1}=R_t-
  \frac{R_tz_{t+1}z_{t+1}^\top R_t}{z_{t+1}^\top R_tz_{t+1}}.
\]
These identities follow by writing the Nystr\"om approximation as
$K^{1/2}P_{\range(K^{1/2}Z_t)}K^{1/2}$ and decomposing the enlarged
range into its old part and the orthogonalized new column.
For the rectangular problem, cyclicity of trace also gives exactly
\[
 \tr R_t
 =\tr\bigl[A^\top(I-P_{\range(AZ_t)})A\bigr]
 =\norm{(I-P_{\range(AZ_t)})A}_\F^2.
\]

Let $f_t=\E\tr R_t$. The rank-$r$ tail of $R_t$ is no larger than
$\tau$, because $R_t\preceq K$. Thus
\[
 \tr(R_t^2)\ge\frac{(\tr R_t-\tau)_+^2}{r}.
\]
The function $g(x)=(x-\tau)_+^2/x$, with $g(0)=0$, is convex:
above $\tau$ it equals $x-2\tau+\tau^2/x$, and its derivative agrees
with zero at $\tau$. Conditioning \eqref{app:eq:angular-drift} on
$R_t$ and applying Jensen therefore yields
\[
 f_{t+1}\le f_t-\frac{(f_t-\tau)_+^2}
                              {\varkappa_{\rm ang}r f_t}.
\]
Before $u_t=f_t/\tau-1$ reaches $\varepsilon$, this implies
\[
 u_t-u_{t+1}\ge
       \frac{u_t^2}{\varkappa_{\rm ang}r(1+u_t)}.
\]
The decreasing convex function $F(u)=u^{-1}-\log u$ satisfies,
as long as $u_{t+1}>\varepsilon$,
\[
 F(u_{t+1})-F(u_t)
 \ge F'(u_t)(u_{t+1}-u_t)
 \ge\frac1{\varkappa_{\rm ang}r}.
\]
If the integer $m$ in \eqref{app:eq:angular-budget-exact} were reached
with $u_m>\varepsilon$, summing this inequality would give
$F(u_m)\ge F(\varepsilon)$, a contradiction. Once the target is
reached, monotonicity preserves it. This proves
\eqref{app:eq:angular-range-error}; using
$-u_0^{-1}+\log u_0\le\log(T/\tau)$ gives
\eqref{app:eq:angular-budget-simple}.

Finally, each Gaussian factor is a positive independent radius times
its spherical direction. Under this coupling a Gaussian product column
is a nonzero scalar multiple of the corresponding spherical product
column. All sampled ranges and all full Nystr\"om outputs are therefore
unchanged. This establishes the Gaussian assertion as well.
\end{proof}

\begin{corollary}[Expected Frobenius error of the full PSD residual]
\label{app:cor:angular-frobenius}
Let $K\succeq0$ be fixed, $T=\tr K$, and
$R_m=K-\mathcal N_K(\Omega_m)$ for the first $m$ product probes of
either law in Theorem~\ref{app:thm:angular-range}. Then, for every
integer $m\ge0$,
\begin{equation}\label{app:eq:angular-frobenius}
 \E\norm{R_m}_\F^2
 \le\frac{\varkappa_{\rm ang}T^2}{m+1}.
\end{equation}
This bound concerns the full PSD Nystr\"om residual, with no
rank-$r$ truncation.
\end{corollary}
\begin{proof}
\phantomsection\label{proof:angular-frobenius}
The case $T=0$ is immediate. For $T>0$, couple the probes to an
infinite independent sequence and put $f_t=\E\tr R_t$.
Since $\tr R_t\le T$, the conditional drift
\eqref{app:eq:angular-drift} gives
\[
 \E\tr(R_t^2)
 \le\varkappa_{\rm ang}T(f_t-f_{t+1}).
\]
Moreover $R_t-R_{t+1}\succeq0$ and $R_t+R_{t+1}\succeq0$, so
\[
 \tr(R_t^2)-\tr(R_{t+1}^2)
 =\tr[(R_t-R_{t+1})(R_t+R_{t+1})]\ge0.
\]
Sum the drift through $t=m$ and use this monotonicity:
\[
 (m+1)\E\tr(R_m^2)
 \le\sum_{t=0}^m\E\tr(R_t^2)
 \le\varkappa_{\rm ang}T(f_0-f_{m+1})
 \le\varkappa_{\rm ang}T^2.
\]
The independent continuation step is used only in the proof; the
output $R_m$ uses exactly the first $m$ probes.
\end{proof}

For comparison, $\varkappa_{\rm ang}\le\prod_j c_{n_j}^{\rm sph}(2)$ by
log-convexity and $H_{\rm ang}(1)=1$. The inequality is strict whenever
some $n_j>1$. At $n_j=2$, each factor contributes $e/2$ to
$\varkappa_{\rm ang}$ and $3/2$ to the spherical fourth-moment bound.
In fact, uniformly over the mode dimensions,
\[
 \varkappa_{\rm ang}\le
 \left(\frac{e^{2-\gamma_{\rm E}}}{2}\right)^d,
 \qquad \frac{e^{2-\gamma_{\rm E}}}{2}=2.0743278106\ldots.
\]
To see this, put $R^2\sim\chi_n^2$ and
$\gmom(p)=2^p\Gamma(p+1/2)/\Gamma(1/2)$.
Gaussian polar decomposition gives
$c_n^{\rm sph}(p)=\gmom(p)/\E(R^2/n)^p\le\gmom(p)$ for $p>1$.
Both sides equal one at $p=1$, so their right logarithmic derivatives
satisfy $(\log c_n^{\rm sph})'(1)\le\log2+\psi(3/2)=2-\gamma_{\rm E}-\log2$.
Each per-mode factor tends to this upper limit as $n\to\infty$,
whereas the Gaussian fourth-moment contribution is $3$.
These are sufficient constants for the specified full-range error criterion.


\subsection{A confidence obstruction for raw regression}
\begin{proposition}[Raw regression needs $1/(\varepsilon\delta)$ under fourth moments]
\label{app:prop:regression-lower}
Let $N\ge2$, $0<\varepsilon\le1/64$, $0<\delta<1/40$, and $\Lambda\ge400$.
For every integer $m$ with
\[
 1\le m<\frac{15\Lambda}{2048\,\varepsilon\delta},
\]
there are an isotropic probe law satisfying \eqref{app:eq:fourth} with
$\Cfour\le\Lambda+4$, a unit design $a\in\R^N$, and a unit response
$b\perp a$ such that the sketch-and-solve solution of
Theorem~\ref{app:thm:relative} from $m$ iid probes satisfies
\[
 \Pp\bigl\{\norm{a\widetilde x-b}^2>1+\varepsilon\bigr\}>\delta.
\]
The optimal squared residual is one. Hence any sketch size that is
sufficient uniformly over isotropic laws with $\Cfour\le\Lambda+4$ must be
at least $15\Lambda/(2048\varepsilon\delta)$ in this regime: for raw
regression, the term $8\gamma(\Cfour)r/(\varepsilon\delta)$ in
\eqref{app:eq:relative-budget} cannot be replaced by one logarithmic in
$1/\delta$.
\end{proposition}
\begin{proof}
First suppose $m\ge\Lambda/(4\varepsilon)$. Let $a=e_1$, $b=e_2$, and $w=(e_1+e_2)/\sqrt2$, and put
\[
 B^2=8\sqrt\varepsilon\,m,\qquad\pi=\Lambda/B^4,\qquad
 \Sigma=(1-\pi)^{-1}(I_N-\pi B^2ww^\top).
\]
Since $m\ge\Lambda/(4\varepsilon)$, we have $\pi\le\varepsilon/(4\Lambda)$,
$\pi B^2\le1/2$, and $m\pi\le1/16$. Each probe equals $\sigma Bw$ with
probability $\pi$, where $\sigma$ is an independent fair sign, and
otherwise equals $\Sigma^{1/2}g$ with $g\sim N(0,I_N)$. Then
$\E zz^\top=\pi B^2ww^\top+(1-\pi)\Sigma=I_N$, and for unit $u$,
\[
 \E\ip uz^4\le\pi B^4+3(1-\pi)\norm\Sigma^2
 \le\Lambda+\frac3{1-\pi}\le\Lambda+4.
\]

Write $a_i=\ip{e_1}{z_i}$ and $b_i=\ip{e_2}{z_i}$. The sketched problem
gives $\widetilde x=\sum_ia_ib_i/\sum_ia_i^2$ whenever the denominator is
positive, and the original squared residual is $1+\widetilde x^2$. It
therefore suffices to show $\widetilde x>\sqrt\varepsilon$ with probability
greater than $\delta$. Let $E$ be the event that exactly one probe, with
index $J$, is of the first kind. Then
$a_Jb_J=a_J^2=B^2/2=4\sqrt\varepsilon\,m$ and
\[
 \Pp(E)=m\pi(1-\pi)^{m-1}\ge m\pi(1-m\pi)\ge\tfrac{15}{16}m\pi.
\]
Conditional on $E$ and $J$, the other $m-1$ pairs $(a_i,b_i)$ are iid
centered Gaussians with $\E a_i^2=s(1-\alpha)$ and $\E a_ib_i=-s\alpha$,
where $s=(1-\pi)^{-1}\le1+2\pi$ and $\alpha=\pi B^2/2$. The sums
$S_{ab}=\sum_{i\ne J}a_ib_i$ and $S_{aa}=\sum_{i\ne J}a_i^2$ have
variances at most $ms^2$ and $2ms^2$. With $t=\sqrt\varepsilon\,m/2$,
Chebyshev's inequality shows that both lie within $t$ of their means,
except with conditional probability at most $12s^2/(\varepsilon m)\le1/2$.
On this event $ms\alpha=s\Lambda/(16\sqrt\varepsilon)\le t$, and since
$\sqrt\varepsilon\le1/8$,
\[
 \widetilde x=\frac{B^2/2+S_{ab}}{B^2/2+S_{aa}}
 \ge\frac{4\sqrt\varepsilon\,m-2t}{4\sqrt\varepsilon\,m+sm+t}
 \ge\frac{3\sqrt\varepsilon}{9/16+s}>\sqrt\varepsilon.
\]
Hence the failure probability is at least
$\tfrac12\Pp(E)\ge15\Lambda/(2048\varepsilon m)$, which exceeds $\delta$
in the stated range.


It remains to treat $m<\Lambda/(4\varepsilon)$. Put
$p_0=\max\{3/\Lambda,1/m\}$, and let a probe equal
$p_0^{-1/2}g$, with $g\sim N(0,I_N)$, with probability $p_0$,
and zero otherwise. It is isotropic and its directional fourth-moment
constant is $3/p_0\le\Lambda$. Keep $a=e_1$ and $b=e_2$.
The number of active probes is $K\sim\operatorname{Binomial}(m,p_0)$,
with $1\le\mu:=\E K<3/(4\varepsilon)$. Conditional on $K=k\ge1$,
the minimizer is unique almost surely and has the distribution
$G/\sqrt{X_k}$, where $G\sim N(0,1)$ and $X_k\sim\chi_k^2$ are
independent. Denote the standard Gaussian distribution function by
$\Phi_{\mathrm N}$, and set
$f(t)=2\Phi_{\mathrm N}(-\sqrt{\varepsilon t})$ for $t\ge0$.
This function is convex, since, for $t>0$,
\[
 f''(t)=\frac{\sqrt\varepsilon\,\phi_{\mathrm N}(\sqrt{\varepsilon t})}{2}
        (t^{-3/2}+\varepsilon t^{-1/2})>0,
\]
where $\phi_{\mathrm N}$ is the Gaussian density. Applying Jensen's
inequality first to $X_k$ and then to $K$, and discarding $K=0$, gives
\[
 \Pp\{|\widetilde x|>\sqrt\varepsilon\}
 \ge \E f(K)-\Pp(K=0)\ge f(\mu)-e^{-\mu}.
\]
For $1\le\mu\le4$, this is at least
$1-1/(2\sqrt{2\pi})-e^{-1}>1/4$; for $\mu>4$, it exceeds
$1-\sqrt{3/(2\pi)}-e^{-4}>1/4$.
Here we used $\phi_{\mathrm N}(t)\le1/\sqrt{2\pi}$,
$\varepsilon\le1/64$, and $\varepsilon\mu<3/4$.
Thus failure exceeds $\delta$ throughout the remaining range.
The event $K=0$ was not used, so the conclusion holds for every choice
of a minimizer on singular samples.
\end{proof}


\subsection{Regression under Gaussian scale mixtures}\label{appendix:weighted}
\begin{proof}[Proof of Proposition~\ref{weighted:prop}]
Use an orthonormal basis for $\range A$ and the unit residual
$e/\norm e$. Orthogonality and Gaussian rotational invariance give
independent design coordinates $\xi_i$ and residual coordinates
$\eta_i$. The normal equations and Pythagoras give
\eqref{weighted:eq:law}; $S_m$ is invertible almost surely for $m\ge r$.
Conditioning leaves only the independent Gaussian $\eta_i$, proving
\eqref{weighted:eq:covariance}. Conditional on $(w_i,\xi_i)$,
$\mathcal E_m^{\rm reg}$ has the law $\sum_j\lambda_j(C_m)G_j^2$. Gaussian
integration gives, for $0<s<(2\norm{C_m})^{-1}$,
\[
 \log\E[e^{s(\mathcal E_m^{\rm reg}-\tr C_m)}\mid(w_i,\xi_i)]
 \le\frac{s^2\tr(C_m^2)}{1-2s\norm{C_m}}.
\]
The exponential Markov inequality proves
\eqref{weighted:eq:conditional-tail}. The strong law gives
$S_m/m\to I_r$ and
$m^{-1}\sum_iw_i^2\xi_i\xi_i^\top\to\mu_2I_r$.
The asserted covariance limit follows, and the conditional Gaussian
characteristic functions give the distributional limit.

The projected scale-mixture row has directional fourth-moment constant
$C_w=3\mu_2$. The injection proof needs only isotropy and directional
fourth moments of the projected rows. Apply Lemma~\ref{app:lem:injection}
in $\R^r$, with isometry $I_r$ and rows $\sqrt{w_i}\xi_i$, to obtain
$S_m\succeq(m/2)I_r$ except with probability $\delta/2$.
On this event and $\max_iw_i\le b_m(\delta)$,
\[
 C_m\preceq b_m(\delta)S_m^{-1}
       \preceq\frac{2b_m(\delta)}m I_r.
\]
For standard $h\in\R^r$, the Gaussian norm bound
$\Pp\{\norm h>\sqrt r+\sqrt{2u}\}\le e^{-u}$, with
$u=\log(4/\delta)$, gives $\mathcal E_m^{\rm reg}\le
2b_m(\delta)(\sqrt r+\sqrt{2u})^2/m$ on the conditional event.
A union bound proves the claim. When $e=0$, the original residual is
zero almost surely for $m\ge r$, so no ratio is needed.
\end{proof}

For the common-factor Gaussian law, $\mu_2=3^{d-1}$ and one may take
\[
 b_m(\delta)=\bigl[2\log(8(d-1)m/\delta)\bigr]^{d-1}.
\]
Indeed, a union bound over the $(d-1)m$ scalar Gaussian factors using
$\Pp\{G^2>2v\}\le2e^{-v}$ controls every factor.
For the order-two common-Schmidt law, write
$\rho=\sum_b\lambda_b^2$ and $L=\max_b\lambda_b$. Then
\[
 \mu_2=1+2\rho,\qquad
 b_m(\delta)=1+2\sqrt{\rho\log(4m/\delta)}
                      +2L\log(4m/\delta),
\]
by the scalar Gaussian tail estimate in Appendix~\ref{geom:appendix}.
Thus the certificate has polylogarithmic confidence cost in these
two support classes. For example, in the common-factor case the
explicit sufficient order is
\[
 m\ge C_d\frac{r+\log(1/\delta)}{\varepsilon}
 \left[\log\frac{e(r+\log(1/\delta))}{\varepsilon\delta}\right]^{d-1}.
\]
To close the logarithm, insert this expression as $M$: its logarithm
is at most a $d$-dependent multiple of the bracketed logarithm, plus
$\log C_d$. Increasing $C_d$ makes every term in
\eqref{weighted:eq:budget} at most $M$; the function
$\log^{d-1}(cm/\delta)/m$ is decreasing thereafter. The same argument
with exponent one applies to the common-Schmidt law, uniformly in its
weights. These are sufficient bounds, with no claim of sharpness in
rank, accuracy, or confidence.


\end{document}
